\documentclass[11pt,a4paper]{article}

% ── Encoding & Fonts ─────────────────────────────────────────────
\usepackage{fontspec}
\setmonofont{DejaVu Sans Mono}[Scale=0.85]
% microtype removed (requires scalable fonts)

% ── Page Geometry ────────────────────────────────────────────────
\usepackage[margin=1in]{geometry}
\usepackage{parskip}

% ── Math ─────────────────────────────────────────────────────────
\usepackage{amsmath,amssymb,amsthm,mathtools}
\usepackage{stmaryrd}
\usepackage{bussproofs}

% ── Diagrams ─────────────────────────────────────────────────────
\usepackage{tikz}
\usetikzlibrary{trees,positioning,arrows.meta,decorations.pathreplacing}
\usepackage{forest}

% ── Code Highlighting ────────────────────────────────────────────
\usepackage[outputdir=build]{minted}
\setminted{
  fontsize=\small,
  breaklines,
  linenos=false,
  frame=single,
  framesep=2mm,
  baselinestretch=1.1,
  bgcolor=codebg,
}
\usemintedstyle{tango}

% ── Colors & Boxes ───────────────────────────────────────────────
\usepackage[dvipsnames]{xcolor}
\definecolor{codebg}{RGB}{248,248,248}
\definecolor{accentblue}{RGB}{30,90,160}
\definecolor{darkgreen}{RGB}{0,120,60}

\usepackage[most]{tcolorbox}
\tcbuselibrary{minted,skins,breakable}

\newtcolorbox{defnbox}[1][]{
  enhanced, breakable,
  colback=blue!3, colframe=accentblue,
  fonttitle=\bfseries, title=#1,
  boxrule=0.8pt, arc=2pt,
  left=6pt, right=6pt, top=4pt, bottom=4pt,
}

\newtcolorbox{intuition}[1][Intuition]{
  enhanced, breakable,
  colback=green!3, colframe=darkgreen,
  fonttitle=\bfseries, title=#1,
  boxrule=0.8pt, arc=2pt,
  left=6pt, right=6pt, top=4pt, bottom=4pt,
}

\newtcolorbox{workedexample}[1][Example]{
  enhanced, breakable,
  colback=orange!3, colframe=orange!70!black,
  fonttitle=\bfseries, title=#1,
  boxrule=0.8pt, arc=2pt,
  left=6pt, right=6pt, top=4pt, bottom=4pt,
}

\newtcolorbox{warning}[1][Warning]{
  enhanced, breakable,
  colback=red!3, colframe=red!70!black,
  fonttitle=\bfseries, title=#1,
  boxrule=0.8pt, arc=2pt,
  left=6pt, right=6pt, top=4pt, bottom=4pt,
}

\newtcolorbox{exercisebox}[1][Exercises]{
  enhanced, breakable,
  colback=violet!3, colframe=violet!60!black,
  fonttitle=\bfseries, title=#1,
  boxrule=0.8pt, arc=2pt,
  left=6pt, right=6pt, top=4pt, bottom=4pt,
}

% ── Exercise numbering ──────────────────────────────────────────
\newcounter{exercise}[section]
\renewcommand{\theexercise}{\thesection.\arabic{exercise}}

% ── Theorem environments ─────────────────────────────────────────
\theoremstyle{definition}
\newtheorem{definition}{Definition}[section]
\newtheorem{theorem}[definition]{Theorem}
\newtheorem{lemma}[definition]{Lemma}

% ── Links & References ───────────────────────────────────────────
\usepackage{hyperref}
\hypersetup{
  colorlinks=true,
  linkcolor=accentblue,
  urlcolor=accentblue,
  citecolor=accentblue,
}

% ── Headers ──────────────────────────────────────────────────────
\usepackage{fancyhdr}
\pagestyle{fancy}
\fancyhf{}
\fancyhead[L]{\small\textit{Lambda Calculus \& Lean\,4}}
\fancyhead[R]{\small\thepage}
\renewcommand{\headrulewidth}{0.4pt}

% ── Shortcuts ────────────────────────────────────────────────────
\newcommand{\lam}{\lambda}
\newcommand{\Lam}{\Lambda}
\newcommand{\reduce}{\longrightarrow_\beta}
\newcommand{\reduces}{\twoheadrightarrow_\beta}
\newcommand{\af}{\longrightarrow_\alpha}
\newcommand{\ef}{\longrightarrow_\eta}
\newcommand{\aeq}{\equiv_\alpha}
\newcommand{\FV}{\operatorname{FV}}
\newcommand{\BV}{\operatorname{BV}}
\newcommand{\proves}{\vdash}
\newcommand{\nproves}{\nvdash}
\newcommand{\Proves}{\Vdash}
\newcommand{\entails}{\models}
\newcommand{\Nat}{\mathbb{N}}
\newcommand{\church}[1]{\overline{#1}}
\newcommand{\TR}{\textsc{True}}
\newcommand{\FL}{\textsc{False}}
\newcommand{\step}{\longrightarrow}
\newcommand{\mstep}{\twoheadrightarrow}
\newcommand{\rmark}[1]{\quad\text{\small(#1)}}
\newcommand{\kw}[1]{\texttt{\textbf{#1}}}

\usepackage{enumitem}

% ── Title ────────────────────────────────────────────────────────
\title{%
  \vspace{-1cm}
  {\LARGE\bfseries A Programmer's Guide to Lambda Calculus}\\[8pt]
  {\large From First Principles to a Working Interpreter in Lean\,4}
}
\author{}
\date{}

\begin{document}
\maketitle
\thispagestyle{fancy}

\begin{abstract}
\noindent
This document is a self-contained introduction to the lambda calculus aimed at
working programmers who are comfortable reading code but have minimal background
in formal logic. We cover the untyped and simply typed lambda calculus, the
three core conversions ($\alpha$, $\beta$, $\eta$), and the formal
notation used in type theory---including \emph{judgments}, \emph{typing contexts}, and
the \emph{turnstile} ($\proves$). Every concept is paired with worked reduction
examples and executable Lean\,4 code.
\end{abstract}

\tableofcontents
\newpage

%======================================================================
\part{Foundations}
%======================================================================

No prior knowledge of formal logic is assumed. If you can write a function in any programming language, you have everything you need.

%======================================================================
\section{What Is Lambda Calculus?}\label{sec:what}
%======================================================================

\subsection{The Big Picture}

The story begins with a crisis. In 1928, the German mathematician David
Hilbert posed the \emph{Entscheidungsproblem} (``decision problem''): is there
a mechanical procedure that can determine, given any mathematical statement,
whether it is true or false? Hilbert believed the answer was yes---that
mathematics could, in principle, be fully automated.

He was wrong, and proving him wrong required answering a more basic question
first: \emph{what exactly is a ``mechanical procedure''?} Before you can show
that no procedure exists for a task, you need a rigorous definition of what a
procedure is. In the mid-1930s, several mathematicians independently proposed
answers.

\textbf{Alonzo Church} (1903--1995), a logician at Princeton, proposed
the lambda calculus in 1936. His idea was radical: computation is nothing more
than function application. Define functions, apply them to arguments, simplify
the results---that's it. No machine, no tape, no steps through a state
table. Just the algebra of functions.

\textbf{Alan Turing} (1912--1954), then a graduate student at Princeton
(Church was his PhD advisor!), independently proposed the Turing machine in
the same year---an imaginary device with a tape, a read/write head, and a table
of rules. His approach was more ``mechanical'': it modeled computation as a
physical process with states and transitions.

Church and Turing proved their models equivalent: anything computable by a
Turing machine can be expressed in $\lam$-calculus, and vice versa. This
remarkable fact is called the \textbf{Church--Turing thesis}. Both men also used
their formalisms to prove Hilbert wrong: the Entscheidungsproblem has no
solution. There are mathematical statements that no mechanical procedure can
decide.

Kurt G\"odel, who had already shaken the foundations of mathematics with his
incompleteness theorems (1931), showed his own system of recursive functions
was equivalent to both. So three independent approaches---functions, machines,
and recursion---all captured the same notion of computability. This convergence
is strong evidence that the definition is ``right.''

The word ``calculus'' here does not mean derivatives and integrals. In
mathematics, a \textbf{calculus} is just a formal system of rules for
manipulating symbols. The lambda calculus is a calculus of functions: its rules
tell you how to create functions, apply them, and simplify the results.

\subsection{Why Should a Programmer Care?}

Every functional language---Haskell, OCaml, Lean, even
the functional parts of JavaScript---is essentially $\lam$-calculus with
conveniences bolted on. Understanding it tells you what closures really are,
why currying works, and what happens under the hood when you pass a function
to \texttt{map}.

Here is a concrete example. In JavaScript:
\begin{minted}{javascript}
const add = x => y => x + y;
const increment = add(1);   // returns y => 1 + y
increment(5);               // returns 6
\end{minted}

This is directly analogous to the lambda calculus expression $\lam x.\,\lam y.\, x + y$.
The \texttt{=>} in JavaScript is literally the $\lam$. The pattern of a
function returning a function (\textbf{currying}) is how $\lam$-calculus
handles multiple arguments. And the mechanism by which \texttt{increment}
``remembers'' that $x = 1$---a \textbf{closure}---corresponds to the
$\beta$-reduction operation we will study in depth.

\begin{intuition}[Programmer's Mental Model]
Think of the $\lam$-calculus as an absurdly minimal programming language with
exactly \emph{three} constructs: variables, function definitions, and function
calls. No \texttt{if}, no \texttt{while}, no numbers, no strings. Despite this,
it is Turing-complete: it can compute anything that any programming language
can. We will see how to encode all of those ``missing'' features using nothing
but functions.
\end{intuition}

%======================================================================
\section{Reading the Notation}\label{sec:notation}
%======================================================================

Before diving in, here is a quick guide to the mathematical notation used
throughout this document. If you are comfortable with set notation and logic
symbols, feel free to skip ahead to Section~\ref{sec:what}. Otherwise, this
page is your cheat sheet---refer back to it whenever a symbol looks unfamiliar.

\subsection{Set Notation}

Sets are collections of distinct elements, written with curly braces. In this
document, we mainly use sets to talk about which variables appear in a term.

\begin{center}
\renewcommand{\arraystretch}{1.6}
\begin{tabular}{ccp{8.5cm}}
\textbf{Symbol} & \textbf{Name} & \textbf{Meaning and example} \\
\hline
$\{x, y\}$ & set literal &
  The set containing $x$ and $y$. Like a Python \texttt{set}: \texttt{\{x, y\}}. \\

$\in$ & element of &
  $x \in S$ means ``$x$ is a member of $S$.'' Like Python's \texttt{x in S}. \\

$\notin$ & not element of &
  $x \notin S$ means ``$x$ is \emph{not} in $S$.'' Like \texttt{x not in S}. \\

$\cup$ & union &
  $A \cup B$ is the set of elements in $A$ \emph{or} $B$ (or both).
  $\{x\} \cup \{y,z\} = \{x,y,z\}$. \\

$\setminus$ & set difference &
  $A \setminus B$ is $A$ with all elements of $B$ removed.
  $\{x,y,z\} \setminus \{y\} = \{x,z\}$.
  Think: Python's \texttt{A - B} for sets. \\

$\varnothing$ & empty set &
  The set with no elements. Like \texttt{set()} in Python.
  $\FV(e) = \varnothing$ means ``$e$ has no free variables.'' \\

$\subseteq$ & subset &
  $A \subseteq B$ means every element of $A$ is also in $B$. \\
\end{tabular}
\end{center}

\subsection{Arrows and Relations}

Arrows describe how terms transform during computation:

\begin{center}
\renewcommand{\arraystretch}{1.6}
\begin{tabular}{ccp{8.5cm}}
\textbf{Symbol} & \textbf{Name} & \textbf{Meaning} \\
\hline
$\longrightarrow_\beta$ & beta reduces to &
  One step of computation. $e \reduce e'$ means ``$e$ reduces to $e'$ in exactly one step.'' \\

$\twoheadrightarrow_\beta$ & multi-step reduces &
  Zero or more steps. $e \reduces e'$ means ``$e$ eventually becomes $e'$ after some number of reductions.'' \\

$\longrightarrow_\alpha$ & alpha converts to &
  A renaming of bound variables.
  $\lam x.\,x \;\af\; \lam y.\,y$. \\

$\equiv_\alpha$ & alpha-equivalent &
  Two terms are ``the same up to renaming.''
  $\lam x.\,x \;\aeq\; \lam y.\,y$. \\

$\longrightarrow_\eta$ & eta converts to &
  Wrapping/unwrapping a redundant lambda. $\lam x.\,f\;x \;\ef\; f$. \\

$\triangleq$ & defined as &
  Introduces a new name for an expression.
  $\TR \triangleq \lam t\,f.\,t$ means ``we \emph{define} \TR{} to be $\lam t\,f.\,t$.'' \\
\end{tabular}
\end{center}

\subsection{Functions and Substitution}

\begin{center}
\renewcommand{\arraystretch}{1.6}
\begin{tabular}{ccp{8.5cm}}
\textbf{Symbol} & \textbf{Name} & \textbf{Meaning} \\
\hline
$\FV(e)$ & free variables &
  The set of variables in $e$ that are not bound by any $\lam$.
  It is a \emph{function} from terms to sets of variables. \\

$e[x := a]$ & substitution &
  ``In term $e$, replace every free occurrence of variable $x$ with the term $a$.''
  This is the formal version of ``passing an argument.'' \\

$\lam x.\, e$ & lambda abstraction &
  A function with parameter $x$ and body $e$.
  Read as ``the function that takes $x$ and returns $e$.'' \\

$e_1\; e_2$ & application &
  Apply function $e_1$ to argument $e_2$.
  Juxtaposition (writing terms next to each other) means function call. \\
\end{tabular}
\end{center}

\subsection{Logic and Typing (used from Section~\ref{sec:stlc} onward)}

These symbols appear when we add types to the calculus:

\begin{center}
\renewcommand{\arraystretch}{1.6}
\begin{tabular}{ccp{8.5cm}}
\textbf{Symbol} & \textbf{Name} & \textbf{Meaning} \\
\hline
$\Gamma$ & context &
  A list of variable--type assumptions, like a symbol table.
  $\Gamma = x{:}\textsf{Nat},\; f{:}\textsf{Nat}{\to}\textsf{Bool}$. \\

$\proves$ & turnstile &
  Separates assumptions from conclusions.
  $\Gamma \proves e : \tau$ means ``given $\Gamma$, we can derive that $e$ has type $\tau$.''
  Detailed explanation in Section~\ref{sec:turnstile}. \\

$\nproves$ & negated turnstile &
  ``We \emph{cannot} derive\ldots'' Used to indicate a type error. \\

$\tau_1 \to \tau_2$ & function type &
  The type of functions taking $\tau_1$ and returning $\tau_2$.
  Right-associative: $A \to B \to C = A \to (B \to C)$. \\

$\forall$ & for all &
  Universal quantification. $\forall \alpha.\, \alpha \to \alpha$ means
  ``for every type $\alpha$, a function from $\alpha$ to $\alpha$.'' \\

$\exists$ & there exists &
  $\exists\, e'$ means ``there is some $e'$ such that\ldots'' \\
\end{tabular}
\end{center}

\subsection{Reading Inference Rules}

Typing rules are written as \textbf{inference rules}: premises above a
horizontal line, conclusion below, with an optional rule name on the right.

\[
\dfrac{\text{premise}_1 \qquad \text{premise}_2}
      {\text{conclusion}}
\;\textsc{RuleName}
\]

Read this as: ``\textit{if} premise$_1$ and premise$_2$ are both established,
\textit{then} the conclusion follows (by the rule named \textsc{RuleName}).''

Multiple rules can be stacked into a \textbf{proof tree} that builds a
complete derivation from axioms at the top down to the final conclusion at the
bottom. We will see many examples in Section~\ref{sec:typrules}.

\subsection{Grammar Notation}

Formal definitions use a grammar notation called BNF:

\begin{center}
\renewcommand{\arraystretch}{1.6}
\begin{tabular}{ccp{8.5cm}}
\textbf{Symbol} & \textbf{Name} & \textbf{Meaning} \\
\hline
$::=$ & ``is defined as'' &
  Introduces the possible forms of something. $e ::= \ldots$ means
  ``an expression $e$ can be any of the following shapes.'' \\

$\mid$ & ``or'' &
  Separates alternatives. $e ::= x \mid \lam x.\, e \mid e_1\;e_2$ means
  ``$e$ is either a variable, or a lambda, or an application.'' \\

$f^n\; x$ & iterated application &
  Apply $f$ a total of $n$ times: $f^0\;x = x$, $f^1\;x = f\;x$,
  $f^3\;x = f\;(f\;(f\;x))$. \\

$\church{n}$ & Church numeral &
  The overline notation $\church{n}$ denotes the Church encoding of
  the natural number $n$ (see Section~\ref{sec:church}). \\

$\uparrow^d_c(e)$ & shift &
  De Bruijn index adjustment: add $d$ to all free variable indices
  $\geq c$ in term $e$ (see Section~\ref{sec:debruijn}). \\
\end{tabular}
\end{center}

\subsection{Reading an Equation Step by Step}\label{sec:readingeqs}

Let's take one of the scariest-looking equations in this document and break it
down piece by piece:
\[
  \FV(\lam x.\, e) \;=\; \FV(e) \setminus \{x\}
\]

\begin{center}
\renewcommand{\arraystretch}{1.6}
\begin{tabular}{cp{10.5cm}}
\textbf{Piece} & \textbf{What it means} \\
\hline
$\FV(\ldots)$ &
  $\FV$ is a \emph{function} that takes a lambda term and returns the
  \emph{set of variables} that appear free (unbound) in it. \\

$\lam x.\, e$ &
  The input to $\FV$ is a lambda abstraction: a function with parameter $x$
  and body $e$. \\

$=$ &
  ``The result is\ldots'' \\

$\FV(e)$ &
  First, recursively compute the free variables of the body $e$. \\

$\setminus$ &
  Set difference (remove elements). In Python: \texttt{A - B}. \\

$\{x\}$ &
  The singleton set containing just the variable $x$. \\
\end{tabular}
\end{center}

\medskip\noindent
Putting it all together in English: ``The free variables of $\lam x.\, e$
are the free variables of the body $e$, \emph{minus} $x$ (because the
$\lam$ binds $x$, so it is no longer free).''

\medskip\noindent
Here is another example---the de Bruijn $\beta$-reduction rule from
Section~\ref{sec:debruijn}:
\[
  (\lam.\, e)\; s \;\;\reduce\;\; \uparrow^{-1}_0\!\bigl(e[0 := \uparrow^1_0(s)]\bigr)
\]

\begin{center}
\renewcommand{\arraystretch}{1.6}
\begin{tabular}{cp{10.5cm}}
\textbf{Piece} & \textbf{What it means} \\
\hline
$(\lam.\, e)\; s$ &
  A function (with body $e$) applied to argument $s$.
  This is the ``before'' side. \\

$\reduce$ &
  ``Beta-reduces to'' (one computation step). \\

$\uparrow^1_0(s)$ &
  Shift $s$ up by 1: increase all free variable indices in $s$ by 1,
  because $s$ is about to move under where a binder used to be. \\

$e[0 := \ldots]$ &
  In body $e$, replace every occurrence of variable 0 (the parameter)
  with the shifted argument. \\

$\uparrow^{-1}_0(\ldots)$ &
  Shift the whole result down by 1: we removed one $\lam$ binder,
  so all free variable indices decrease by 1 to compensate. \\
\end{tabular}
\end{center}

\medskip\noindent
In English: ``To apply a function, plug the argument into the body (adjusting
index numbering on the way in and out).''

%======================================================================
\section{Lean\,4 in 60 Seconds}\label{sec:lean}
%======================================================================

If you haven't used Lean, here is the minimum you need:

\begin{minted}{lean}
-- Inductive types (algebraic data types)
inductive MyNat where
  | zero : MyNat
  | succ : MyNat → MyNat

-- Pattern matching + recursion
def add : MyNat → MyNat → MyNat
  | .zero,   m => m
  | .succ n, m => .succ (add n m)

-- Option for partiality, do notation for chaining
def safeDivide (a b : Nat) : Option Nat := do
  if b == 0 then none else some (a / b)
\end{minted}

Key points: Lean requires termination by default. For $\beta$-reduction (which
may loop), we'll use a \textit{fuel} parameter.

%======================================================================
\section{The Syntax of Lambda Terms}\label{sec:syntax}
%======================================================================

\subsection{The Three Building Blocks}

\begin{defnbox}[Definition: Lambda Terms]
Given a countably infinite set of variables
$\mathcal{V} = \{x, y, z, \ldots\}$, the set $\Lambda$ of
\textbf{lambda terms} is defined by the grammar:
\[
  e \;::=\;
    x                              \quad\text{(variable)}
  \;\;\mid\;\; \lam x.\, e            \quad\text{(abstraction)}
  \;\;\mid\;\; e_1\; e_2              \quad\text{(application)}
\]
where $x \in \mathcal{V}$.
\end{defnbox}

\begin{intuition}[Reading the grammar]
The notation $e ::= \ldots$ defines what shapes an expression $e$ can take.
The vertical bar $\mid$ means ``or.'' So this says: ``an expression is
\emph{either} a variable $x$, \emph{or} an abstraction $\lam x.\, e$,
\emph{or} an application $e_1\; e_2$---and nothing else.'' This is recursive:
the body of a $\lam$ is itself an expression $e$, so terms nest to arbitrary
depth. If you have used algebraic data types (\texttt{enum} in Rust,
\texttt{data} in Haskell, \texttt{inductive} in Lean), this is the same idea.
\end{intuition}

Let's unpack each construct in detail.

\paragraph{Variables.} A variable is just a name---$x$, $y$, $z$, $f$,
\texttt{foo}, any identifier you like. It stands for a value that hasn't been
supplied yet, exactly like a function parameter in a programming language. By
itself, a variable is the simplest possible lambda term.

\paragraph{Abstraction (function definition).} This is how you create a
function. The notation $\lam x.\, e$ is read ``lambda $x$ dot $e$,'' and it
means ``the function that takes a parameter called $x$ and returns the
expression $e$.'' The $\lam$ (Greek letter lambda---hence the name of the
whole system) announces that a function is being defined. The variable $x$
after the $\lam$ is the \textbf{parameter name}. The dot separates the
parameter from the \textbf{body}---the expression $e$ that gets computed when
the function is called.

In different languages, this same concept is written differently:

\begin{center}
\renewcommand{\arraystretch}{1.4}
\begin{tabular}{ll}
\textbf{Language} & \textbf{Syntax for ``$\lam x.\, x + 1$''} \\
\hline
Lambda calculus & $\lam x.\, x + 1$ \\
JavaScript & \texttt{x => x + 1} \\
Python & \texttt{lambda x: x + 1} \\
Haskell & \texttt{\textbackslash x -> x + 1} \\
Lean\,4 & \texttt{fun x => x + 1} \\
\end{tabular}
\end{center}

They all mean the same thing. In each case, there is a marker (``$\lam$'',
``\texttt{=>}'', ``\texttt{lambda}'', etc.) that says ``a function is being
defined here.''

The word \textbf{abstraction} is used because the function ``abstracts over''
its parameter: $\lam x.\, x + 1$ is a general description of ``add one to
something,'' without specifying what that something is.

\paragraph{Application (function call).} Written $e_1\; e_2$, this means
``apply the function $e_1$ to the argument $e_2$.'' In most programming
languages you would write $f(a)$ with parentheses; in lambda calculus, you
simply write $f\; a$ (juxtaposition---putting things next to each other).
The function comes first, the argument comes second.

Note that $e_1$ and $e_2$ are both arbitrary lambda terms. The function can
be a variable (like $f\; x$), a lambda (like $(\lam x.\, x)\; 5$), or even
a complex expression. The argument can also be anything. This is what makes
lambda calculus so flexible: functions are values just like anything else, so
you can pass functions to functions, return functions from functions, and so on.

\medskip
\begin{center}
\renewcommand{\arraystretch}{1.5}
\begin{tabular}{lcl}
  \textbf{Lambda term} & \textbf{JavaScript analog} & \textbf{Meaning} \\
  \hline
  $x$ & \texttt{x} & a variable reference \\
  $\lam x.\, x$ & \texttt{x => x} & the identity function \\
  $(\lam x.\, x)\; y$ & \texttt{(x => x)(y)} & apply identity to $y$ \\
  $\lam f.\, \lam x.\, f\;(f\;x)$ & \texttt{f => x => f(f(x))} & apply $f$ twice \\
\end{tabular}
\end{center}

\subsection{Parenthesization Conventions}

Without conventions, expressions get buried in parentheses. For instance,
$f\; g\; h$ is ambiguous: do we mean ``apply $f$ to $g$, then apply the result
to $h$'' (i.e., $(f\;g)\;h$), or ``apply $f$ to the result of $g$ applied to
$h$'' (i.e., $f\;(g\;h)$)? Two standard conventions resolve all ambiguity:

\begin{enumerate}[leftmargin=2em]
\item \textbf{Application is left-associative.} This means $f\;g\;h$ is parsed
  as $(f\;g)\;h$. Read left to right: first apply $f$ to $g$, then apply
  the result to $h$. This is analogous to how \texttt{f(g)(h)} works in
  JavaScript (call $f$ with $g$, then call the result with $h$).

\item \textbf{The body of a lambda extends as far right as possible.} This
  means $\lam x.\, x\; y$ is parsed as $\lam x.\,(x\; y)$---``the function
  that takes $x$ and returns $x$ applied to $y$.'' It is \emph{not}
  $(\lam x.\, x)\; y$ (``apply the identity to $y$''). The lambda
  ``gobbles up'' everything to its right, until it hits a closing parenthesis
  or the end of the expression.
\end{enumerate}

\begin{workedexample}[Parsing practice]
How does $\lam f.\, \lam x.\, f\; x\; x$ parse?
\begin{enumerate}[leftmargin=2em]
\item The outer $\lam f$ grabs everything to its right: $\lam f.\,(\lam x.\, f\; x\; x)$.
\item The inner $\lam x$ grabs everything to \emph{its} right: $\lam f.\,(\lam x.\, (f\; x\; x))$.
\item Application is left-associative: $f\; x\; x = (f\; x)\; x$.
\item Fully parenthesized: $\lam f.\,(\lam x.\, ((f\; x)\; x))$.
\end{enumerate}
In English: ``the function that takes $f$ and returns a function that takes $x$
and returns $f$ applied to $x$, with the result applied to $x$ again.''
\end{workedexample}

\subsection{Multi-argument Functions and Currying}\label{sec:currying}

Lambda calculus functions take \textbf{exactly one argument}. There is no way to
write $\lam(x, y).\, e$---the grammar simply does not allow it. So how do we
handle functions that need more than one input?

The trick is called \textbf{currying} (named after logician Haskell Curry,
though the technique was first described by Moses Sch\"onfinkel in 1924 and
used by Gottlob Frege even earlier in the 1890s---Curry himself acknowledged
this, but the name stuck). Instead of a function of
two arguments, you write a function that takes the first argument and
\emph{returns another function} that takes the second argument.

Let's see this concretely. Suppose we want an ``add'' function that takes two
numbers. We cannot write $\lam(x, y).\, x + y$. Instead, we write:
\[
  \lam x.\, \lam y.\, x + y
\]
This is a function that takes $x$ and returns $\lam y.\, x + y$---a new
function that ``remembers'' $x$ and waits for $y$. Let's trace through using it
to add 3 and 4:

\begin{workedexample}[Currying in action]
\begin{align*}
  & \underbrace{(\lam x.\, \lam y.\, x + y)}_{\text{``add''}} \;\; 3 \;\; 4 \\[4pt]
  \reduce\;\; & \underbrace{(\lam y.\, 3 + y)}_{\text{``add 3 to\ldots''}} \;\; 4
    \rmark{substituted 3 for $x$; got a function waiting for $y$} \\[4pt]
  \reduce\;\; & 3 + 4 = 7
    \rmark{substituted 4 for $y$}
\end{align*}
The intermediate result $\lam y.\, 3 + y$ is a \textbf{partially applied}
function---it has received one argument but is still waiting for another. In
JavaScript, this pattern looks like \texttt{add(3)(4)} or equivalently:
\begin{minted}{javascript}
const add = x => y => x + y;
const add3 = add(3);    // y => 3 + y  (partial application!)
add3(4);                // 7
\end{minted}
\end{workedexample}

This extends to any number of arguments: a function of $n$ arguments is a chain
of $n$ nested lambdas. As shorthand, we write $\lam x\, y\, z.\, e$ for
$\lam x.\, \lam y.\, \lam z.\, e$.

\subsection{Let Bindings Are Just Application}\label{sec:let}

Most programming languages have \texttt{let} for naming intermediate values:
\begin{minted}{javascript}
let x = 2 + 3 in x * x    // evaluates to 25
\end{minted}

Lambda calculus has no built-in \texttt{let}, but it doesn't need one.
A \texttt{let} expression is just syntactic sugar for a function application:
\[
  \underbrace{\texttt{let } x = e_1 \texttt{ in } e_2}
  _{\text{let binding}}
  \quad\equiv\quad
  \underbrace{(\lam x.\, e_2)\; e_1}
  _{\text{application}}
\]

This works because $(\lam x.\, e_2)\; e_1$ does exactly what \texttt{let} does:
it evaluates $e_1$, binds the result to $x$, then evaluates $e_2$ with $x$
available. Let's verify:
\[
  (\lam x.\, x \cdot x)\;(2 + 3) \;\reduce\; (2 + 3) \cdot (2 + 3) = 5 \cdot 5 = 25
\]

\begin{intuition}[Why this matters]
Whenever you see \texttt{let} in a functional language, remember: it's not a
new primitive. It's just a $\lam$ being applied immediately. This is one of the
beauties of lambda calculus---seemingly essential features dissolve into the
three basic constructs. We add \texttt{let} support to our parser in
Section~\ref{sec:parser} as a desugaring step.
\end{intuition}

%======================================================================

indices, type checking, and ultimately a dependent type checker that verifies
proofs.

\subsection{Implementing It: Representing Terms}\label{sec:namedimpl}
%======================================================================

Your first design decision: how do you represent a $\lam$-term as a data
structure?

\subsubsection{Think About It}

Go back to Section~\ref{sec:syntax}. A $\lam$-term has exactly three possible
shapes: it's a variable, a lambda abstraction, or an application. That's it---
there are no other cases.

\begin{exercisebox}[Design challenge]
\stepcounter{exercise}\textbf{\theexercise.}\; Before reading further, write down (on paper or in your head) a data type that can represent any $\lam$-term. What fields does each case need? How would you represent $\lam x.\, x\;y$?
\end{exercisebox}

\subsubsection{The Answer: An Algebraic Data Type}

Each shape becomes a constructor:

\begin{minted}{lean}
inductive Term where
  | var : String → Term            -- a variable carries its name
  | lam : String → Term → Term     -- a lambda carries parameter name + body
  | app : Term → Term → Term       -- an application carries two sub-terms
  deriving Repr, DecidableEq, Inhabited
\end{minted}

The \texttt{deriving} clause asks Lean to auto-generate three typeclass
instances for \texttt{Term}:

\begin{itemize}[leftmargin=2em]
\item \texttt{Repr}: Produces string representations so \texttt{\#eval} can
  display terms in the console (e.g., \texttt{Term.lam "x" (Term.var "x")}).
\item \texttt{DecidableEq}: Lets us use \texttt{==} to compare terms for
  structural equality, and also unlocks \texttt{if x == y} and list membership
  tests like $x \in \mathit{xs}$---both of which we'll need in substitution.
\item \texttt{Inhabited}: Provides a default value for the type (here,
  \texttt{Term.var ""}). Lean sometimes requires this when using certain
  combinators or partial functions.
\end{itemize}

The term $\lam x.\, x\;y$ becomes \texttt{lam "x" (app (var "x") (var "y"))}.
This is a tree---the same tree structure we visualized in
Section~\ref{sec:beta}. Variables are leaves, applications are internal nodes
with two children, and lambdas are internal nodes with one child plus a name
annotation.

Notice how the data type mirrors the grammar almost exactly:
\begin{center}
\renewcommand{\arraystretch}{1.4}
\begin{tabular}{lcl}
\textbf{Grammar} & & \textbf{Lean} \\
\hline
$e ::= x$          & $\longleftrightarrow$ & \texttt{var : String → Term} \\
$e ::= \lam x.\, e$ & $\longleftrightarrow$ & \texttt{lam : String → Term → Term} \\
$e ::= e_1\; e_2$    & $\longleftrightarrow$ & \texttt{app : Term → Term → Term} \\
\end{tabular}
\end{center}

This is a recurring pattern: \emph{grammars translate directly into algebraic
data types}. If you can write the grammar, you can write the type.

%======================================================================
\section{Free and Bound Variables}\label{sec:fv}
%======================================================================

\subsection{Scope in Lambda Calculus}

Every programmer understands \textbf{scope}: a variable declared inside a
function is only visible within that function. Lambda calculus has the same
concept.

When you write $\lam x.\, e$, the $\lam$ creates a \textbf{binding} for the
variable $x$, and the \textbf{scope} of that binding is the body $e$. Within
$e$, the variable $x$ is said to be \textbf{bound}---it has a definition
(it's the function's parameter). A variable that is \emph{not} bound by any
enclosing $\lam$ is \textbf{free}---it refers to something defined elsewhere,
like a global variable or a variable from an outer scope.

\begin{workedexample}[Identifying free and bound variables]
Consider $\lam x.\, x\; y$. This is a function that takes $x$ and applies it
to $y$. Here:
\begin{enumerate}[leftmargin=2em]
\item $x$ is \textbf{bound}---it's the parameter of the $\lam$.
\item $y$ is \textbf{free}---no $\lam$ binds it, so its meaning depends on
  the surrounding context.
\end{enumerate}
As a Python analogy:
\begin{minted}{python}
y = 10            # y is defined outside
f = lambda x: x + y  # x is bound (parameter), y is free (from outer scope)
\end{minted}
If $y$ is not defined anywhere, the expression is incomplete---it has an
undefined variable. A term with no free variables is self-contained.
\end{workedexample}

\subsection{Formal Definition}

\begin{defnbox}[Definition: Free Variables]
\vspace{-0.8em}
\begin{align*}
\FV(x)           &= \{x\} \\
\FV(\lam x.\, e) &= \FV(e) \setminus \{x\} \\
\FV(e_1\; e_2)   &= \FV(e_1) \cup \FV(e_2)
\end{align*}
\end{defnbox}

\begin{intuition}[Reading the definition line by line]
This is a recursive function from terms to sets of variable names. Reading each
line (see Section~\ref{sec:notation} for the symbols):
\begin{enumerate}[leftmargin=2em]
\item $\FV(x) = \{x\}$: A bare variable $x$ is free---it's not under any
  $\lam$, so the set of free variables is just $\{x\}$.
\item $\FV(\lam x.\, e) = \FV(e) \setminus \{x\}$: A lambda \emph{binds} $x$,
  so take the free variables of the body $e$ and \emph{remove} $x$ from them
  (the ``$\setminus$'' is set difference, like subtracting $\{x\}$ from the set).
\item $\FV(e_1\; e_2) = \FV(e_1) \cup \FV(e_2)$: An application's free
  variables are the \emph{union} ($\cup$) of both subterms' free variables---anything free on either side is free in the whole expression.
\end{enumerate}
\end{intuition}

\begin{workedexample}[Computing $\FV$ step by step]
Let's compute $\FV(\lam x.\, x\; y)$ from scratch:
\begin{align*}
  \FV(\lam x.\, x\; y)
  &= \FV(x\; y) \setminus \{x\}
    \rmark{rule 2: peel off the $\lam x$ and remove $x$} \\
  &= \bigl(\FV(x) \cup \FV(y)\bigr) \setminus \{x\}
    \rmark{rule 3: union both sides of the application} \\
  &= \bigl(\{x\} \cup \{y\}\bigr) \setminus \{x\}
    \rmark{rule 1: each variable is its own singleton} \\
  &= \{x, y\} \setminus \{x\}
    \rmark{perform the union} \\
  &= \{y\}
    \rmark{remove $x$ from the set}
\end{align*}
Only $y$ is free. This matches our earlier analysis: $x$ is the parameter
(bound), $y$ comes from outside (free).
\end{workedexample}

\subsection{Closed Terms (Combinators)}

A term with $\FV(e) = \varnothing$ (no free variables at all) is
\textbf{closed}---it's completely self-contained, like a standalone program
with no undefined variables. Closed terms are also called \textbf{combinators}.
The three most famous are:
\[
  \mathbf{I} = \lam x.\, x
  \qquad
  \mathbf{K} = \lam x\, y.\, x
  \qquad
  \mathbf{S} = \lam f\, g\, x.\, f\;x\;(g\;x)
\]
$\mathbf{I}$ is the identity (returns its argument unchanged). $\mathbf{K}$
takes two arguments and returns the first, ignoring the second. $\mathbf{S}$
takes three arguments and ``distributes'' the third to the first two.

These three combinators form a basis: every $\lam$-term can be expressed using
just $\mathbf{S}$ and $\mathbf{K}$ (since $\mathbf{I} = \mathbf{S}\;\mathbf{K}\;\mathbf{K}$, as we will verify later).

\begin{intuition}[Historical note: Combinatory logic]
The $\mathbf{S}$, $\mathbf{K}$, $\mathbf{I}$ combinators actually predate
lambda calculus. They were introduced by Moses Sch\"onfinkel in a 1924 paper,
``On the building blocks of mathematical logic,'' in which he showed that
all of logic could be expressed using just two operations ($\mathbf{S}$ and
$\mathbf{K}$). Haskell Curry independently rediscovered and greatly expanded
these ideas starting in 1927, developing them into a full theory called
\textbf{combinatory logic}. Church's lambda calculus (1936) took a different
but equivalent path---using variable binding rather than combinators as the
primitive. The two systems are interconvertible and equally expressive.
\end{intuition}

\subsection{The Danger: Variable Capture}

\begin{warning}[Variable Capture]
When we perform substitution (which we will define formally in
Section~\ref{sec:beta}), we must be careful not to accidentally change the
meaning of a term by allowing a free variable to become bound. This problem is
called \textbf{variable capture}.

Consider: what happens if we substitute $y$ for $x$ in $\lam y.\, x$? The
original term is a function that ignores its argument and returns whatever $x$
is (a free variable). Naively replacing $x$ with $y$ gives $\lam y.\, y$---but
that's the \emph{identity} function, a completely different thing! The free $y$
got ``captured'' by the $\lam y$ binder.

The fix: \textbf{rename the binder first}. Change $\lam y.\, x$ to
$\lam z.\, x$ (which is the same function---the parameter name doesn't matter),
and \emph{then} substitute to get $\lam z.\, y$. This correctly preserves the
meaning: a function that ignores its argument and returns $y$.

This renaming operation is called \textbf{alpha conversion}, and it's the
subject of the next section.
\end{warning}

%======================================================================

\subsection{Implementing It: Free Variables in Lean}\label{sec:impl-fv}

Let's turn the formal definition of free variables into code.

\subsubsection{Free Variables}

We need $\FV(e)$ in code. Go back to the formal definition from
Section~\ref{sec:fv}:
$\FV(x) = \{x\}$, \;$\FV(\lam x.\, e) = \FV(e) \setminus \{x\}$, \;
$\FV(e_1\; e_2) = \FV(e_1) \cup \FV(e_2)$.

\begin{exercisebox}[Implementation challenge]
\stepcounter{exercise}\textbf{\theexercise.}\; Translate the free variables definition into a function \texttt{freeVars : Term → List String}. Each case of the formal definition becomes one pattern-match branch. Try it yourself before reading on.
\end{exercisebox}

Look at the formal definition again. It has three cases, one for each shape of
term. Your data type has three constructors. So the function will have three
branches. Try writing them as comments first:

\begin{minted}{text}
def freeVars : Term → List String
  | .var x     => ???   -- FV(x) = {x}
  | .lam x e   => ???   -- FV(λx.e) = FV(e) \ {x}
  | .app e₁ e₂ => ???   -- FV(e₁ e₂) = FV(e₁) ∪ FV(e₂)
\end{minted}

Now fill in each \texttt{???}. For the variable case, the free variables are
just the variable itself: a one-element list. For the lambda case, you need
``all free variables of the body, except $x$''---that's a filter. For the
application case, you need the union of both sides---concatenation with
duplicates removed.

The translation is almost mechanical:

\begin{minted}{lean}
def freeVars : Term → List String
  | .var x     => [x]                                  -- FV(x) = {x}
  | .lam x e   => (freeVars e).filter (· != x)         -- FV(λx.e) = FV(e)\{x}
  | .app e₁ e₂ => (freeVars e₁ ++ freeVars e₂).eraseDups  -- union
\end{minted}

This is a pattern you'll see repeatedly: \textbf{formal definitions translate
almost line-for-line into code}. If you understand the math, you can write the
program.
\section{Alpha Conversion ($\alpha$)}\label{sec:alpha}
%======================================================================

\subsection{The Core Idea: Names Don't Matter}

Here is a question: are these two JavaScript functions the same?

\begin{minted}{javascript}
const f = x => x + 1;
const g = y => y + 1;
\end{minted}

Obviously yes. They both add 1 to their argument. The fact that one calls its
parameter \texttt{x} and the other calls it \texttt{y} is irrelevant---parameter
names are arbitrary labels.

\textbf{Alpha conversion} ($\alpha$-conversion) is the formal version of this
observation. It says that you can rename a bound variable at any time, as long
as you do it consistently and don't create a name clash.

\begin{defnbox}[The $\alpha$-conversion Rule]
\[
  \lam x.\, e \;\;\af\;\; \lam y.\, e[x := y]
  \qquad\text{provided } y \notin \FV(e) \text{ and } y \notin \BV(e)
\]
Two terms differing only in bound variable names are $\alpha$-equivalent:
$\lam x.\,x \;\aeq\; \lam y.\,y \;\aeq\; \lam z.\,z$.
\end{defnbox}

The notation $e[x := y]$ here means ``replace every free $x$ in $e$ with $y$''
(the same substitution operation from the next section). The two conditions
ensure safety: $y$ must not already appear in $e$ either free or bound, so the
renaming doesn't clash with any existing names.

\subsection{Why Alpha Conversion Matters}

Alpha conversion is not just a theoretical nicety. As we saw in
Section~\ref{sec:fv}, without it, substitution can go wrong due to variable
capture. Every lambda calculus implementation must handle $\alpha$-equivalence.
In practice, there are three common approaches:
\begin{enumerate}[leftmargin=2em]
\item Rename variables on the fly during substitution (the classical approach,
  which we implement in our named Lean\,4 interpreter).
\item Use \textbf{de Bruijn indices}, where variables are numbers rather than
  names, so $\alpha$-equivalence becomes structural equality (see
  Section~\ref{sec:debruijn}).
\item Use a locally nameless representation, combining the best of both.
\end{enumerate}

\begin{workedexample}[Alpha-converting $\lam x.\, \lam y.\, x\; y$]
Rename the outer $x$ to $z$:
\begin{enumerate}[leftmargin=2em]
\item Check that $z$ is not free in the body $(\lam y.\, x\;y)$. We have $\FV(\lam y.\, x\;y) = \{x\}$, and $z \notin \{x\}$. \checkmark
\item Replace every free $x$ in the body with $z$: $\lam y.\, z\; y$.
\item Result: $\lam z.\, \lam y.\, z\; y$.
\end{enumerate}
Now rename the inner $y$ to $w$:
\begin{enumerate}[leftmargin=2em]
\item Check $w \notin \FV(z\;y) = \{z, y\}$. Since $w$ is neither $z$ nor $y$: \checkmark
\item Replace: $z\; w$.
\item Result: $\lam z.\, \lam w.\, z\; w$.
\end{enumerate}
The original $\lam x.\, \lam y.\, x\; y$ and the final
$\lam z.\, \lam w.\, z\; w$ are $\alpha$-equivalent ($\aeq$). They are the
same function---just with different parameter names.
\end{workedexample}

\begin{exercisebox}[Exercises: Variables and Renaming]
\stepcounter{exercise}\textbf{\theexercise.}\; Compute $\FV(\lam x.\, \lam y.\, x\; z\; y)$. Which variables are free? Which are bound?

\stepcounter{exercise}\textbf{\theexercise.}\; Is the term $\lam f.\, \lam x.\, f\;(f\;x)$ closed? Why or why not?

\stepcounter{exercise}\textbf{\theexercise.}\; Alpha-convert $\lam x.\, \lam x.\, x$ so that no variable name is used twice. What does this function do?

\stepcounter{exercise}\textbf{\theexercise.}\; Can you alpha-convert $\lam x.\, x\;y$ to $\lam y.\, y\;y$? Why or why not?
\end{exercisebox}

%======================================================================
\section{Beta Reduction ($\beta$)}\label{sec:beta}
%======================================================================

\subsection{The Heart of Computation}

Alpha conversion renames things. \textbf{Beta reduction} is where actual
\emph{computation} happens. It is the single rule that makes lambda calculus a
model of computation, and it corresponds to something every programmer does
constantly: \textbf{calling a function}.

When you write $(\lam x.\, x + 1)\; 5$ and evaluate it, you get $5 + 1 = 6$.
What happened? You took the function $\lam x.\, x + 1$ (``add 1 to the
argument''), and \emph{plugged in} 5 for $x$ in the body. That ``plugging in''
is beta reduction.

\begin{defnbox}[The $\beta$-reduction Rule]
\[
  \boxed{\;(\lam x.\, e)\; a \;\;\reduce\;\; e[x := a]\;}
\]
Read this as: ``a function $\lam x.\, e$ applied to an argument $a$
reduces to the body $e$ with every free $x$ replaced by $a$.''

\medskip
The expression $(\lam x.\, e)\; a$ is called a \textbf{redex} (short for
``reducible expression''---a spot where a reduction can happen). The result
$e[x := a]$ is called the \textbf{reduct}.
\end{defnbox}

\begin{intuition}[What is happening mechanically?]
Beta reduction is function calling, broken into two parts:
\begin{enumerate}[leftmargin=2em]
\item \textbf{Pattern match.} Look at the term. If it has the shape
  ``something applied to something,'' \emph{and} the left side is a
  $\lam$-abstraction, then you have a redex.
\item \textbf{Substitute.} Take the body of the $\lam$, find every place the
  parameter appears, and replace it with the argument. Erase the $\lam$
  wrapper.
\end{enumerate}
That's it. All of computation in the $\lam$-calculus is just pattern-matching
on this shape and performing substitution.
\end{intuition}

\subsection{Why Is This ``Computation''?}

At first, beta reduction might seem too simple to deserve the name
``computation.'' All we're doing is substituting one thing for another---how
can that express loops, conditionals, data structures, and everything else a
real program needs?

This question is worth sitting with. In 1933, when Church first defined the
$\lam$-calculus, it was not at all obvious that it had anything to do with
computation. Church was working on the foundations of logic, trying to build
a system where mathematical functions could be treated as first-class objects.
The substitution rule was just the natural way functions should behave:
if you define $f(x) = x + 1$ and then ask for $f(5)$, you substitute $5$
for $x$ and get $6$. Nothing controversial there.

The surprise came when Church realized that this simple mechanism was
\emph{universal}. The key insight unfolds in three steps.

\paragraph{Step 1: You can encode data as behavior.}
Lambda calculus has no built-in data types. But you can represent any piece of
data as a function that describes \emph{how it behaves}. Define
$\TR \triangleq \lam t.\, \lam f.\, t$ and $\FL \triangleq \lam t.\,
\lam f.\, f$. These are just functions, but they behave like booleans:
$\TR$ selects the first of two arguments, $\FL$ selects the second.
Now ``if $b$ then $x$ else $y$'' is just $b\;x\;y$---a beta reduction:
\[
  \TR\; \texttt{"yes"}\; \texttt{"no"}
  \;\reduce\; (\lam f.\, \texttt{"yes"})\; \texttt{"no"}
  \;\reduce\; \texttt{"yes"}
\]
We just evaluated a conditional using nothing but substitution. Similarly,
define the number $n$ as a function that applies $f$ exactly $n$ times:
$\church{3} = \lam f.\, \lam x.\, f\;(f\;(f\;x))$. Now ``add 1'' is just
``apply $f$ one more time''---again, a beta reduction. Arithmetic from
substitution.

\paragraph{Step 2: You can express recursion.}
The term $(\lam x.\, x\;x)\;(\lam x.\, x\;x)$ reduces to itself. This
\emph{is} an infinite loop, expressed as substitution. And by modifying this
pattern slightly (the $\mathbf{Y}$ combinator, Section~\ref{sec:recursion}),
we get controlled recursion: a function that calls itself as many times as
needed and then stops.

This is where the magic happens. With data encoding and recursion, you can
express every computable function as a chain of beta reductions.

\paragraph{Step 3: This is \emph{all} of computation.}
Church's 1936 paper made this precise. He defined a class of functions called
``$\lam$-definable functions''---functions on natural numbers that could be
expressed as $\lam$-terms using Church numerals. He then showed that this class
included all functions that mathematicians intuitively considered ``effectively
calculable'': addition, multiplication, exponentiation, primality testing,
and crucially, the recursive functions defined earlier by G\"odel and Herbrand.

This was a bold claim. Church was saying: \emph{if a function can be computed
at all, by any method whatsoever, then it can be computed by beta reduction.}

Turing, independently working on the same problem, proved the same thing for
his Turing machines using a different approach. When Church (1936) and Turing
(1937) proved that $\lam$-definable functions and Turing-computable functions
were exactly the same class, the result was electrifying. Two completely
different formalizations of ``computation''---one based on abstract
symbol manipulation, the other on an imaginary physical machine---captured
the same set of functions. This convergence was strong evidence that both
had found something fundamental.

This is the \textbf{Church--Turing thesis}: any effectively calculable function
is computable by a Turing machine / expressible in $\lam$-calculus. It is not a
theorem (you cannot prove it, because ``effectively calculable'' is an informal
notion), but no counterexample has ever been found. Every alternative model of
computation---register machines, cellular automata, quantum computers,
programming languages---has turned out to be equivalent in power.

\paragraph{The philosophical punchline.}
A system with nothing but function definitions and function calls, where the
only operation is substituting arguments into function bodies, is powerful
enough to express \emph{all of computation}. Beta reduction is not just
``a rule''---it is \emph{the} rule. Every program you have ever written,
in any language, could in principle be expressed as a lambda term and
evaluated by iterated substitution. The elaborate machinery of modern
computers---registers, memory, instruction sets---is, in a precise
mathematical sense, no more powerful than pen-and-paper substitution.

\subsection{Visualizing Reduction as Tree Surgery}

Lambda terms have a natural tree structure. Seeing them as trees makes
$\beta$-reduction much more concrete: it is literally a tree transformation
where you splice the argument subtree into the body.

Consider $(\lam x.\, f\; x)\; a$. As a tree, the top node is an application
(``@''), whose left child is $\lam x.\, f\;x$ and whose right child is $a$:

\begin{center}
\begin{forest}
  for tree={
    circle, draw, minimum size=1.8em, inner sep=1pt,
    s sep=18pt, l sep=14pt,
    edge={->,>=stealth},
    font=\small
  }
  [@
    [$\lam x$
      [@
        [$f$]
        [$x$, fill=orange!25]
      ]
    ]
    [$a$, fill=blue!15]
  ]
\end{forest}
\end{center}

$\beta$-reduction works in three steps: (1) identify the redex (the ``@'' whose
left child is a $\lam$), (2) take the right child ($a$, shaded blue) and
substitute it for every leaf labeled $x$ (shaded orange), (3) remove the
$\lam$-node and the ``@'' above it. The result:

\begin{center}
\begin{forest}
  for tree={
    circle, draw, minimum size=1.8em, inner sep=1pt,
    s sep=18pt, l sep=14pt,
    edge={->,>=stealth},
    font=\small
  }
  [@
    [$f$]
    [$a$, fill=blue!15]
  ]
\end{forest}
\end{center}

The variable leaf $x$ has been replaced by the argument $a$. This is
$\beta$-reduction visualized as tree surgery.

\medskip
Here is a more complex example. The $\mathbf{K}$ combinator applied to two
arguments: $(\lam x.\, \lam y.\, x)\; a\; b$. Because application is
left-associative, this is $((\lam x.\, \lam y.\, x)\; a)\; b$:

\begin{center}
\begin{forest}
  for tree={
    circle, draw, minimum size=1.8em, inner sep=1pt,
    s sep=15pt, l sep=13pt,
    edge={->,>=stealth},
    font=\small
  }
  [@
    [@
      [$\lam x$
        [$\lam y$
          [$x$, fill=orange!25]
        ]
      ]
      [$a$, fill=blue!15]
    ]
    [$b$, fill=green!15]
  ]
\end{forest}
$\;\;\xrightarrow{\;\beta\;}\;\;$
\begin{forest}
  for tree={
    circle, draw, minimum size=1.8em, inner sep=1pt,
    s sep=15pt, l sep=13pt,
    edge={->,>=stealth},
    font=\small
  }
  [@
    [$\lam y$
      [$a$, fill=blue!15]
    ]
    [$b$, fill=green!15]
  ]
\end{forest}
$\;\;\xrightarrow{\;\beta\;}\;\;$
\begin{forest}
  for tree={
    circle, draw, minimum size=1.8em, inner sep=1pt,
    s sep=15pt, l sep=13pt,
    edge={->,>=stealth},
    font=\small
  }
  [$a$, fill=blue!15]
\end{forest}
\end{center}

First reduction: the inner ``@'' fires, substituting $a$ for $x$ inside
$\lam y.\, x$, giving $\lam y.\, a$. Second reduction: the outer ``@'' fires,
substituting $b$ for $y$ in $a$---but $a$ has no $y$, so $b$ is discarded.
Final result: $a$.

\subsection{Substitution: The Engine Behind Beta Reduction}

Beta reduction says $(\lam x.\, e)\;a \reduce e[x := a]$. The hard part is
not the pattern-matching (finding a $\lam$ applied to something). The hard part
is the operation $e[x := a]$: ``replace every free $x$ in $e$ with $a$.'' This
is called \textbf{substitution}, and getting it right is the most subtle aspect
of lambda calculus.

\subsubsection{Why Not Just Find-and-Replace?}

Your first instinct might be: ``substitution is just find-and-replace---search
for every $x$ in $e$ and swap in $a$.'' And in many cases, that works perfectly:

\[
  (x + 1)[x := 5] \;=\; 5 + 1 \quad\checkmark
\]

But lambda calculus has \textbf{variable binding}---the $\lam$ operator creates
a new scope. This makes naive find-and-replace dangerous. Consider three
scenarios that show why:

\paragraph{Problem 1: Don't replace bound variables.}
In $(\lam x.\, x)[x := 5]$, should we replace $x$ with $5$? No! The $x$ inside
$\lam x.\, x$ is \emph{bound} by that $\lam$---it's a fresh parameter, a local
variable with its own scope. The outer substitution should not touch it. Just as
in this Python code:
\begin{minted}{python}
x = 5
f = lambda x: x    # this x is a NEW x, shadowing the outer one
f(10)               # returns 10, not 5
\end{minted}

The correct result is: $(\lam x.\, x)[x := 5] = \lam x.\, x$. The substitution
stops at the $\lam x$ because $x$ is \textbf{shadowed}.

\paragraph{Problem 2: Do substitute into other binders.}
What about $(\lam y.\, x)[x := 5]$? Here the $\lam$ binds $y$, not $x$. The
$x$ in the body is still free---it refers to the outer $x$, which we're
substituting. So we should go inside:

\[
  (\lam y.\, x)[x := 5] \;=\; \lam y.\, 5 \quad\checkmark
\]

\paragraph{Problem 3: Variable capture---the real danger.}
Now the tricky case. What about $(\lam y.\, x)[x := y]$? We want to replace the
free $x$ with $y$. Naively:
\[
  (\lam y.\, x)[x := y] \;\overset{\text{WRONG}}{=}\; \lam y.\, y
\]
We got $\lam y.\, y$---the identity function! But the original term
$\lam y.\, x$ is \emph{not} the identity. It's a function that ignores its
argument and returns the (free) variable $x$. By blindly substituting $y$
into a scope where $y$ is already bound, we accidentally \textbf{captured} the
free $y$, turning it into a bound variable. The $y$ we substituted in now refers
to the lambda's parameter instead of the outer variable.

This is called \textbf{variable capture}, and it produces wrong results. The
fix is to rename the conflicting binder before substituting:
\[
  (\lam y.\, x)[x := y]
    \;=\; (\lam z.\, x)[x := y]     \quad\text{(rename $y \to z$ first)}
    \;=\; \lam z.\, y               \quad\checkmark
\]
Now the free $y$ stays free, as intended.

\subsubsection{The Complete Rules}

With the motivation clear, here is the formal definition. Each rule handles one
shape of expression:

\begin{defnbox}[Capture-Avoiding Substitution]
\vspace{-0.8em}
\begin{align*}
x[x := a]              &= a
  & \text{\small(1) target variable: replace it} \\
y[x := a]              &= y
  & \text{\small(2) different variable: leave it} \\
(\lam x.\, e)[x := a]  &= \lam x.\, e
  & \text{\small(3) shadowed: stop} \\
(\lam y.\, e)[x := a]  &= \lam y.\, (e[x := a])
  & \text{\small(4) safe: go inside} \\
  &\qquad\qquad \text{if } y \neq x,\; y \notin \FV(a) & \\
(\lam y.\, e)[x := a]  &= \lam z.\, (e[y := z][x := a])
  & \text{\small(5) capture: rename first} \\
  &\qquad\qquad \text{if } y \neq x,\; y \in \FV(a),\; z \text{ fresh} & \\
(e_1\; e_2)[x := a]    &= (e_1[x := a])\;(e_2[x := a])
  & \text{\small(6) application: both sides}
\end{align*}
\end{defnbox}

\subsubsection{Rule by Rule, in Depth}

Let's go through each rule carefully, with the reasoning behind it and a
concrete example.

\paragraph{Rule 1: $x[x := a] = a$.}
The simplest case. We found the variable we're substituting for, so we replace
it. Example: $x[x := 5] = 5$.

\paragraph{Rule 2: $y[x := a] = y$ when $y \neq x$.}
We found a variable, but it's not the one we're looking for. Leave it alone.
Example: $y[x := 5] = y$.

\paragraph{Rule 3: $(\lam x.\, e)[x := a] = \lam x.\, e$ (shadowing).}
The $\lam$ creates a \emph{new} $x$ that shadows the one we're substituting
for. Inside this $\lam$, every $x$ refers to the parameter, not the outer
variable. So we stop---the substitution has no effect.

Think of it like nested function definitions:
\begin{minted}{python}
x = 10
def f(x):       # this x shadows the outer x
    return x    # refers to the parameter, not 10
\end{minted}

Example: $(\lam x.\, x\;x)[x := \mathbf{I}] = \lam x.\, x\;x$. The
$x$'s inside are bound, not free.

\paragraph{Rule 4: $(\lam y.\, e)[x := a] = \lam y.\, (e[x := a])$ (safe descent).}
The $\lam$ binds a \emph{different} variable $y$, so it doesn't shadow $x$. We
can safely substitute inside the body. But there is a condition: $y \notin
\FV(a)$. This means the term $a$ (what we're plugging in) does not mention $y$
freely. Since $a$ doesn't mention $y$, going under $\lam y$ can't accidentally
capture anything.

Example: $(\lam y.\, x\;y)[x := \lam z.\, z] = \lam y.\, (\lam z.\, z)\;y$.
Here $a = \lam z.\, z$, and $y \notin \FV(\lam z.\, z) = \varnothing$, so it's
safe.

\paragraph{Rule 5: $(\lam y.\, e)[x := a] = \lam z.\, (e[y := z][x := a])$
(capture avoidance).}
This is the critical rule---the one that makes substitution ``capture-avoiding.''
We need to substitute inside $\lam y$, but $a$ contains $y$ as a free variable.
If we went ahead and substituted, the $\lam y$ would capture the $y$ in $a$,
changing its meaning.

The fix is a two-step process:
\begin{enumerate}[leftmargin=2em]
\item $\alpha$-convert: rename the bound $y$ to a fresh name $z$ that appears
  nowhere in the expression (specifically, $z \notin \FV(e) \cup \FV(a) \cup \{x\}$).
  This gives $\lam z.\, e[y := z]$.
\item Now it's safe to substitute $a$ into the body, because the binder is $z$,
  not $y$, so no capture can occur.
\end{enumerate}

\paragraph{Rule 6: $(e_1\;e_2)[x := a] = (e_1[x := a])\;(e_2[x := a])$ (application).}
An application has no binders---it's just two expressions side by side. So we
substitute into both independently. Example:
$(f\;x)[x := 5] = f\;5$.

\subsubsection{Worked Example: Capture Avoidance in Action}

Let's trace through the dangerous case step by step. We want to compute:
\[
  (\lam y.\, x\;y)[x := y]
\]

In English: we're substituting $y$ for $x$ inside the term $\lam y.\, x\;y$.

\textbf{Step 1: Identify which rule applies.}
The outermost structure is $\lam y.\, (\ldots)$, and we're substituting for $x$
(not $y$), so rules 1--3 don't apply. We need either rule 4 or 5. Check: is
$y \in \FV(a)$? Here $a = y$, so $\FV(a) = \{y\}$. Yes, $y \in \FV(a)$. This
means rule 5 applies---we're in the dangerous case.

\textbf{Step 2: Pick a fresh name.}
We need $z$ that doesn't appear free in $e$, $a$, or equal $x$. We have
$\FV(e) = \FV(x\;y) = \{x, y\}$, $\FV(a) = \{y\}$, and $x$ is the substitution
variable. So we need $z \notin \{x, y\}$. Pick $z$.

\textbf{Step 3: Rename the binder.}
$\alpha$-convert $\lam y.\, x\;y$ to $\lam z.\, x\;z$ (replacing the bound $y$
with $z$ in the body):
\[
  (\lam y.\, x\;y) \;\aeq\; (\lam z.\, x\;z)
\]

\textbf{Step 4: Now substitute safely.}
\[
  (\lam z.\, x\;z)[x := y] \;=\; \lam z.\, (x\;z)[x := y]
\]
Now apply rule 6 (application): $(x\;z)[x := y] = x[x := y]\;z[x := y] =
y \cdot z$ (by rules 1 and 2). So:
\[
  (\lam y.\, x\;y)[x := y] \;=\; \lam z.\, y\;z
\]

The free $y$ in the result is still free---it refers to whatever $y$ was in the
surrounding context. If we had skipped the renaming, we would have gotten
$\lam y.\, y\;y$ (the self-application function), which is completely wrong.

\begin{intuition}[Why these rules exist]
The six rules encode a simple principle: substitution should respect scope.
Rules 1--2 handle variables. Rule 6 handles applications (which have no scope).
The interesting rules are 3--5, which handle binders:

Rule 3 says: if the binder re-introduces the same variable we're substituting
for, the outer variable is no longer visible here. This is \textbf{shadowing},
exactly like local variables in any programming language.

Rules 4--5 say: if the binder introduces a \emph{different} variable, we can
go inside---but only if it's safe. The safety check ($y \notin \FV(a)$) prevents
variable capture. If the check fails, we rename first.

The deep reason these rules are necessary is that lambda calculus uses
\textbf{names} to connect variables to their binders. Names are convenient for
humans but create a bookkeeping problem: different binders can accidentally use
the same name, and substitution can introduce new name collisions. De Bruijn
indices (Section~\ref{sec:debruijn}) eliminate this problem entirely by using
positions instead of names---which is why most implementations prefer them.
\end{intuition}

\subsubsection{A Flowchart for Substitution}

When computing $e[x := a]$, follow this decision process:

\begin{enumerate}[leftmargin=2em]
\item \textbf{Is $e$ a variable?}
  \begin{itemize}
  \item Is it $x$? $\longrightarrow$ Replace with $a$. \hfill(Rule 1)
  \item Is it something else? $\longrightarrow$ Leave it alone. \hfill(Rule 2)
  \end{itemize}
\item \textbf{Is $e$ an application $e_1\;e_2$?}
  $\longrightarrow$ Substitute into both sides: $(e_1[x:=a])\;(e_2[x:=a])$.
  \hfill(Rule 6)
\item \textbf{Is $e$ a $\lam$-abstraction $\lam y.\, \mathit{body}$?}
  \begin{itemize}
  \item Is $y = x$? $\longrightarrow$ Stop. The $\lam$ shadows $x$. \hfill(Rule 3)
  \item Is $y \neq x$ and $y \notin \FV(a)$?
    $\longrightarrow$ Go inside: $\lam y.\, (\mathit{body}[x := a])$.
    \hfill(Rule 4)
  \item Is $y \neq x$ and $y \in \FV(a)$?
    $\longrightarrow$ \textbf{Danger!} Rename $y$ to a fresh $z$ first,
    \emph{then} go inside. \hfill(Rule 5)
  \end{itemize}
\end{enumerate}

\subsection{Worked Reductions}

Let's work through several beta reductions step by step, starting from the
simplest and building up.

\begin{workedexample}[The simplest reduction]
\[
  (\lam x.\, x)\; 5 \;\;\reduce\;\; 5
\]
The identity function $\lam x.\, x$ applied to $5$: replace $x$ with $5$ in
the body ($x$), giving $5$. The function returns its argument unchanged.
\end{workedexample}

\begin{workedexample}[Identity applied to identity]
\[
  (\lam x.\, x)\;(\lam y.\, y)
    \;\reduce\; \lam y.\, y
    \rmark{substitute $(\lam y.\, y)$ for $x$ in body $x$}
\]
The argument is itself a function---that's fine! In lambda calculus, functions
are values like any other. The result is just $\lam y.\, y$ (the identity
function), returned unchanged.
\end{workedexample}

\begin{workedexample}[The $\mathbf{K}$ combinator discards its second argument]
\begin{align*}
  (\lam x.\, \lam y.\, x)\; a\; b
    &\reduce (\lam y.\, a)\; b  \rmark{substitute $a$ for $x$; the inner $\lam y$ remains} \\
    &\reduce a                  \rmark{substitute $b$ for $y$ in body $a$; but $y \notin \FV(a)$, so nothing changes}
\end{align*}
The $\mathbf{K}$ combinator takes two arguments and returns the first,
discarding the second. Notice that the second reduction ``does nothing''
because $a$ doesn't mention $y$---the argument $b$ is simply thrown away.
\end{workedexample}

\begin{workedexample}[A multi-step reduction: $\mathbf{S}\;\mathbf{K}\;\mathbf{K} = \mathbf{I}$]
$\mathbf{S} = \lam f\, g\, x.\, f\;x\;(g\;x)$. Let's apply it to $\mathbf{K}$ twice:
\begin{align*}
  \mathbf{S}\;\mathbf{K}\;\mathbf{K}
    &= (\lam f\, g\, x.\, f\;x\;(g\;x))\;\mathbf{K}\;\mathbf{K}
    \rmark{expand $\mathbf{S}$} \\
    &\reduce (\lam g\, x.\, \mathbf{K}\;x\;(g\;x))\;\mathbf{K}
    \rmark{plug in first $\mathbf{K}$ for $f$} \\
    &\reduce \lam x.\, \mathbf{K}\;x\;(\mathbf{K}\;x)
    \rmark{plug in second $\mathbf{K}$ for $g$}
\end{align*}
Now simplify the body. Recall $\mathbf{K}\;x\;(\mathbf{K}\;x)$: $\mathbf{K}$
takes two arguments and returns the first. So $\mathbf{K}\;x\;(\text{anything})
\reduces x$:
\[
  \lam x.\, \mathbf{K}\;x\;(\mathbf{K}\;x)
  \;\reduces\; \lam x.\, x \;=\; \mathbf{I}
\]
We've derived the identity function from $\mathbf{S}$ and $\mathbf{K}$ alone!
\end{workedexample}

\begin{workedexample}[Non-termination: $\Omega$]
Now something surprising. Consider the term $\lam x.\, x\;x$: a function that
applies its argument \emph{to itself}. What happens when we apply this to a copy
of itself?
\begin{align*}
  \Omega &= (\lam x.\, x\;x)\;(\lam x.\, x\;x) \\
         &\reduce (\lam x.\, x\;x)\;(\lam x.\, x\;x)
         \rmark{substitute $(\lam x.\, x\;x)$ for $x$ in $x\;x$} \\
         &= \Omega
\end{align*}
We got right back where we started! $\Omega$ reduces to itself, forever. This
is the $\lam$-calculus equivalent of an \textbf{infinite loop}. It shows that
not all lambda terms terminate---which makes sense, since any system powerful
enough to express all computable functions must also be able to express
non-terminating ones.
\end{workedexample}

\subsection{Normal Forms: When Computation Is ``Done''}

When you evaluate $(\lam x.\, x)\; 5$ and get $5$, you're done---there's
nothing left to reduce. But when you try $\Omega$, the reduction never ends.
How do we formalize the difference?

A \textbf{redex} is any sub-expression with the shape $(\lam x.\, e)\; a$---a
function directly applied to an argument. A term in \textbf{normal form}
contains \emph{no} redexes anywhere inside it. There is nowhere left to apply
the $\beta$-rule, so computation has finished.

Some examples:
\begin{enumerate}[leftmargin=2em]
\item $x$ is in normal form (just a variable, no redex).
\item $\lam x.\, x$ is in normal form (a function, but not applied to anything).
\item $x\; y$ is in normal form (an application, but $x$ is not a $\lam$).
\item $(\lam x.\, x)\; y$ is \emph{not} in normal form---it contains a redex.
\item $\Omega$ is \emph{not} in normal form, and it never reaches one.
\end{enumerate}

There is also a weaker notion called \textbf{weak head normal form} (WHNF): a
term is in WHNF if its outermost structure is either a $\lam$-abstraction or an
application whose leftmost part is not a redex. Most practical implementations
(including Haskell's runtime) stop at WHNF because fully normalizing under
lambdas is expensive and usually unnecessary.

\begin{theorem}[Church--Rosser / Confluence]
If $e \reduces e_1$ and $e \reduces e_2$ (possibly via different reduction
paths), then there exists some $e'$ such that $e_1 \reduces e'$ and
$e_2 \reduces e'$. In particular, \textbf{if a normal form exists, it is
unique} up to $\aeq$.
\end{theorem}

\begin{intuition}[What Church--Rosser means]
This theorem says: it doesn't matter what \emph{order} you reduce things in.
You might reduce the left sub-expression first, or the right one first, or
something deep inside the term. As long as you \emph{eventually} reach a normal
form, it will be \textbf{the same} normal form regardless of the path you took.
Different paths may take different numbers of steps (and some paths may diverge
entirely), but they can never reach \emph{different} normal forms. This is a
very reassuring property---it means the ``answer'' is well-defined.

Alonzo Church and J. Barkley Rosser proved this in 1936. It was a crucial
result: without it, the $\lam$-calculus would be ambiguous as a model of
computation, since different evaluation orders could yield different ``answers.''
The proof was notoriously difficult---Church himself called it one of the
hardest things he had worked on---and it took decades for truly elegant proofs
to emerge (notably Tait and Martin-L\"of's parallel reduction technique in the
1970s).
\end{intuition}

\begin{exercisebox}[Exercises: Beta Reduction]
\stepcounter{exercise}\textbf{\theexercise.}\; Reduce $(\lam x.\, x\;x)\;(\lam y.\, y)$ to normal form.

\stepcounter{exercise}\textbf{\theexercise.}\; Reduce $(\lam f.\, \lam x.\, f\;(f\;x))\;(\lam y.\, y\;y)$ step by step. Does it terminate?

\stepcounter{exercise}\textbf{\theexercise.}\; How many redexes are in the term $(\lam x.\, x)\;((\lam y.\, y)\;z)$? List them.

\stepcounter{exercise}\textbf{\theexercise.}\; Is $\lam x.\, (\lam y.\, y)\;x$ in normal form? If not, reduce it. Can you $\eta$-reduce the result?

\stepcounter{exercise}\textbf{\theexercise.}\; Perform the substitution $(\lam y.\, x\;y)[x := \lam y.\, y]$ using the capture-avoiding rules. Which rule (1--6) applies at each step?
\end{exercisebox}

%======================================================================

\subsection{Implementing It: Substitution and Evaluation}\label{sec:impl-subst}

Now let's implement the substitution and evaluation machinery we just studied.

\subsubsection{Substitution: Where It Gets Interesting}

Now the hard part. We need $e[x := a]$---the six rules from
Section~\ref{sec:beta}. The first two rules (variables) and the last rule
(application) are straightforward. The challenge is the lambda cases.

\begin{exercisebox}[The key challenge]
\stepcounter{exercise}\textbf{\theexercise.}\; Try writing \texttt{subst (x : String) (s : Term) : Term → Term} yourself. Start with the skeleton---three branches for three constructors. Handle the variable and application cases first (they're direct translations of Rules 1, 2, 6). Then think carefully about the lambda case ($\lam y.\, e)[x := s]$: there are three sub-cases (Rules 3, 4, 5). What conditions distinguish them?
\end{exercisebox}

Before looking at the answer, think through the lambda case. You have
$(\lam y.\, e)[x := s]$ and three questions to ask, \emph{in order}:

\begin{enumerate}[leftmargin=2em]
\item Is $y = x$? If yes, the $\lam$ shadows the variable we're substituting
  for. Stop. (Rule 3.)
\item Is $y \neq x$ and $y \notin \FV(s)$? If yes, it's safe to go inside.
  (Rule 4.)
\item Otherwise $y \in \FV(s)$---danger! Rename $y$ to something fresh before
  going inside. (Rule 5.)
\end{enumerate}

These three conditions map directly to an \texttt{if}/\texttt{else if}/\texttt{else} chain. You also need a helper to generate fresh names
that don't collide with anything in scope.

Here is the implementation. Read it against the substitution rules---each
\texttt{if}/\texttt{else} branch corresponds to one rule:

\begin{minted}{lean}
def freshVar (avoid : List String) (base : String := "x") : String :=
  let rec go (n : Nat) : String :=
    let candidate := s!"{base}{n}"
    if candidate ∈ avoid then go (n + 1) else candidate
  if base ∉ avoid then base else go 0

def subst (x : String) (s : Term) : Term → Term
  | .var y     => if y == x then s else .var y         -- Rules 1,2
  | .app e₁ e₂ => .app (subst x s e₁) (subst x s e₂) -- Rule 6
  | .lam y e   =>
    if y == x then .lam y e                           -- Rule 3: shadowed
    else if y ∉ freeVars s then
      .lam y (subst x s e)                            -- Rule 4: safe
    else
      let z := freshVar (freeVars s ++ freeVars e ++ [x])
      .lam z (subst x s (subst y (.var z) e))          -- Rule 5: rename
\end{minted}

The \texttt{freshVar} function generates a name that doesn't collide with
anything in scope. This is the ``pick a fresh $z$'' step from Rule~5.

\begin{intuition}[Substitution is not ``calling the function'']
A common early mistake is to test substitution by writing something like
\texttt{subst "x" (.var "a") K}, where K $= \lam x.\, \lam y.\, x$, and
expecting the result to be $\lam y.\, a$ (``K applied to a''). But the
result is K unchanged---because $x$ is \emph{not free} in K. It's bound by K's
own outermost $\lam$, so the substitution hits Rule 3 (shadowing) and stops.

This isn't a bug. Substitution is a textual operation: it replaces \emph{free}
occurrences of a variable. To actually compute ``K applied to a,'' you need
\textbf{beta reduction}, which first \emph{strips} K's outer $\lam$ and then
substitutes into the \emph{body}:
\[
  (\lam x.\, \lam y.\, x)\;a
  \;\;\reduce\;\;
  (\lam y.\, x)[x := a]
  \;\;=\;\;
  \lam y.\, a
\]
Beta reduction calls substitution as one of its internal steps, but they are
different operations. We'll implement beta reduction next, in the evaluator.
\end{intuition}

%======================================================================
\subsection{Evaluation}\label{sec:evalimpl}
%======================================================================

\subsubsection{Single-Step Reduction}

You have terms and substitution. Now: how do you actually \emph{evaluate}?

The core idea: scan the term for a redex---an application whose left side is a
$\lam$. When you find one, fire the $\beta$-rule.

\begin{exercisebox}[Design challenge]
\stepcounter{exercise}\textbf{\theexercise.}\; A \texttt{betaStep} function should find \emph{one} redex and reduce it, returning the updated term. If no redex exists, return \texttt{none}. Think about: what order should you search for redexes? What if the redex is nested deep inside the term?
\end{exercisebox}

Before coding, think about what this function must do by case analysis on the
term's shape:

\begin{enumerate}[leftmargin=2em]
\item \textbf{Variable.} No redex here---a variable isn't an application. Return
  \texttt{none}.
\item \textbf{Application $e_1\;e_2$.} This is where the action is. Ask: is
  $e_1$ a $\lam$? If yes, you found a redex---fire the rule! If no, try to
  reduce $e_1$ recursively. If $e_1$ is already stuck, try $e_2$.
\item \textbf{Lambda $\lam x.\, e$.} Not a redex itself, but there might be one
  inside the body $e$. Try recursing into it.
\end{enumerate}

Notice how the application case has a special check at the top (``is the left
side a $\lam$?'') before falling through to the recursive cases. This is what
makes the evaluation strategy \textbf{normal order}---we try the outermost
redex first.

Here is normal-order evaluation (leftmost outermost first):

\begin{minted}{lean}
def betaStep : Term → Option Term
  | .app (.lam x body) arg => some (subst x arg body)  -- found a redex!
  | .app e₁ e₂ =>
    match betaStep e₁ with                           -- try left side
    | some e₁' => some (.app e₁' e₂)
    | none     => match betaStep e₂ with             -- then right side
                  | some e₂' => some (.app e₁ e₂')
                  | none     => none
  | .lam x e => match betaStep e with                 -- reduce under λ
               | some e' => some (.lam x e')
               | none    => none
  | _ => none                                        -- variable: nothing to do
\end{minted}

The first branch is the redex case---the $\beta$-rule in action. The remaining
branches recursively search for a redex deeper inside the term.

\subsubsection{Multi-Step with Fuel}

One step isn't enough---we want to keep reducing until we reach normal form (or
give up). But here's a problem: lambda calculus evaluation might
\textbf{not terminate} ($\Omega$!). Lean requires all functions to terminate.

\begin{exercisebox}[Termination challenge]
\stepcounter{exercise}\textbf{\theexercise.}\; How do you write a ``keep reducing until done'' function in a language that demands termination? Think about what parameter you could add that guarantees the loop eventually stops.
\end{exercisebox}

The standard trick: a \textbf{fuel} parameter that decreases every step.

\begin{minted}{lean}
def eval (fuel : Nat) (t : Term) : Term :=
  match fuel with
  | 0         => t                       -- out of fuel: return as-is
  | fuel' + 1 => match betaStep t with
                 | none    => t          -- normal form reached!
                 | some t' => eval fuel' t'
\end{minted}

Let's test it:

\begin{minted}{lean}
-- (λx. x)(λy. y) →β (λy. y)
#eval eval 100 (.app (.lam "x" (.var "x")) (.lam "y" (.var "y")))
-- Result: (λy. y)

-- K a b = (λx. λy. x) a b →β a
#eval eval 100
  (.app (.app (.lam "x" (.lam "y" (.var "x"))) (.var "a")) (.var "b"))
-- Result: a
\end{minted}

It works! We have a complete $\lam$-calculus interpreter. But there's a problem
lurking beneath the surface\ldots

%======================================================================
\section{Evaluation Strategies}\label{sec:strategies}
%======================================================================

\subsection{The Problem: Which Redex First?}

A lambda term can contain more than one redex at the same time. For example,
consider the expression $(\lam x.\, x)\;((\lam y.\, y)\; z)$. There are two
redexes here:
\begin{enumerate}[leftmargin=2em]
\item The \textbf{outer} redex: the whole expression is $(\lam x.\, x)$ applied
  to $(\lam y.\, y)\; z$.
\item The \textbf{inner} redex: inside the argument, $(\lam y.\, y)\; z$ is
  itself a function applied to an argument.
\end{enumerate}
Which one do we reduce first? The Church--Rosser theorem guarantees that if we
eventually reach a normal form, it will be the same one regardless of order.
But the \emph{path} we take can matter enormously---some orders terminate while
others don't!

An \textbf{evaluation strategy} is a rule for choosing which redex to reduce
next. Different strategies correspond to different programming language designs.

\subsection{The Main Strategies}

\paragraph{Normal order (reduce the leftmost outermost redex first).}
``Outermost'' means we reduce the biggest enclosing redex before diving into its
parts. In programming terms, we pass arguments into functions \emph{before}
evaluating those arguments. This is the strategy behind \textbf{lazy}
evaluation.

\paragraph{Applicative order (reduce the leftmost innermost redex first).}
``Innermost'' means we fully evaluate arguments \emph{before} passing them into
a function. This is how most familiar languages work: in \texttt{f(g(3))},
Python evaluates \texttt{g(3)} first, then passes the result to \texttt{f}.
This is \textbf{eager} (or strict) evaluation.

\paragraph{Call-by-need (normal order with sharing).}
Like normal order, but if an argument is used multiple times, it's evaluated
only once and the result is shared. This gives the termination benefits of
normal order without redundant computation. This is what Haskell does.

\begin{center}
\renewcommand{\arraystretch}{1.5}
\begin{tabular}{lll}
\textbf{Strategy} & \textbf{Rule} & \textbf{Languages} \\
\hline
Normal order     & Leftmost outermost first & Lazy evaluation \\
Applicative order & Leftmost innermost first & JS, Python, C, Java, Lean \\
Call-by-need     & Normal order + memoization & Haskell \\
\end{tabular}
\end{center}

\subsection{Why It Matters: A Life-or-Death Example}

\begin{workedexample}[Strategy matters: $\mathbf{K}\;\mathbf{I}\;\Omega$]
Recall $\mathbf{K} = \lam x\, y.\, x$ (takes two arguments, returns the
first). Consider $\mathbf{K}\;\mathbf{I}\;\Omega$, where $\Omega$ is the
non-terminating term from Section~\ref{sec:beta}.

\textbf{Normal order} reduces the outermost redex first---the $\mathbf{K}$
applied to $\mathbf{I}$:
\begin{align*}
  (\lam x\, y.\, x)\;\mathbf{I}\;\Omega
  \;\reduce\; (\lam y.\, \mathbf{I})\;\Omega
  \;\reduce\; \mathbf{I} \quad\checkmark
\end{align*}
$\mathbf{K}$ discards its second argument, so $\Omega$ is thrown away
\emph{before it's ever evaluated}. Done in two steps.

\textbf{Applicative order} tries to evaluate the arguments first, including
$\Omega$:
\[
  (\lam x\, y.\, x)\;\mathbf{I}\;\underbrace{\Omega}_{\text{evaluate this first!}}
  \;\longrightarrow\; (\lam x\, y.\, x)\;\mathbf{I}\;\Omega
  \;\longrightarrow\; \cdots
  \quad\text{(loops forever!)}
\]
It gets stuck trying to evaluate $\Omega$ and never finishes. The program
\emph{hangs}, even though the answer ($\mathbf{I}$) is perfectly well-defined
and doesn't depend on $\Omega$ at all.
\end{workedexample}

This example demonstrates a fundamental theorem:

\begin{theorem}[Standardization]
If a term has a normal form, normal-order reduction will find it. Applicative
order may diverge even when a normal form exists.
\end{theorem}

This theorem was proved by Curry and Feys in 1958 and sharpened
by Plotkin in 1975, who formalized the precise relationship between
different evaluation strategies and their corresponding programming language
semantics. Plotkin's paper ``Call-by-name, call-by-value, and the
$\lambda$-calculus'' is one of the most cited papers in programming language
theory.

\begin{intuition}[The trade-off]
If normal order is ``safer'' (it always finds a normal form when one exists),
why don't all languages use it? The trade-off is \textbf{efficiency}. Under
applicative order, each argument is evaluated exactly once. Under normal order,
if an argument is used three times in the body, it might be evaluated three
times. Call-by-need (Haskell's approach) gets the best of both worlds: lazy
semantics with no redundant work---but at the cost of implementation complexity.
\end{intuition}

%======================================================================
\section{Eta Conversion ($\eta$)}\label{sec:eta}
%======================================================================

We have seen two operations so far: $\alpha$-conversion (renaming) and
$\beta$-reduction (function calling). There is one more:
$\eta$-conversion (pronounced ``EH-tah''), which captures a subtler
equivalence.

\subsection{The Core Idea: Pointless Wrappers}

Consider this JavaScript code:

\begin{minted}{javascript}
const f = x => Math.sqrt(x);
\end{minted}

This creates a function that takes \texttt{x} and passes it to
\texttt{Math.sqrt}. But \texttt{Math.sqrt} is \emph{already} a function that
takes one argument! So the wrapper is redundant---you could just write:

\begin{minted}{javascript}
const f = Math.sqrt;
\end{minted}

These two definitions behave identically on all inputs. Eta conversion is the
formal rule that says these are the same.

\begin{defnbox}[The $\eta$-conversion Rule]
\[
  \lam x.\, f\; x \;\;\ef\;\; f
  \qquad\text{provided } x \notin \FV(f)
\]
\end{defnbox}

The condition $x \notin \FV(f)$ is essential. If $f$ mentions $x$ freely,
then $\lam x.\, f\; x$ is not just a wrapper around $f$---the $\lam$
actually binds an $x$ that $f$ uses. For example, $\lam x.\, (x + 1)\; x$
cannot be $\eta$-reduced because the body uses $x$ in two places.

\begin{intuition}[When to use $\eta$-conversion]
$\eta$-conversion expresses \textbf{extensionality}: if two functions produce
the same output for every input, they are the same function. Most functional
programmers encounter this as the ``remove pointless wrappers'' refactoring:
\texttt{x => f(x)} simplifies to \texttt{f}. In Haskell, this is called
\textbf{eta reduction} and is a common style recommendation
(\texttt{map (\textbackslash x -> f x) xs} simplifies to \texttt{map f xs}).
\end{intuition}

%======================================================================
\section{Church Encodings}\label{sec:church}
%======================================================================

\subsection{The Big Idea: Everything Is a Function}

Lambda calculus has no built-in data types. No booleans, no numbers, no strings,
no lists. There is \emph{nothing} except variables, function definitions, and
function calls. So how could this system possibly be useful for computation?

The answer is Church's remarkable insight: you don't \emph{need} built-in data.
You can \emph{encode} any data type as a pure function, by representing each
value as a function that describes ``what you do with it.'' A boolean
is a function that \emph{chooses}. A number is a function that
\emph{repeats}. A pair is a function that \emph{holds two things and gives
you access to them}.

These function-based representations are called \textbf{Church encodings}, and
they are one of the most beautiful ideas in computer science. They show that
data and behavior are two sides of the same coin.

Church developed these encodings out of necessity. To prove that his
$\lam$-calculus could serve as a foundation for mathematics---and specifically
to show it could express arithmetic, which was needed for his attack on
Hilbert's Entscheidungsproblem---he had to demonstrate that numbers, logical
operations, and recursive definitions could all be expressed using nothing but
functions. The encodings were his proof that the system was powerful enough.
His student Stephen Kleene then extended the work, showing how to encode more
complex operations (including the predecessor function, as we'll see).

In the definitions below, the symbol $\triangleq$ means ``is defined as''
(see Section~\ref{sec:notation}). We use it instead of $=$ to emphasize that
we are \emph{introducing a name} for a term, not stating an equation we
discovered.

\subsection{Church Booleans}

What is a boolean? At its core, it's something that lets you choose between
two alternatives. ``If true, do this; if false, do that.'' A Church boolean
\emph{is} that choice, expressed directly as a function.

\begin{defnbox}[Church Booleans]
\vspace{-0.8em}
\begin{align*}
\TR   &\triangleq \lam t\, f.\, t \rmark{select first} &
\FL  &\triangleq \lam t\, f.\, f \rmark{select second}
\end{align*}
\end{defnbox}

Both $\TR$ and $\FL$ are functions that take two arguments, $t$ and $f$
(short for ``true branch'' and ``false branch''). $\TR$ returns the first one;
$\FL$ returns the second one.

\begin{intuition}[Data as behavior]
In most languages, \texttt{if} is a language construct and booleans are
primitive values. In Church encoding, there is no \texttt{if}---the boolean
itself \emph{is} the conditional. When you ``apply a boolean to two choices,''
it picks one. There is no need for a separate \texttt{if} statement because
the boolean already knows how to choose.

This is the fundamental insight of Church encodings: \textbf{data is
represented by what you do with it}, not by some internal structure.
\end{intuition}

Using these booleans, we can define the standard logical operations. Each one
works by cleverly applying a boolean to the right pair of arguments:

\vspace{-0.5em}
\begin{align*}
\textsc{And}  &\triangleq \lam p\, q.\, p\; q\; \FL &
\textsc{Or}   &\triangleq \lam p\, q.\, p\; \TR\; q \\
\textsc{Not}  &\triangleq \lam p.\, p\; \FL\; \TR &
\textsc{If}   &\triangleq \lam p\, a\, b.\, p\; a\; b
\end{align*}

How do these work? Take $\textsc{And}$: it applies $p$ to two arguments, $q$
and $\FL$. If $p$ is $\TR$, it selects the first argument ($q$), so the result
depends on $q$---exactly right for ``and.'' If $p$ is $\FL$, it selects the
second argument ($\FL$), giving false immediately---also correct.

$\textsc{Not}$ swaps the two arguments: it gives a boolean $\FL$ as its ``true
branch'' and $\TR$ as its ``false branch,'' so $\TR$ becomes $\FL$ and vice
versa.

$\textsc{If}$ is the most elegant---it's literally just function application!
Applying a Church boolean to two choices \emph{is} an if-then-else.

\begin{workedexample}[$\textsc{And}\;\TR\;\FL = \FL$]
Let's verify step by step that ``true and false = false'':
\begin{align*}
  \textsc{And}\;\TR\;\FL
  &= (\lam p\, q.\, p\; q\; \FL)\;\TR\;\FL
    \rmark{expand the definition of \textsc{And}} \\
  &\reduce (\lam q.\, \TR\; q\; \FL)\;\FL
    \rmark{substitute $\TR$ for $p$} \\
  &\reduce \TR\;\FL\;\FL
    \rmark{substitute $\FL$ for $q$} \\
  &= (\lam t\, f.\, t)\;\FL\;\FL
    \rmark{expand the definition of $\TR$} \\
  &\reduce (\lam f.\, \FL)\;\FL
    \rmark{$\TR$ selects its first argument: $\FL$} \\
  &\reduce \FL \quad\checkmark
    \rmark{the remaining $\FL$ argument is discarded}
\end{align*}
\end{workedexample}

\subsection{Church Numerals}\label{sec:churchnums}

How do you represent numbers as functions? Church's idea: a number $n$ is a
function that \emph{applies another function $n$ times}. The number 0 applies
$f$ zero times (just returns the starting value). The number 1 applies $f$
once. The number 3 applies $f$ three times. The number \emph{is} the
repetition count.

\begin{defnbox}[Church Numerals]
$\church{n}$ applies its first argument $n$ times to its second:
\[
\church{0} = \lam f\, x.\, x
\qquad
\church{1} = \lam f\, x.\, f\;x
\qquad
\church{n} = \lam f\, x.\, \underbrace{f\;(f\;(\cdots(f}_{n}\; x)\cdots))
\]
\end{defnbox}

\begin{intuition}[Reading Church numeral notation]
The overline $\church{n}$ is shorthand for ``the Church encoding of number
$n$.'' The notation $f^n\;x$ means ``apply $f$ a total of $n$ times to $x$'':
$f^0\;x = x$, $f^1\;x = f\;x$, $f^2\;x = f\;(f\;x)$, etc. So $\church{n}$
is a function that takes $f$ and $x$ and applies $f$ exactly $n$ times. The
number \emph{is} the iteration.
\end{intuition}

Writing out the first few concretely:
\begin{align*}
\church{0} &= \lam f\, x.\, x &\rmark{apply $f$ zero times: just return $x$} \\
\church{1} &= \lam f\, x.\, f\;x &\rmark{apply $f$ once} \\
\church{2} &= \lam f\, x.\, f\;(f\;x) &\rmark{apply $f$ twice} \\
\church{3} &= \lam f\, x.\, f\;(f\;(f\;x)) &\rmark{apply $f$ three times}
\end{align*}

Notice something: $\church{0} = \lam f\, x.\, x$ has the same structure as
$\FL$ (select the second argument). This is not a coincidence---zero and
``false'' are deeply connected in Church encoding.

Now we can define arithmetic. Each operation works by manipulating how many
times $f$ gets applied:

\begin{align*}
\textsc{Succ}  &\triangleq \lam n\, f\, x.\, f\;(n\; f\; x) &
\textsc{Add}   &\triangleq \lam m\, n\, f\, x.\, m\; f\;(n\; f\; x) \\
\textsc{Mul}   &\triangleq \lam m\, n\, f.\, m\;(n\; f) &
\textsc{Exp}   &\triangleq \lam m\, n.\, n\; m
\end{align*}

How does $\textsc{Succ}$ (successor, i.e., ``add 1'') work? It takes a Church
numeral $n$ and returns a new function that applies $f$ one extra time:
first do $n\;f\;x$ (apply $f$ a total of $n$ times), then wrap one more $f$
around the result.

$\textsc{Add}$ works similarly: apply $f$ a total of $n$ times (via
$n\;f\;x$), then apply $f$ another $m$ times on top (via $m\;f\;(\ldots)$).
The total is $m + n$ applications of $f$.

$\textsc{Mul}$ is more subtle: $n\;f$ is a function that applies $f$ exactly
$n$ times. Then $m\;(n\;f)$ applies \emph{that combined function} $m$ times.
$n$ applications repeated $m$ times gives $m \times n$ applications total.

\begin{workedexample}[$\textsc{Succ}\;\church{2} = \church{3}$]
\begin{align*}
  \textsc{Succ}\;\church{2}
  &= (\lam n\, f\, x.\, f\;(n\; f\; x))\;\church{2}
    \rmark{expand $\textsc{Succ}$} \\
  &\reduce \lam f\, x.\, f\;\bigl(\church{2}\; f\; x\bigr)
    \rmark{plug in $\church{2}$ for $n$} \\
  &\reduces \lam f\, x.\, f\;(f\;(f\;x))
    \rmark{$\church{2}\;f\;x = f\;(f\;x)$; wrap one more $f$} \\
  &= \church{3} \quad\checkmark
\end{align*}
\end{workedexample}

\subsection{Church Pairs}

A pair holds two values and lets you extract either one. In Church encoding, a
pair is a function that ``closes over'' its two elements and waits for you to
tell it which one you want---by passing it either $\TR$ (to get the first) or
$\FL$ (to get the second).

\begin{align*}
\textsc{Pair} &\triangleq \lam a\, b\, f.\, f\; a\; b &
\textsc{Fst} &\triangleq \lam p.\, p\; \TR &
\textsc{Snd} &\triangleq \lam p.\, p\; \FL
\end{align*}

$\textsc{Pair}\;a\;b$ creates a function $\lam f.\, f\;a\;b$ that remembers
$a$ and $b$ in its closure. To extract the first element, $\textsc{Fst}$ passes
$\TR$ (which selects the first of two arguments) as the extraction function. To
extract the second, $\textsc{Snd}$ passes $\FL$.

\begin{workedexample}[$\textsc{Fst}\;(\textsc{Pair}\;x\;y) = x$]
\begin{align*}
  \textsc{Fst}\;(\textsc{Pair}\;x\;y)
  &\reduces (\lam p.\, p\;\TR)\;(\lam f.\, f\;x\;y)
    \rmark{$\textsc{Pair}\;x\;y$ becomes $\lam f.\, f\;x\;y$} \\
  &\reduce (\lam f.\, f\;x\;y)\;\TR
    \rmark{pass the pair its extractor ($\TR$)} \\
  &\reduce \TR\;x\;y
    \rmark{the pair hands both elements to $\TR$} \\
  &\reduces x \quad\checkmark
    \rmark{$\TR$ selects the first: $x$}
\end{align*}
\end{workedexample}

\subsection{IsZero and Predecessor}\label{sec:pred}

With booleans and numerals defined, we need a way to \emph{test} whether a
number is zero. This turns out to be surprisingly elegant:

\[
\textsc{IsZero} \triangleq \lam n.\, n\; (\lam \_.\, \FL)\; \TR
\]

How does it work? Recall that a Church numeral $\church{n}$ applies its first
argument $n$ times to its second. If $n = 0$, it applies the function zero
times, just returning $\TR$. If $n \geq 1$, it applies $\lam\_.\,\FL$ at least
once---and that function ignores its argument and returns $\FL$ regardless.

\begin{workedexample}[$\textsc{IsZero}\;\church{0} = \TR$ and $\textsc{IsZero}\;\church{2} = \FL$]
\begin{align*}
  \textsc{IsZero}\;\church{0}
  &= (\lam n.\, n\;(\lam \_.\,\FL)\;\TR)\;(\lam f\,x.\,x)
  \;\reduce\; (\lam f\,x.\,x)\;(\lam \_.\,\FL)\;\TR
  \;\reduces\; \TR \;\checkmark \\[4pt]
  \textsc{IsZero}\;\church{2}
  &\reduces\; (\lam \_.\,\FL)\;\bigl((\lam \_.\,\FL)\;\TR\bigr)
  \;\reduces\; \FL \;\checkmark
\end{align*}
\end{workedexample}

\textbf{Predecessor} is famously the hardest basic Church encoding. Successor
is easy---just wrap one more $f$ around the result. But \emph{removing} one
layer of $f$ is not straightforward, because a Church numeral only knows how to
apply $f$ forwards; there is no built-in way to ``undo'' an application.

The solution was discovered by Stephen Kleene in 1932, when he was a PhD
student of Church's at Princeton. Legend has it that the insight came to him
while sitting in the dentist's chair. The problem had stumped Church's group
for months---without predecessor, they couldn't define subtraction, comparison,
or any of the recursive functions needed to show that $\lam$-calculus could
express all of arithmetic. Kleene's breakthrough was the ``pair-walking
trick'': iterate a pair-updating function $n$ times,
starting from a carefully chosen initial pair, so that after $n$ steps the
desired value sits in the first component.

\[
\textsc{Pred} \triangleq \lam n.\, \textsc{Fst}\;\bigl(
  n\;\Phi\;(\textsc{Pair}\;\church{0}\;\church{0})
\bigr)
\qquad\text{where}\quad
\Phi \triangleq \lam p.\, \textsc{Pair}\;(\textsc{Snd}\;p)\;(\textsc{Succ}\;(\textsc{Snd}\;p))
\]

The helper $\Phi$ takes a pair $(a, b)$ and produces $(\text{old } b,\;
b + 1)$. Starting from $(0, 0)$:

\begin{center}
\renewcommand{\arraystretch}{1.3}
\begin{tabular}{cccc}
\textbf{Step} & \textbf{Pair before} & $\Phi$ & \textbf{Pair after} \\
\hline
0 & --- & --- & $(0, 0)$ \\
1 & $(0, 0)$ & $(\text{snd}, \text{succ snd})$ & $(0, 1)$ \\
2 & $(0, 1)$ & & $(1, 2)$ \\
3 & $(1, 2)$ & & $(2, 3)$ \\
$n$ & $(n{-}2,\; n{-}1)$ & & $(n{-}1,\; n)$ \\
\end{tabular}
\end{center}

After $n$ applications of $\Phi$, the first component is $n - 1$. Taking
$\textsc{Fst}$ of the final pair gives the predecessor. (For $n = 0$, the pair
never changes from $(0, 0)$, so $\textsc{Pred}\;\church{0} = \church{0}$.)

\begin{intuition}[Why Pred is hard]
Successor adds a layer of $f$; predecessor must remove one. But Church numerals
are ``write-only''---they apply $f$ a fixed number of times and give you no way
to ``un-apply'' it. Kleene's trick works around this by never removing anything:
instead, it rebuilds the count from scratch using a pair that always trails one
step behind. This is why predecessor is $O(n)$ even though the number $n$ is
``already there'' in the encoding.
\end{intuition}

\subsection{Church-Encoded Lists}\label{sec:churchlists}

Lists round out the data structures story. A Church-encoded list is a function
that takes two arguments---a ``cons handler'' and a ``nil value''---and folds
itself. This is exactly the same pattern as \texttt{foldr} in Haskell or
\texttt{reduce} in JavaScript.

\begin{defnbox}[Church Lists]
\vspace{-0.8em}
\begin{align*}
\textsc{Nil}  &\triangleq \lam c\, n.\, n \\
\textsc{Cons} &\triangleq \lam h\, t\, c\, n.\, c\; h\; (t\; c\; n)
\end{align*}
\end{defnbox}

$\textsc{Nil}$ is a list that ignores the cons handler and returns the nil
value directly---like an empty list that has nothing to fold over.
$\textsc{Cons}\;h\;t$ is a list with head $h$ and tail $t$: when folded, it
applies the cons handler $c$ to $h$ and the result of folding the tail.

\begin{intuition}[Lists are their own fold]
In Haskell, \texttt{foldr f z [1,2,3]} computes \texttt{f 1 (f 2 (f 3 z))}.
A Church list \emph{is} this computation. The list
$\textsc{Cons}\;1\;(\textsc{Cons}\;2\;(\textsc{Cons}\;3\;\textsc{Nil}))$
is a function that, when given $c$ and $n$, computes
$c\;1\;(c\;2\;(c\;3\;n))$. The list doesn't store data passively and wait
for you to traverse it---it \emph{is} the traversal.
\end{intuition}

With this encoding, list operations become simple:

\begin{align*}
\textsc{Head}  &\triangleq \lam l.\, l\; (\lam h\, \_.\, h)\; \FL
  &\rmark{fold with ``take first''} \\
\textsc{IsNil} &\triangleq \lam l.\, l\; (\lam \_\, \_.\, \FL)\; \TR
  &\rmark{any cons returns $\FL$} \\
\textsc{Map}   &\triangleq \lam f\, l\, c\, n.\, l\; (\lam h\, t.\, c\; (f\; h)\; t)\; n
  &\rmark{apply $f$ in the cons handler} \\
\textsc{Append} &\triangleq \lam l_1\, l_2\, c\, n.\, l_1\; c\; (l_2\; c\; n)
  &\rmark{fold $l_1$ then $l_2$}
\end{align*}

$\textsc{Head}$ folds the list with a handler that returns the first element
it sees (ignoring the rest). $\textsc{IsNil}$ works like $\textsc{IsZero}$:
if no cons is ever applied, return $\TR$; otherwise $\FL$.
$\textsc{Append}$ is particularly elegant: fold the first list with $c$, but
instead of stopping at $n$, continue folding into the second list.

\begin{workedexample}[Extracting the head of a list]
Let $L = \textsc{Cons}\;a\;(\textsc{Cons}\;b\;\textsc{Nil})$.
\begin{align*}
  \textsc{Head}\;L
  &= L\; (\lam h\, \_.\, h)\; \FL \\
  &\reduces (\lam h\, \_.\, h)\; a\; \bigl((\lam h\, \_.\, h)\; b\; \FL\bigr)
  \;\reduce\; a \quad\checkmark
\end{align*}
\end{workedexample}

\begin{exercisebox}[Exercises: Church Encodings]
\stepcounter{exercise}\textbf{\theexercise.}\; Reduce $\textsc{IsZero}\;(\textsc{Succ}\;\church{0})$ step by step. What is the result?

\stepcounter{exercise}\textbf{\theexercise.}\; Verify that $\textsc{Add}\;\church{2}\;\church{3} = \church{5}$ by reducing it to normal form.

\stepcounter{exercise}\textbf{\theexercise.}\; Define $\textsc{Sub}$ (subtraction) using $\textsc{Pred}$:
$\textsc{Sub}\;\church{m}\;\church{n}$ should apply $\textsc{Pred}$ a total of $n$ times to $\church{m}$.

\stepcounter{exercise}\textbf{\theexercise.}\; Define $\textsc{Length}$ for Church lists. \emph{Hint:} fold with a handler that ignores the element and increments a counter.

\stepcounter{exercise}\textbf{\theexercise.}\; Is $\textsc{Nil}$ the same term as $\church{0}$? Explain why or why not, and what this says about Church encodings.
\end{exercisebox}

In most programming languages, writing a recursive function is straightforward:
you give the function a name and call that name inside its own body.

\begin{minted}{javascript}
function factorial(n) {
    if (n === 0) return 1;
    return n * factorial(n - 1);  // calls itself by name
}
\end{minted}

But lambda calculus has \emph{no names}. There is no way to say ``define a
function called \texttt{factorial} and refer to it from within itself.'' Every
function is anonymous---it's just $\lam x.\, \mathit{body}$, with no way for
the body to refer back to the function that contains it.

This seems like a fatal limitation. Without recursion, how can we express
loops? How can we compute factorials, traverse lists, or do anything that
requires repeating a computation?

\subsection{The Solution: Fixed-Point Combinators}\label{sec:recursion}

The trick is both ingenious and strange. Instead of a function calling itself
by name, we use a helper called a \textbf{fixed-point combinator} that
``feeds a function to itself'' automatically.

The key insight: suppose you write your recursive function in a special way.
Instead of calling yourself directly, you take an \emph{extra parameter}
called \texttt{self} and call \emph{that} whenever you would have made a
recursive call:

\begin{minted}{javascript}
// Instead of:  function fact(n) { ... fact(n-1) ... }
// Write:       self => n => ... self(n-1) ...
const F = self => n => (n === 0) ? 1 : n * self(n - 1);
\end{minted}

Now $F$ is not recursive---it doesn't call itself. But if we could somehow pass
$F$ a version of itself that ``works,'' we'd have recursion. A fixed-point
combinator does exactly this: given $F$, it produces a value $v$ such that
$v = F(v)$. That $v$ is the recursive function we wanted.

\begin{defnbox}[Fixed-Point Property]
$\mathbf{Y}$ is a \textbf{fixed-point combinator} if, for every function $f$:
\[
  \mathbf{Y}\; f \;=\; f\;(\mathbf{Y}\; f)
\]
In other words, $\mathbf{Y}\;f$ is a \textbf{fixed point} of $f$: a value
that, when you apply $f$ to it, gives you back the same value.
\end{defnbox}

\begin{intuition}[Why ``fixed point''?]
In mathematics, a fixed point of a function $f$ is a value $v$ where
$f(v) = v$. For example, 0 is a fixed point of the function $f(x) = x^2$
because $f(0) = 0$. The $\mathbf{Y}$ combinator finds a fixed point of
\emph{any} function---even functions on lambda terms. When $f$ is our
``factorial template'' $F$ from above, the fixed point
$\mathbf{Y}\;F$ is the actual factorial function: $F$ applied to it gives back
the same function.
\end{intuition}

\subsection{The Y Combinator}

The $\mathbf{Y}$ combinator was discovered by Haskell Curry in the 1940s,
though Church's group at Princeton had studied fixed-point constructions
earlier. Curry showed that \emph{every} function in the $\lam$-calculus has a
fixed point---a consequence of the system's ability to express self-application.
This was a double-edged sword: it made recursion possible, but it also meant
Church's original goal of using $\lam$-calculus as a consistent foundation for
\emph{logic} was doomed (Kleene and Rosser showed in 1935 that the untyped
$\lam$-calculus is logically inconsistent). The $\lam$-calculus survived by
pivoting: it became a theory of \emph{computation} rather than logic.

Here is the actual definition:

\begin{align*}
\mathbf{Y} &= \lam f.\, (\lam x.\, f\;(x\;x))\;(\lam x.\, f\;(x\;x))
  & \text{(normal order only)}
\end{align*}

This looks intimidating, but it's built from a familiar pattern. Remember
$\Omega = (\lam x.\, x\;x)\;(\lam x.\, x\;x)$, which loops forever because
it applies itself to itself? The $\mathbf{Y}$ combinator uses the same
self-application trick, but inserts $f$ so the self-application does something
useful (calling $f$ each time) instead of just spinning in circles.

\begin{workedexample}[$\mathbf{Y}\;g = g\;(\mathbf{Y}\;g)$]
Let's verify the fixed-point property by reducing $\mathbf{Y}\;g$ step by step:
\begin{align*}
  \mathbf{Y}\;g
  &= (\lam f.\, (\lam x.\, f\;(x\;x))\;(\lam x.\, f\;(x\;x)))\;g
    \rmark{expand $\mathbf{Y}$} \\
  &\reduce (\lam x.\, g\;(x\;x))\;(\lam x.\, g\;(x\;x))
    \rmark{substitute $g$ for $f$} \\
  &\reduce g\;\bigl(\underbrace{(\lam x.\, g\;(x\;x))\;(\lam x.\, g\;(x\;x))}_{\text{this is } \mathbf{Y}\;g \text{ again!}}\bigr) \\
  &= g\;(\mathbf{Y}\;g) \quad\checkmark
\end{align*}
The key moment is the last step: after two $\beta$-reductions, we get
$g$ applied to an expression that is identical to $\mathbf{Y}\;g$. So
$\mathbf{Y}\;g$ ``unrolls'' into $g\;(\mathbf{Y}\;g)$, which unrolls into
$g\;(g\;(\mathbf{Y}\;g))$, and so on---as many times as needed.
\end{workedexample}

\subsection{The Z Combinator (Call-by-Value)}

The $\mathbf{Y}$ combinator only works under normal-order (lazy) evaluation.
Under applicative-order (eager) evaluation---which is what most languages
use---$\mathbf{Y}\;f$ diverges because the self-application $x\;x$ is
evaluated immediately, before $f$ can ``decide'' to stop.

The fix is the \textbf{Z combinator}, which wraps the self-application in a
$\lam$ to delay it:

\[
\mathbf{Z} = \lam f.\, (\lam x.\, f\;(\lam y.\, x\;x\;y))\;(\lam x.\, f\;(\lam y.\, x\;x\;y))
\]

The extra $\lam y.\, \ldots\; y$ acts as a ``thunk''---it prevents $x\;x$
from being evaluated until the result is actually needed (i.e., until the
function is called with an argument $y$). This is the combinator used in
eager languages like JavaScript, Python, and our call-by-value Lean
interpreter.

\subsection{Factorial: Putting It All Together}

\begin{workedexample}[Factorial via $\mathbf{Y}$]
Define the factorial ``template'' (a non-recursive function that takes a
$\mathit{self}$ parameter):
\[
F = \lam \mathit{self}.\, \lam n.\,
\textsc{If}\;(\textsc{IsZero}\;n)\;\church{1}\;(\textsc{Mul}\;n\;(\mathit{self}\;(\textsc{Pred}\;n)))
\]
Then $\textsc{Fact} = \mathbf{Y}\;F$ is the actual factorial function. Tracing
$\textsc{Fact}\;\church{3}$:
\begin{align*}
\textsc{Fact}\;\church{3}
  &= (\mathbf{Y}\;F)\;\church{3}
   = F\;(\mathbf{Y}\;F)\;\church{3}
    \rmark{fixed-point property: $\mathbf{Y}\;F = F\;(\mathbf{Y}\;F)$} \\
  &\reduces \textsc{Mul}\;\church{3}\;(\textsc{Fact}\;\church{2})
    \rmark{$\church{3}$ is not zero, so take the else branch} \\
  &\reduces \textsc{Mul}\;\church{3}\;(\textsc{Mul}\;\church{2}\;(\textsc{Fact}\;\church{1}))
    \rmark{unroll one more time} \\
  &\reduces \textsc{Mul}\;\church{3}\;(\textsc{Mul}\;\church{2}\;(\textsc{Mul}\;\church{1}\;\church{1}))
    \rmark{base case: $\church{0}$ gives $\church{1}$} \\
  &\reduces \church{6} \quad\checkmark
\end{align*}
Each recursive call ``unrolls'' the $\mathbf{Y}$ combinator one more time,
feeding $\textsc{Fact}$ back into $F$ as the $\mathit{self}$ argument. The
recursion terminates when $n$ reaches $\church{0}$ and \textsc{IsZero} selects
the base case ($\church{1}$).
\end{workedexample}

\begin{warning}[In typed calculi]
$\mathbf{Y}$ cannot be typed in the simply typed $\lam$-calculus. The problem
is the self-application $x\;x$: for this to type-check, $x$ would need to
simultaneously have type $T \to T$ (because it's being used as a function) and
type $T$ (because it's the argument to itself). No finite type can satisfy
both requirements. This is why typed languages like Lean, Haskell, and OCaml
provide recursion through dedicated language features (\texttt{fix},
\texttt{partial}, or structural recursion) rather than through combinators.
\end{warning}

%======================================================================
\part{Going Deeper}
%======================================================================

%======================================================================
\section{De Bruijn Indices}\label{sec:debruijn}
%======================================================================

Named variables require constant $\alpha$-conversion. In 1972, the Dutch
mathematician Nicolaas Govert de Bruijn proposed an elegant solution while
developing the \textbf{Automath} system---one of the earliest computerized
proof checkers. De Bruijn's problem was practical: when a computer manipulates
$\lam$-terms, it must constantly rename variables to avoid capture, which is
both error-prone and expensive. His insight was to eliminate names entirely: a
variable becomes a natural number counting binders outward.

\begin{defnbox}[De Bruijn Syntax]
\[
  e \;::=\; n \;\mid\; \lam.\, e \;\mid\; e_1\; e_2
  \qquad (n \in \Nat)
\]
Variable $n$ refers to the $n$-th enclosing $\lam$ (0-indexed).
\end{defnbox}

\begin{center}
\renewcommand{\arraystretch}{1.5}
\begin{tabular}{lcl}
\textbf{Named} & & \textbf{De Bruijn} \\
\hline
$\lam x.\, x$                            && $\lam.\, 0$ \\
$\lam x.\, \lam y.\, x$                  && $\lam.\, \lam.\, 1$ \\
$\lam x.\, \lam y.\, y$                  && $\lam.\, \lam.\, 0$ \\
$\lam x.\, \lam y.\, x\; y$              && $\lam.\, \lam.\, 1\; 0$ \\
$\lam f\, g\, x.\, f\;x\;(g\;x)$        && $\lam.\, \lam.\, \lam.\, 2\;0\;(1\;0)$ \\
\end{tabular}
\end{center}

Alpha equivalence becomes \textbf{structural equality}---no renaming needed.

\paragraph{The most common misconception.}
Your first instinct with the K combinator $\lam x.\, \lam y.\, x$ might be to
assign $x = 0$ and $y = 1$ (treating indices like serial numbers). This gives
$\lam.\, \lam.\, 0$---which is \textbf{wrong}. That's
$\lam x.\, \lam y.\, y$, the function that returns its second argument!

De Bruijn indices are \emph{not names}. They count distance. Each variable
looks \emph{upward} from where it sits and counts how many $\lam$'s it must
skip to reach its binder:

\begin{center}
\begin{tikzpicture}[scale=0.9, every node/.style={transform shape},
  node distance=0.3cm,
  binder/.style={font=\ttfamily\large},
  varnode/.style={font=\ttfamily\large},
  arrow/.style={-{Stealth[length=5pt]}, thick, rounded corners=3pt},
  label/.style={font=\small\itshape, text=black!60}
]
  % === LEFT: Wrong mental model ===
  \node[font=\large\bfseries] at (-4.5, 2.2) {Wrong: indices are names};
  
  \node[binder] (l1w) at (-5.5, 1.2) {$\lambda$};
  \node[label, right=0.05cm of l1w] (x0w) {``$x$ is the 0th parameter''};
  \node[binder, below=0.4cm of l1w] (l2w) {$\lambda$};
  \node[label, right=0.05cm of l2w] (y1w) {``$y$ is the 1st parameter''};
  \node[varnode, below=0.4cm of l2w] (vw) {0};
  \node[label, right=0.15cm of vw] {``$x$ was assigned 0''};
  
  \draw[arrow, red!70!black] (vw.west) -- +(-0.5, 0) |- (l2w.west);
  \node[red!70!black, font=\small\bfseries] at (-7.5, -0.25) {points to $y$!};
  
  \node[red!70!black, font=\Large\bfseries] at (-4.5, -1.2) {$\lambda.\,\lambda.\, 0 = \lambda x.\,\lambda y.\, y$ \;\; {\LARGE\texttimes}};

  % === RIGHT: Correct mental model ===
  \node[font=\large\bfseries] at (4.2, 2.2) {Right: indices count distance};
  
  \node[binder] (l1r) at (3, 1.2) {$\lambda$};
  \node[label, right=0.05cm of l1r] {binder for $x$};
  \node[binder, below=0.4cm of l1r] (l2r) {$\lambda$};
  \node[label, right=0.05cm of l2r] {binder for $y$ \;\; (skip this: $+1$)};
  \node[varnode, below=0.4cm of l2r] (vr) {1};
  \node[label, right=0.15cm of vr] {``my binder is 1 $\lambda$ up''};
  
  \draw[arrow, green!50!black] (vr.west) -- +(-0.5, 0) |- (l1r.west);
  \node[green!50!black, font=\small\bfseries] at (1.2, 0.4) {skip 1};
  
  \node[green!50!black, font=\Large\bfseries] at (4.2, -1.2) {$\lambda.\,\lambda.\, 1 = \lambda x.\,\lambda y.\, x$ \;\; \checkmark};
\end{tikzpicture}
\end{center}

Here is the rule: stand at the variable, look upward, count $\lam$'s
starting from 0. The \emph{nearest} $\lam$ is 0. The next one up is 1.
The variable's index is the number of the $\lam$ that \emph{binds} it.

\begin{center}
\begin{tikzpicture}[
  binder/.style={font=\ttfamily\large},
  varnode/.style={font=\ttfamily\large, fill=blue!8, rounded corners=2pt, inner sep=3pt},
  arrow/.style={-{Stealth[length=5pt]}, thick, rounded corners=3pt},
  countlabel/.style={font=\footnotesize, text=black!50}
]
  % S combinator: λf. λg. λx. f x (g x)
  \node[font=\large\bfseries] at (0, 3.3) {S combinator: $\lambda f.\, \lambda g.\, \lambda x.\, f\;x\;(g\;x)$};
  
  \node[binder] (l1) at (0, 2.4) {$\lambda$};
  \node[countlabel, left=0.3cm of l1] {binds $f$};
  
  \node[binder] (l2) at (0, 1.6) {$\lambda$};
  \node[countlabel, left=0.3cm of l2] {binds $g$};
  
  \node[binder] (l3) at (0, 0.8) {$\lambda$};
  \node[countlabel, left=0.3cm of l3] {binds $x$};
  
  % Body: f x (g x) = 2 0 (1 0)
  \node[varnode] (f) at (-2.5, -0.2) {2};
  \node[countlabel, below=0.05cm of f] {$f$};
  \node[varnode] (x1) at (-1.2, -0.2) {0};
  \node[countlabel, below=0.05cm of x1] {$x$};
  \node[varnode] (g) at (1.2, -0.2) {1};
  \node[countlabel, below=0.05cm of g] {$g$};
  \node[varnode] (x2) at (2.5, -0.2) {0};
  \node[countlabel, below=0.05cm of x2] {$x$};
  
  \node at (0, -0.2) {(};
  \node at (1.85, -0.2) {};
  \node at (3.1, -0.2) {)};
  
  % Arrows from variables to their binders
  \draw[arrow, blue!60!black] (f.north) -- ++(0, 0.15) -| ([xshift=-0.8cm]l1.south west) -- (l1.south west);
  \draw[arrow, red!60!black] (x1.north) -- (l3.south);
  \draw[arrow, orange!70!black] (g.north) -- ++(0, 0.15) -| ([xshift=0.5cm]l2.south east) -- (l2.south east);
  \draw[arrow, red!60!black] (x2.north) -- ++(0, 0.15) -| ([xshift=0.8cm]l3.south east) -- (l3.south east);
  
  % Count annotations
  \node[blue!60!black, font=\scriptsize] at (-3.5, 1.6) {skip 2};
  \node[red!60!black, font=\scriptsize] at (-1.2, 0.4) {skip 0};
  \node[orange!70!black, font=\scriptsize] at (2.2, 1.2) {skip 1};
  
  \node[font=\large] at (0, -1.2) {$= \lambda.\,\lambda.\,\lambda.\, 2\;\;0\;\;(1\;\;0)$};
\end{tikzpicture}
\end{center}

\begin{intuition}[The elevator rule]
Imagine each $\lam$ is a floor in a building, numbered from the bottom up.
A variable sits on the ground floor and needs to ride the elevator to its
binder's floor. The de Bruijn index is \emph{how many floors it rides}.
The nearest $\lam$ (floor 0) costs 0. The next one up costs 1. And so on.
The index tells you the distance, not the identity.
\end{intuition}

\subsection{Shifting and Substitution}

\begin{defnbox}[Shifting: $\uparrow^d_c(e)$]
Add $d$ to all variables $\geq c$:
\vspace{-0.5em}
\begin{align*}
  \uparrow^d_c(n) &=
    \begin{cases} n + d & \text{if } n \geq c \\ n & \text{if } n < c \end{cases} \\
  \uparrow^d_c(\lam.\, e) &= \lam.\, \uparrow^d_{c+1}(e) \\
  \uparrow^d_c(e_1\; e_2) &= \uparrow^d_c(e_1)\;\; \uparrow^d_c(e_2)
\end{align*}
\end{defnbox}

\begin{intuition}[Reading the shift notation]
$\uparrow^d_c(e)$ is a function of three things: the \textbf{shift amount} $d$
(an integer---can be negative), the \textbf{cutoff} $c$ (a natural number), and
the \textbf{term} $e$. The curly-brace $\begin{cases}\cdots\end{cases}$
notation is a piecewise definition (an if/else):
if the variable index $n$ is $\geq c$, add $d$; otherwise leave it alone.
The cutoff $c$ starts at 0 and increases by 1 each time we go under a $\lam$,
which is how we track ``how many binders are above us'' and only adjust
\emph{free} variables.

\medskip
\textbf{Why shifting?} When we substitute a term $s$ into a body that has a
$\lam$ removed, the indices of $s$'s free variables may be off by one. Shifting
corrects them. Think of it as updating line numbers after inserting or deleting
a line of code.
\end{intuition}

\begin{defnbox}[{De Bruijn Substitution}]
\vspace{-0.5em}
\begin{align*}
  n[k := s]                &= \begin{cases} s & \text{if } n = k \\ n & \text{otherwise} \end{cases} \\
  (\lam.\, e)[k := s]      &= \lam.\, \bigl(e[k{+}1 := \uparrow^1_0(s)]\bigr) \\
  (e_1\; e_2)[k := s]      &= (e_1[k := s])\;(e_2[k := s])
\end{align*}
\end{defnbox}

\begin{intuition}[Reading the De Bruijn substitution]
$e[k := s]$ means ``in term $e$, replace variable index $k$ with term $s$.''
\begin{enumerate}[leftmargin=2em]
\item If $e$ is a variable $n$: replace it with $s$ when $n = k$, otherwise leave it.
  (The $\begin{cases}\cdots\end{cases}$ is an if/else.)
\item If $e$ is $\lam.\,\mathit{body}$: going under a binder means bumping $k$
  to $k{+}1$ (the target variable is one level deeper now) \emph{and} shifting $s$
  up by 1 ($\uparrow^1_0(s)$) so that $s$'s free variables still point to the right
  binders from inside the extra $\lam$.
\item Applications: substitute into both sides independently.
\end{enumerate}
\end{intuition}

\begin{defnbox}[$\beta$-reduction with De Bruijn Indices]
\[
  \boxed{\;(\lam.\, e)\; s \;\;\reduce\;\; \uparrow^{-1}_0\bigl(e[0 := \uparrow^1_0(s)]\bigr)\;}
\]
Shift $s$ up by 1, substitute for index 0, shift result down by 1. \textbf{No $\alpha$-conversion needed.}
\end{defnbox}

\begin{intuition}[Unpacking the formula step by step]
Reading the formula right-to-left, inside-out:
\begin{enumerate}[leftmargin=2em]
\item $\uparrow^1_0(s)$: Shift the argument $s$ up by 1. This adjusts $s$'s
  free variables because $s$ is about to go \emph{under} the $\lam$ binder in
  the body~$e$.
\item $e[0 := \uparrow^1_0(s)]$: In the body $e$, replace every occurrence of
  variable index 0 (the parameter bound by the $\lam$ we're eliminating) with
  the shifted $s$.
\item $\uparrow^{-1}_0(\ldots)$: Shift the whole result \emph{down} by 1.
  The $\lam$ binder has been removed, so all remaining free variable indices
  need to decrease by 1 to compensate.
\end{enumerate}
\end{intuition}

\begin{workedexample}[De Bruijn: $(\lam.\,\lam.\, 1)\;(\lam.\, 0)$ \;{\small(i.e., $\mathbf{K}\;\mathbf{I}$)}]
Body $e = \lam.\, 1$, \; arg $s = \lam.\, 0$.
\begin{align*}
  \uparrow^1_0(\lam.\, 0) &= \lam.\, \uparrow^1_1(0) = \lam.\, 0
    \rmark{$0 < 1$, no change} \\
  (\lam.\, 1)[0 := \lam.\, 0]
    &= \lam.\, \bigl(1[1 := \uparrow^1_0(\lam.\, 0)]\bigr)
     = \lam.\, \bigl(1[1 := \lam.\, 0]\bigr)
     = \lam.\, (\lam.\, 0)
    \rmark{$1 = 1$, replace} \\
  \uparrow^{-1}_0(\lam.\, \lam.\, 0) &= \lam.\, \lam.\, 0
    \rmark{no free vars to shift}
\end{align*}
Result: $\lam.\,\lam.\, 0$, which in named form is $\lam a.\,\lam b.\, b$. But
wait---$\mathbf{K}\;\mathbf{I}$ should give $\lam y.\, \mathbf{I} = \lam y.\, \lam z.\, z = \lam.\,\lam.\, 0$. The inner 0 refers to the $z$ binder. So this \textbf{is correct}! $\checkmark$

\medskip
\textit{Moral: De Bruijn arithmetic is fiddly by hand---let the computer handle it.}
\end{workedexample}


%======================================================================


\subsection{Implementing It: De Bruijn in Lean}\label{sec:dbimpl}
%======================================================================

\subsubsection{When Names Fight Back}

Try reducing $(\lam x.\, \lam y.\, x)\; y$ in your head. Step 1: substitute
$y$ for $x$ in $\lam y.\, x$. But wait---$y$ is free in the argument, and
$\lam y$ would capture it! Rule 5 kicks in: rename the bound $y$ to $z$
first, giving $\lam z.\, y$.

This works, but it's \emph{annoying}. Every substitution requires checking
whether capture might happen, generating fresh names, and renaming. Worse, the
``same'' term can look different depending on which names were chosen during
renaming. $\lam x.\, x$ and $\lam y.\, y$ and $\lam z.\, z$ are all the
identity function, but they're different strings of symbols. If you want to
check whether two terms are equal, you need $\alpha$-equivalence, which itself
requires comparing structure modulo renaming.

\begin{exercisebox}[Design challenge]
\stepcounter{exercise}\textbf{\theexercise.}\; Imagine you're designing a system where $\alpha$-equivalent terms are \emph{literally identical}---not just ``the same up to renaming,'' but the same data structure with the same bits. What would you use instead of names to refer to variables?

\emph{Hint:} What stays the same about $\lam x.\, x$ and $\lam y.\, y$? It's not the name. What is it?
\end{exercisebox}

\subsubsection{The Insight: Count, Don't Name}

Here's what $\lam x.\, x$ and $\lam y.\, y$ have in common: the variable in
the body refers to the \emph{nearest enclosing binder}. Not ``the binder named
$x$'' or ``the binder named $y$''---just ``the binder that's one $\lam$ up.''

If we replace names with numbers that count binders outward, both become the
same thing: $\lam.\, 0$ (``the variable bound by the nearest $\lam$, zero
layers up''). This is exactly the de Bruijn representation from
Section~\ref{sec:debruijn}.

\begin{minted}{lean}
inductive Expr where
  | bvar : Nat → Expr        -- bound variable: de Bruijn index
  | lam  : Expr → Expr       -- lambda: no name needed!
  | app  : Expr → Expr → Expr
  deriving Repr, DecidableEq, Inhabited
\end{minted}

Notice: \texttt{lam} has \emph{no} \texttt{String} parameter. The binder
doesn't need a name---variables find their binder by counting.

\subsubsection{The New Challenge: Shifting}

With names, substitution needed renaming to avoid capture. With de Bruijn
indices, renaming is gone---but a new challenge appears. When you move a term
under or out from a binder, the indices must be adjusted.

\begin{exercisebox}[Shifting challenge]
\stepcounter{exercise}\textbf{\theexercise.}\; Consider the term $\lam.\, 0\; 1$ (which in named form is $\lam x.\, x\;y$ for some free $y$). If we substitute this into a context that has one more binder above it, the free variable $1$ (referring to $y$) needs to increase by 1 to still point to the same thing. But $0$ (the bound variable) must \emph{not} change. How do you decide which indices to adjust?
\end{exercisebox}

Think about this carefully. The index $0$ points to the $\lam$ right above it---
it's a \textbf{bound} variable. The index $1$ points past the $\lam$ to
something external---it's \textbf{free}. When we move the term under another
binder, all the free indices shift up by 1 (they now have one more binder to
``skip over''), but the bound indices stay the same (they still point to the
same $\lam$ they always did).

So the question becomes: how do you distinguish free indices from bound ones?

The answer: track how many binders you've descended through. Call this the
\textbf{cutoff} $c$. Start at $c = 0$. Each time you enter a $\lam$, increment
$c$. An index $n$ is free (and needs adjusting) if $n \geq c$, and bound (and
stays fixed) if $n < c$.

Why does this work? If you're inside $c$ nested binders, then indices
$0, 1, \ldots, c{-}1$ all point to one of those binders (they're bound). Any
index $\geq c$ points past all of them (it's free).

\begin{center}
\begin{tikzpicture}[
  binder/.style={font=\ttfamily\large},
  varnode/.style={font=\ttfamily\large},
  arrow/.style={-{Stealth[length=5pt]}, thick, rounded corners=3pt},
  shifted/.style={font=\ttfamily\large, red!70!black},
  same/.style={font=\ttfamily\large, green!50!black},
  label/.style={font=\small\itshape, text=black!60}
]
  % === LEFT: Before shifting ===
  \node[font=\large\bfseries] at (-3.5, 3.5) {Before: $\lambda.\;0\;\;1$};
  
  \node[binder] (l1) at (-3.5, 2.5) {$\lambda$};
  \node[label, right=0.1cm of l1] {(the term's own $\lambda$)};
  
  \node[varnode] (v0) at (-4.5, 1.5) {0};
  \node[label, below=0cm of v0] {bound};
  \node[varnode] (v1) at (-2.5, 1.5) {1};
  \node[label, below=0cm of v1] {free};
  
  \draw[arrow, green!50!black] (v0.north) -- (l1.south west);
  \draw[arrow, blue!60!black] (v1.north) -- ++(0, 0.3) -- ++(1.5, 0) -- ++(0, 0.9)
    node[right, font=\small, text=blue!60!black] {to external binder};
    
  % === ARROW ===
  \node[font=\Huge] at (0, 2) {$\Rightarrow$};
  \node[font=\small, text=black!60] at (0, 1.3) {moves under};
  \node[font=\small, text=black!60] at (0, 0.9) {a new $\lambda$};
  
  % === RIGHT: After shifting ===
  \node[font=\large\bfseries] at (4.2, 3.5) {After: $\lambda.\;0\;\;2$};
  
  \node[binder] (nl) at (4.2, 2.9) {$\lambda$};
  \node[label, right=0.1cm of nl] {(new outer $\lambda$)};
  \node[binder] (l1r) at (4.2, 2.1) {$\lambda$};
  \node[label, right=0.1cm of l1r] {(the term's own $\lambda$)};
  
  \node[same] (v0r) at (3.2, 1.1) {0};
  \node[label, below=0cm of v0r] {unchanged};
  \node[shifted] (v2r) at (5.2, 1.1) {2};
  \node[label, below=0cm of v2r] {was 1, now 2};
  
  \draw[arrow, green!50!black] (v0r.north) -- (l1r.south west);
  \draw[arrow, blue!60!black] (v2r.north) -- ++(0, 0.3) -- ++(1, 0) -- ++(0, 1.7)
    node[right, font=\small, text=blue!60!black] {same external binder};
\end{tikzpicture}
\end{center}

The bound variable 0 still points to its own $\lam$---nothing changed.
But the free variable had to increase from 1 to 2: it now needs to skip
\emph{two} $\lam$'s to reach the same external binder it always referred to.

\begin{minted}{lean}
def shift (d : Int) (c : Nat) : Expr → Expr
  | .bvar n    => if n ≥ c then .bvar (Int.toNat (n + d)) else .bvar n
  | .lam e     => .lam (shift d (c + 1) e)     -- going under a binder: c increases
  | .app e₁ e₂ => .app (shift d c e₁) (shift d c e₂)
\end{minted}

With shifting in hand, we're ready for substitution. But this is the hardest
part of the de Bruijn implementation---it took real researchers (de Bruijn
himself, then many others refining the details) years to get this right. The
difficulty is that \emph{two} adjustments are needed, and it's hard to see why
until you try the wrong thing and watch it break.

But before we write any code, there is a conceptual trap that catches almost
everyone: confusing \textbf{substitution} with \textbf{beta reduction}. Getting
this distinction clear first will save you from a lot of confusion later.

\paragraph{Substitution is not beta reduction.}
It's natural to think of substitution as ``plugging an argument into a
function.'' But that's beta reduction, which is a \emph{computation step}.
Substitution is a lower-level operation that beta reduction \emph{uses}.

\begin{center}
\renewcommand{\arraystretch}{1.4}
\begin{tabular}{lp{9cm}}
\textbf{Beta reduction} & A computation step: strip the outer $\lam$, then
  substitute the argument into the body, adjusting indices for the removed
  binder. This is what happens when you ``call a function.'' \\
\textbf{Substitution} & A textual operation: find every free occurrence of a
  variable and replace it with a term. It knows nothing about function calls---it
  just does search-and-replace. \\
\end{tabular}
\end{center}

Beta reduction \emph{calls} substitution as one of its internal steps, but they
are different things. Consider the K combinator, $\lam x.\, \lam y.\, x$. If
you try to substitute directly into the whole term:

\begin{minted}{lean}
-- WRONG: trying to use subst as if it were beta reduction
-- This does NOT compute "K a"!
#eval subst "x" (.var "a") K    -- named version
#eval subst 0 (.bvar 99) K_db  -- de Bruijn version
\end{minted}

Both produce the \emph{original term, unchanged}. Why? Because $x$ (or index
$0$) is not free in K---it's bound by K's own outermost $\lam$. Substitution
only replaces \emph{free} occurrences, and there are none.

To actually compute ``K applied to \texttt{a},'' you need beta reduction, which
first \textbf{strips the outer $\lam$} and then substitutes into the
\textbf{body}:

\begin{minted}{text}
Beta reduction of K a:
   (λx. λy. x) a             -- the full application
   = (λy. x)[x := a]         -- strip outer λ, substitute into BODY
   = λy. a                    -- x was free in the body, so it got replaced

De Bruijn version:
   (λ. λ. 1) a
   = (λ. 1)[0 := a]          -- strip outer λ, substitute into BODY
                               -- 0 means "the variable the stripped λ bound"
                               -- from the body's top level, that's index 0
   = ...                       -- (we'll work out the details below)
\end{minted}

The key insight: the \texttt{subst 0} in beta reduction targets index $0$
\emph{relative to the body}, not relative to the original term. The outer
$\lam$ is already gone. From the body's perspective, the variable it used to
bind is the nearest external thing---index $0$.

\begin{intuition}[Substitution is a tool{,} not a computation]
Think of substitution as a screwdriver. Beta reduction is the act of assembling
furniture. You don't hand someone a screwdriver and say ``build me a bookshelf''
---the screwdriver is one tool used during assembly. Similarly, you don't call
\texttt{subst} on a whole $\lam$-term and expect it to ``run'' the function.
Beta reduction orchestrates the process: it strips the $\lam$, prepares the
argument, calls \texttt{subst} to do the replacement, and then cleans up the
indices.
\end{intuition}

With that distinction clear, let's build \texttt{subst} itself. We'll do it in
three stages: a na\"ive version that's almost right, then two fixes that
repair it.

\paragraph{What substitution means, concretely.}
$\texttt{subst}\;k\;s\;e$ means ``in $e$, replace every free occurrence
of index $k$ with the term $s$.'' The caller decides what $k$ and $s$ are.
When beta reduction calls substitution, it always starts with $k = 0$ (the
variable the stripped $\lam$ used to bind). But substitution itself is
general---it can target any index.

\paragraph{Stage 1: the na\"ive attempt.}

Let's try the simplest thing that could possibly work. For each term form:

\begin{itemize}[leftmargin=2em]
\item \textbf{Variable $n$}: if $n = k$, replace it with $s$. Otherwise leave it.
\item \textbf{Application}: recurse into both sides.
\item \textbf{Lambda}: recurse into the body.
\end{itemize}

\begin{minted}{lean}
-- WRONG! But a useful starting point.
def substNaive (k : Nat) (s : Expr) : Expr → Expr
  | .bvar n    => if n == k then s else .bvar n
  | .app e₁ e₂ => .app (substNaive k s e₁) (substNaive k s e₂)
  | .lam e     => .lam (substNaive k s e)  -- BUG: k and s unchanged
\end{minted}

This works for terms with no nested lambdas. Let's test it on a simple case:

\begin{minted}{text}
substNaive 0 (bvar 5) (bvar 0)       = bvar 5     ✓  replaced
substNaive 0 (bvar 5) (bvar 3)       = bvar 3     ✓  left alone
substNaive 0 (bvar 5) (app (bvar 0) (bvar 1))
                                      = app (bvar 5) (bvar 1)  ✓
\end{minted}

But now try it on a lambda:

\begin{minted}{text}
substNaive 0 (bvar 5) (λ. 0)         = λ. (bvar 5)     ✗ WRONG!
\end{minted}

The term $\lam.\, 0$ is $\lam x.\, x$---the identity function. Index $0$
inside the body refers to the $\lam$'s \emph{own} parameter, not to the
external variable $0$ we wanted to replace. But our na\"ive code doesn't know
the difference: it sees index $0$, matches $k = 0$, and replaces it.

\begin{exercisebox}[Diagnosing the bug]
\stepcounter{exercise}\textbf{\theexercise.}\; The na\"ive version also fails on $\texttt{substNaive}\;0\;(\texttt{bvar}\;5)\;(\lam.\, \lam.\, 0\;\;1\;\;2)$. The term $\lam.\,\lam.\, 0\;\;1\;\;2$ is $\lam x.\,\lam y.\, y\;x\;z$ where $z$ is free (index $2$). Which variables get incorrectly replaced?

\emph{Hint:} Inside two $\lam$'s, indices $0$ and $1$ are bound. The free variable $0$ from outside is now at index $2$.
\end{exercisebox}

\paragraph{Stage 2: fix the target index.}

The problem is clear: when we descend under a $\lam$, every index gets pushed
up by one (there's a new binder between us and whatever the index pointed to).
The external variable $0$ that we're looking for doesn't disappear---it becomes
index $1$ inside the first $\lam$, index $2$ inside two $\lam$'s, and so on.

So when we recurse under a $\lam$, we must \emph{increment the target}:

\begin{minted}{lean}
-- STILL WRONG! But closer.
def substV2 (k : Nat) (s : Expr) : Expr → Expr
  | .bvar n    => if n == k then s else .bvar n
  | .app e₁ e₂ => .app (substV2 k s e₁) (substV2 k s e₂)
  | .lam e     => .lam (substV2 (k + 1) s e)  -- FIX 1: target shifts
                  --              ^^^^^
\end{minted}

Let's re-test:

\begin{minted}{text}
substV2 0 (bvar 5) (λ. 0)
  = λ. (substV2 1 (bvar 5) (bvar 0))
  = λ. (bvar 0)                            ✓  bound var untouched!

substV2 0 (bvar 5) (λ. 1)
  = λ. (substV2 1 (bvar 5) (bvar 1))
  = λ. (bvar 5)                            ✓  free var replaced!
\end{minted}

Progress! But there's still a bug. Try this:

\begin{minted}{text}
substV2 0 (bvar 1) (λ. 1)
  = λ. (substV2 1 (bvar 1) (bvar 1))
  = λ. (bvar 1)                            ✗ WRONG!
\end{minted}

We're substituting $s = \texttt{bvar}\;1$ (some free variable) into
$\lam.\, 1$. Index $1$ in the body refers to the external variable $0$---the
one we want to replace. So the result should contain $s$. And it does contain
\texttt{bvar 1}\ldots\ but now that \texttt{bvar 1} is \emph{inside a
$\lam$}, so it points to a different thing! From inside the $\lam$, index $1$
refers to whatever is \emph{one} binder outside, not to the original free
variable $s$ was pointing at.

\begin{exercisebox}[Seeing the capture]
\stepcounter{exercise}\textbf{\theexercise.}\; Work through $\texttt{substV2}\;0\;(\texttt{bvar}\;0)\;(\lam.\, 1)$ in named notation.

The named version is: substitute $y$ for $x$ in $\lam z.\, x$. The result should be $\lam z.\, y$. But index $0$ inside the $\lam$ refers to $z$, not $y$!

What index should the result have inside the $\lam$ to correctly refer to $y$?

\emph{Answer:} It should be $\texttt{bvar}\;1$ inside the $\lam$ (skip past $z$'s binder to reach $y$). The replacement term $\texttt{bvar}\;0$ needed to become $\texttt{bvar}\;1$ because it moved under a binder.
\end{exercisebox}

\paragraph{Stage 3: fix the replacement term.}

The replacement term $s$ is being dropped into a context that has one more
binder than where $s$ came from. Its free variables need to increase by $1$
to skip over that new binder. This is exactly what \texttt{shift} does:

\begin{minted}{lean}
-- CORRECT!
def subst (k : Nat) (s : Expr) : Expr → Expr
  | .bvar n    => if n == k then s else .bvar n
  | .lam e     => .lam (subst (k + 1) (shift 1 0 s) e)
  --                          ^^^^^^^  ^^^^^^^^^^^^
  --                          FIX 1    FIX 2
  --                          target   replacement moves under
  --                          moves    a binder: shift its
  --                          further  free variables up by 1
  --                          away
  | .app e₁ e₂ => .app (subst k s e₁) (subst k s e₂)
\end{minted}

Here is what happens visually when substitution walks under a $\lam$:

\begin{center}
\begin{tikzpicture}[
  binder/.style={font=\ttfamily\large},
  varnode/.style={font=\ttfamily, inner sep=2pt},
  box/.style={draw, rounded corners=4pt, inner sep=6pt, thick},
  arrow/.style={-{Stealth[length=5pt]}, thick},
  fixlabel/.style={font=\small\bfseries, fill=yellow!20, rounded corners=2pt, inner sep=2pt}
]
  % === OUTSIDE THE λ ===
  \node[font=\bfseries] at (-4.8, 3.2) {Outside the $\lambda$:};
  \node[box, blue!60!black] (outer) at (-4.8, 2.3) {
    \begin{tabular}{l}
    hunting for: $k = 0$ \\
    will insert: $s$
    \end{tabular}
  };
  
  % === THE LAMBDA ===
  \node[binder, font=\Huge] (lam) at (0, 2.3) {$\lambda$};
  \draw[arrow, thick, black!40] (-2.8, 2.3) -- (-0.5, 2.3)
    node[midway, above, font=\small\itshape, text=black!50] {entering};
  
  % === INSIDE THE λ ===
  \node[font=\bfseries] at (4.8, 3.2) {Inside the $\lambda$:};
  \node[box, red!60!black] (inner) at (4.8, 2.3) {
    \begin{tabular}{l}
    hunting for: $k = 1$ \\
    will insert: $\mathit{shift}\;1\;0\;s$
    \end{tabular}
  };
  
  \draw[arrow, thick, black!40] (0.5, 2.3) -- (2.6, 2.3);
  
  % === FIX LABELS ===
  \node[fixlabel, blue!70!black] at (1.5, 0.9) {Fix 1: $k \to k + 1$};
  \node[font=\small, text=black!70] at (1.5, 0.3) {The variable we want is now};
  \node[font=\small, text=black!70] at (1.5, -0.1) {one binder further away};
  
  \node[fixlabel, red!70!black] at (1.5, -0.9) {Fix 2: $s \to \mathit{shift}\;1\;0\;s$};
  \node[font=\small, text=black!70] at (1.5, -1.5) {The replacement's free vars must};
  \node[font=\small, text=black!70] at (1.5, -1.9) {skip past this new $\lambda$};
\end{tikzpicture}
\end{center}

Let's verify the case that broke before:

\begin{minted}{text}
subst 0 (bvar 1) (λ. 1)
  = λ. (subst 1 (shift 1 0 (bvar 1)) (bvar 1))
  = λ. (subst 1 (bvar 2) (bvar 1))       -- s shifted: 1 → 2
  = λ. (bvar 2)                           ✓  correctly points past the λ
\end{minted}

\paragraph{But wait---why does 1 become 2?}

This is the most confusing part of de Bruijn substitution, and it's worth
pausing on. You might think: ``The target term $\lam.\, 1$ already uses
indices from \emph{inside} the $\lam$. The 1 already accounts for the $\lam$
being there. So why are we adjusting anything?''

You're right that the target term is already in ``inside coordinates.'' The
problem is that $k$ and $s$ are not. They were decided from \emph{outside} the
$\lam$, in ``outside coordinates.'' When we step inside, we need to translate
them.

\paragraph{The coordinate system analogy.}
Imagine naming each variable by its index from the outside. Suppose
the context above us has $x$ at index 0 and $z$ at index 1:

\begin{minted}{text}
Standing OUTSIDE the lambda:
    index 0 = x       ← k=0 means "replace x"
    index 1 = z       ← s=(bvar 1) means "insert z"

Standing INSIDE the lambda (which binds w):
    index 0 = w       ← the lambda's own parameter (NEW)
    index 1 = x       ← was index 0 from outside
    index 2 = z       ← was index 1 from outside
\end{minted}

The lambda inserted a new floor between us and everything external.
From outside, the instruction was ``replace $x$ (index 0) with $z$ (index 1).''
From inside, the \emph{same} instruction is ``replace $x$ (index 1) with $z$
(index 2).'' In named notation, nothing changed---it's still ``replace $x$ with
$z$.'' But in de Bruijn indices, the \emph{addresses} of $x$ and $z$ both
shifted up by 1 because we moved deeper into the building.

So the substitution transforms from:
\[
  (\lam.\, 1)[0 := \texttt{bvar}\;1]
  \quad\xrightarrow{\text{enter the } \lam}\quad
  (1)[1 := \texttt{bvar}\;2]
\]

And now $1 = 1$, so the substitution fires, producing $\texttt{bvar}\;2$.
Wrapped back in the $\lam$: $\lam.\, 2$.

In named terms, this is just $(\lam w.\, x)[x := z] = \lam w.\, z$. The
result still ignores its parameter and returns a free variable---but $z$
instead of $x$.

\begin{intuition}[The two fixes are a coordinate translation]
The body of the $\lam$ is already expressed in ``inside coordinates.'' But
$k$ and $s$ arrived from outside, in ``outside coordinates.'' The two fixes
translate them into the same coordinate system as the body:

\textbf{Fix 1 ($k + 1$):} Translates the target. ``Replace index 0'' becomes
``replace index 1'' because $x$ moved from position 0 to position 1 when
$w$ took position 0.

\textbf{Fix 2 (\texttt{shift 1 0 s}):} Translates the replacement. ``Insert
bvar 1'' becomes ``insert bvar 2'' because $z$ moved from position 1 to
position 2 for the same reason.

Both fixes do the same thing: add 1 to account for the new binder. Fix 1
does it to a single number ($k$). Fix 2 does it to every free variable
inside a term ($s$). Same reason, same adjustment, applied to different things.
\end{intuition}

\paragraph{Why shifting is separate from substitution.}

You might wonder: why not bake shifting directly into the substitution function
instead of having a separate \texttt{shift}? The reason is that
\texttt{shift} is needed independently in beta reduction (step 1 and step 3
below), and having it as a separate, well-tested function makes each piece
easier to reason about. The separation mirrors a general principle: build
small, correct components and compose them.

\begin{exercisebox}[Building intuition]
\stepcounter{exercise}\textbf{\theexercise.}\; Before reading further, think about what substitution should do for each of the three term forms:

(a) If the term is a \textbf{variable} $n$: when should it be replaced? What happens when $n \neq k$?

(b) If the term is an \textbf{application} $e_1\;e_2$: what do you do?

(c) If the term is a \textbf{lambda} $\lam.\,\mathit{body}$: which of the two fixes applies here, and why?

(d) What would go wrong if you applied fix 1 but not fix 2? What about fix 2 but not fix 1?
\end{exercisebox}

\begin{exercisebox}[Tracing the target]
\stepcounter{exercise}\textbf{\theexercise.}\; Suppose we're doing $\texttt{subst}\;0\;s\;(\lam.\, \lam.\, \lam.\, 3)$.

(a) At the outermost call, what is $k$?

(b) After descending through the first $\lam$, what is $k$?

(c) After all three $\lam$s, what is $k$? Does it match the index $3$?

\emph{Answer:} $k$ goes $0 \to 1 \to 2 \to 3$. Yes, index $3$ inside three binders refers to the same variable as index $0$ outside. The substitution fires correctly.
\end{exercisebox}

\begin{exercisebox}[Tracing substitution]
\stepcounter{exercise}\textbf{\theexercise.}\; Work through $\texttt{subst}\;0\;(\texttt{bvar}\;1)\;(\lam.\, 0\;\;1)$ by hand. This represents substituting a free variable (index 1) for variable 0 in the body $\lam x.\, x\;\;\square$ (where $\square$ is the variable being replaced).

Step 1: We hit the $\lam$. We recurse with $k = 0 + 1 = 1$ and $s = \texttt{shift}\;1\;0\;(\texttt{bvar}\;1) = \texttt{bvar}\;2$.

Step 2: Now we substitute into the application $0\;\;1$:
\begin{itemize}[leftmargin=2em]
\item Left: $\texttt{subst}\;1\;(\texttt{bvar}\;2)\;(\texttt{bvar}\;0)$. Is $0 = 1$? No. Result: $\texttt{bvar}\;0$. (This is the $\lam$'s own variable---correctly untouched.)
\item Right: $\texttt{subst}\;1\;(\texttt{bvar}\;2)\;(\texttt{bvar}\;1)$. Is $1 = 1$? Yes! Result: $\texttt{bvar}\;2$.
\end{itemize}

Final result: $\lam.\, 0\;\;2$. The bound variable $0$ stayed put; the free
variable got replaced and correctly shifted.
\end{exercisebox}

\paragraph{Beta reduction: the full dance.}
This is where the distinction from earlier pays off. Substitution is the
search-and-replace tool; beta reduction is the process that \emph{uses} it.
When we reduce $(\lam.\,\mathit{body})\;\mathit{arg}$, three things happen:

\begin{enumerate}[leftmargin=2em]
\item \textbf{Shift the argument up:} $\mathit{arg}$ is about to enter the body, which is inside the $\lam$. Its free variables need $+1$: $\texttt{shift}\;1\;0\;\mathit{arg}$.
\item \textbf{Substitute:} Replace index $0$ in the body with the shifted argument: $\texttt{subst}\;0\;(\texttt{shift}\;1\;0\;\mathit{arg})\;\mathit{body}$. Here is where substitution gets called---on the \textbf{body}, not the whole $\lam$-term. The outer $\lam$ is already stripped. Index $0$ means ``the variable that $\lam$ used to bind,'' which is the nearest external variable from the body's perspective.
\item \textbf{Shift the result down:} The $\lam$ is gone---we've consumed it. All remaining free variables that pointed past it need $-1$: $\texttt{shift}\;(-1)\;0\;(\ldots)$.
\end{enumerate}

\begin{center}
\begin{tikzpicture}[
  stpbox/.style={draw, rounded corners=4pt, inner sep=5pt, thick, font=\small,
               minimum width=2.2cm, text width=2cm, align=center},
  arrow/.style={-{Stealth[length=5pt]}, thick},
  term/.style={font=\ttfamily\small},
  annot/.style={font=\scriptsize\itshape, text=black!60}
]
  % Example: (λ. λ. 1) (bvar 0)   -- K combinator applied to free var
  \node[font=\small\bfseries] at (0, 2.5) {Example: $(\lambda.\,\lambda.\, 1)\;\;(\texttt{bvar}\;0)$};
  \node[annot] at (0, 2.0) {K combinator applied to a free variable};
  
  \node[term] at (-4.5, 1.0) {arg $= 0$};
  
  \node[stpbox, fill=blue!10] (s1) at (-2.5, 0) {\textbf{Step 1}\\shift arg\\up by 1};
  \node[stpbox, fill=orange!10] (s2) at (0.5, 0) {\textbf{Step 2}\\substitute\\into body};
  \node[stpbox, fill=green!10] (s3) at (3.5, 0) {\textbf{Step 3}\\shift result\\down by 1};
  
  \draw[arrow] (-4.0, 0.5) -- (s1.north west);
  \draw[arrow] (s1) -- (s2) node[midway, above, term] {arg $= 1$};
  \draw[arrow] (s2) -- (s3) node[midway, above, term] {$\lambda.\, 2$};
  \draw[arrow] (s3.east) -- ++(1.5, 0) node[right, term] {$\lambda.\, 1$};
  
  \node[annot] at (-2.5, -0.9) {arg enters $\lambda$:};
  \node[annot] at (-2.5, -1.2) {free vars $+1$};
  \node[annot] at (0.5, -0.9) {replace idx 0};
  \node[annot] at (0.5, -1.2) {in body with arg};
  \node[annot] at (3.5, -0.9) {outer $\lambda$ gone:};
  \node[annot] at (3.5, -1.2) {free vars $-1$};
  
  \node[annot] at (5.8, -0.2) {$= \lambda y.\, z$};
  \node[annot] at (5.8, -0.5) {(ignores $y$, returns};
  \node[annot] at (5.8, -0.8) {the free variable)};
\end{tikzpicture}
\end{center}

\begin{minted}{lean}
def betaReduce (body arg : Expr) : Expr :=
  shift (-1) 0 (subst 0 (shift 1 0 arg) body)
\end{minted}

\paragraph{Why all three steps are needed.}
Each step exists for a specific reason, and skipping any one produces wrong
results. The shift-up and shift-down are \emph{not} redundant---they handle
different lambdas. Substitution's internal shifting (which happens when
\texttt{subst} walks under lambdas \emph{inside} the body) handles lambdas
that are still present. The shift-up and shift-down in \texttt{betaReduce}
handle the lambda that is being \emph{stripped}.

Consider a body that has \textbf{no inner lambdas}---just a bare variable:

\begin{minted}{text}
(λ. 0) (bvar 1)            -- identity applied to a free variable
body = bvar 0, arg = bvar 1

Without shift-up:  subst 0 (bvar 1) (bvar 0) = bvar 1
                   shift (-1) 0 (bvar 1)      = bvar 0    ✗ wrong variable!

With shift-up:     shift 1 0 (bvar 1)         = bvar 2
                   subst 0 (bvar 2) (bvar 0)  = bvar 2
                   shift (-1) 0 (bvar 2)      = bvar 1    ✓ correct!
\end{minted}

Without the shift-up, substitution drops the raw argument into the body. Then
the shift-down corrupts it because it was never adjusted for the lambda's
scope. The shift-up ``pre-adjusts'' the argument so that the shift-down brings
it back to the right value.

Now consider a case where a variable \textbf{survives} substitution. Start
with $(\lam.\, 1)\;(\texttt{bvar}\;5)$, i.e., $(\lam x.\, y)\;\mathit{arg}$
where $y$ is free. The function ignores its argument and just returns $y$.

In the original term $\lam.\, 1$, the $\texttt{bvar}\;1$ is \textbf{bound}---
it points to a binder above the $\lam$ we're about to strip. After stripping,
the body is just $\texttt{bvar}\;1$, and the binder it pointed to is gone.
It's an orphan:

\begin{minted}{text}
(λ. 1) (bvar 5)            -- named: (λx. y) arg; y is bound above

body = bvar 1               -- after stripping, y is FREE (binder gone!)
arg  = bvar 5
target index 0              -- x (the stripped λ's variable)

Step 1: shift arg up      → bvar 6
Step 2: subst 0 (bvar 6) (bvar 1)
        1 ≠ 0, leave it   → bvar 1     -- y survived, x wasn't even here

Without shift-down: bvar 1              ✗ points past a λ that's gone!
With shift-down:    bvar 0              ✓ y is now the nearest external
\end{minted}

The $\texttt{bvar}\;1$ was counting the stripped lambda as one of the floors
it needed to skip. That floor is gone now, so it must drop to
$\texttt{bvar}\;0$.

A richer example makes both roles visible. Take
$(\lam.\, \lam.\, 1\;\;0)\;(\texttt{bvar}\;5)$, i.e.,
$(\lam y.\, \lam x.\, y\;x)\;\mathit{arg}$:

\begin{minted}{text}
body = λ. 1 0,  arg = bvar 5

Step 1: shift arg up       → bvar 6
Step 2: subst 0 (bvar 6) (λ. 1 0)
  enter inner λ: subst 1 (bvar 7) (app (bvar 1) (bvar 0))
    bvar 1 (y): 1 == 1 → replaced with bvar 7       ← target
    bvar 0 (x): 0 ≠ 1 → stays                       ← inner λ's own param
  result: λ. 7 0
Step 3: shift down         → λ. 6 0

Named result: λx. arg x   -- takes x, applies arg to x
\end{minted}

Here $y$ gets replaced, so the shift-down only adjusts the argument's index
($7 \to 6$). The inner $\lam$'s own variable $x$ (index 0) is bound and
unaffected.

\paragraph{The building analogy.}
The clearest way to see shift-down is to think of lambdas as floors in a
building. Each variable stores its binder's floor number, counted from its own
position upward. Beta reduction \emph{demolishes} a floor---and every variable
that pointed above it must update its count.

Consider $(\lam.\, 1)\;(\texttt{bvar}\;5)$, i.e., $(\lam x.\, y)\;\mathit{arg}$:

\begin{center}
\begin{tikzpicture}[
  floor/.style={draw, minimum width=4.5cm, minimum height=0.7cm, thick},
  stripped/.style={draw, minimum width=4.5cm, minimum height=0.7cm, thick,
                   fill=red!15, dashed},
  varbox/.style={font=\ttfamily\small, fill=blue!8, rounded corners=2pt,
                 inner sep=3pt},
  arrow/.style={-{Stealth[length=5pt]}, thick, rounded corners=3pt},
  flabel/.style={font=\small\itshape, text=black!50},
  phase/.style={font=\small\bfseries}
]
  % === BEFORE: (λx. y) arg ===
  \node[phase] at (-5, 4.5) {Before: $(\lambda.\, 1)\;\;\mathit{arg}$};
  \node[font=\small, text=black!60] at (-5, 4) {$y$ is \textbf{bound} above};
  
  % y's binder (external)
  \node[floor, fill=gray!10] (ext) at (-5, 3.1) {};
  \node[flabel] at (-5, 3.1) {floor 1: binds $y$};
  
  % x's lambda (will be stripped)
  \node[stripped] (lx) at (-5, 2.25) {};
  \node[flabel] at (-5, 2.25) {floor 0: $\lambda$ binds $x$};
  \node[font=\scriptsize\bfseries, red!60!black] at (-2.3, 2.25) {$\leftarrow$ stripped!};
  
  % Variable y
  \node[varbox] (vy) at (-5, 1.25) {1};
  \node[font=\scriptsize, below=0.02cm of vy] {$y$ (bound)};
  
  \draw[arrow, orange!70!black] (vy.north) -- ++(0, 0.15)
    -| ([xshift=-0.7cm]ext.south west) -- (ext.south west);
  \node[orange!70!black, font=\scriptsize] at (-6.7, 2.5) {skip 1};
  
  % === MIDDLE: body after stripping ===
  \draw[-{Stealth[length=8pt]}, ultra thick, black!30] (-2.2, 2.5) -- (-0.3, 2.5)
    node[midway, above, font=\scriptsize\itshape, text=black!40] {strip $\lambda$};
  
  \node[phase] at (1.5, 4.5) {Body: $\texttt{bvar}\;1$};
  \node[font=\small, text=red!60!black] at (1.5, 4) {binder is \textbf{gone}!};
  
  \node[floor, fill=gray!10] (ext1b) at (1.5, 3.1) {};
  \node[flabel] at (1.5, 3.1) {floor 0: binds $y$ \;\;(was floor 1)};
  
  % demolished floor
  \node[font=\small, text=red!40, align=center] at (1.5, 2.25) {(demolished)};
  
  \node[varbox] (vy1b) at (1.5, 1.25) {1};
  \node[font=\scriptsize, text=red!60!black, below=0.02cm of vy1b] {orphan!};
  
  \draw[arrow, red!60!black, dashed] (vy1b.north) -- ++(0, 0.65)
    node[right, font=\scriptsize, text=red!50] {points past $y$};
  
  % === AFTER: shift down fixes it ===
  \draw[-{Stealth[length=8pt]}, ultra thick, black!30] (4, 2.5) -- (5.8, 2.5)
    node[midway, above, font=\scriptsize\itshape, text=black!40] {shift $-1$};
  
  \node[phase] at (8, 4.5) {Result: $\texttt{bvar}\;0$};
  \node[font=\small, text=green!50!black] at (8, 4) {index \textbf{corrected}};
  
  % y's binder (now floor 0)
  \node[floor, fill=gray!10] (ext2) at (8, 3.1) {};
  \node[flabel] at (8, 3.1) {floor 0: binds $y$};
  
  % Variable y (fixed)
  \node[varbox] (vy2) at (8, 2.1) {0};
  \node[font=\scriptsize, below=0.02cm of vy2] {$y$ (fixed)};
  
  \draw[arrow, green!50!black] (vy2.north) -- (ext2.south);
  \node[green!50!black, font=\scriptsize] at (9.0, 2.7) {skip 0};
\end{tikzpicture}
\end{center}

The demolished floor is gone, so $y$'s floor (which was floor 1) is now
floor 0. The variable was still saying ``floor 1'' after the demolition---the
shift-down corrects it to ``floor 0.'' In general, every free variable that
pointed past the stripped $\lam$ gets decremented by 1.

\paragraph{The three steps, visualized.}
Here's what happens to each variable during the three steps of
$(\lam.\, \lam.\, 1\;\;0)\;(\texttt{bvar}\;5)$, i.e., $(\lam y.\, \lam x.\, y\;x)\;\mathit{arg}$.
Remember: stripping the outer $\lam$ gives body $= \lam.\, 1\;\;0$, where
$\texttt{bvar}\;1$ \emph{was} bound to $y$ but is now free:

\begin{center}
\begin{tikzpicture}[
  col/.style={minimum width=2.8cm, align=center, font=\small},
  header/.style={col, font=\small\bfseries},
  val/.style={col, font=\ttfamily\small},
  sepline/.style={draw, thick, black!30},
  arrow/.style={-{Stealth[length=5pt]}, thick}
]
  % Headers
  \node[header] at (0, 3.3) {Variable};
  \node[header] at (3, 3.3) {Step 1\\(shift arg up)};
  \node[header] at (6, 3.3) {Step 2\\(substitute)};
  \node[header] at (9, 3.3) {Step 3\\(shift down)};
  
  \draw[sepline] (-1.5, 2.7) -- (10.5, 2.7);
  
  % y row (target, replaced)
  \node[col, font=\small\itshape] at (0, 2.1) {$y$ (target)};
  \node[val, text=black!40] at (3, 2.1) {---};
  \node[val] at (6, 2.1) {bvar 1 $\to$ bvar 7};
  \node[val, text=black!40] at (9, 2.1) {(now part of arg)};
  \node[font=\scriptsize, text=black!50] at (11.3, 2.1) {replaced $\checkmark$};
  
  % arg row
  \node[col, font=\small\itshape] at (0, 1.2) {arg (free var)};
  \node[val] at (3, 1.2) {bvar 5 $\to$ bvar 6};
  \node[val] at (6, 1.2) {replaces bvar 1 ($y$)};
  \node[val] at (9, 1.2) {bvar 7 $\to$ bvar 6};
  \node[font=\scriptsize, text=black!50] at (11.3, 1.2) {adjusted $\checkmark$};
  
  % x row (inner lambda's own variable, stays)
  \node[col, font=\small\itshape] at (0, 0.3) {$x$ (inner $\lambda$'s param)};
  \node[val, text=black!40] at (3, 0.3) {---};
  \node[val] at (6, 0.3) {bvar 0 stays};
  \node[val] at (9, 0.3) {bvar 0 stays};
  \node[font=\scriptsize, text=black!50] at (11.3, 0.3) {untouched $\checkmark$};
  
  \draw[sepline] (-1.5, -0.2) -- (10.5, -0.2);
  
  \node[font=\small] at (4.5, -0.7) {Result: $\lambda.\, 6\;\;0$ \;=\; $\lambda x.\, \mathit{arg}\;x$ \;\; $\checkmark$};
\end{tikzpicture}
\end{center}

Each step handles a different variable's adjustment:
step 1 pre-adjusts the argument so step 3 can bring it back to the right value;
step 2 replaces $y$ (the target) with the argument;
step 3 would fix any surviving free variables whose floor was demolished---but
in this example $y$ was replaced, so only the argument's indices need adjusting.

\begin{intuition}[Why the arg goes up and then back down]
The shift-up (step 1) and shift-down (step 3) look like they cancel out for the
argument---and for the arg's own indices, they do! $5 \to 6 \to 7 \to 6$ in
this example. But the shift-down doesn't \emph{know} which indices came from
the argument and which were already in the body. It decrements \emph{all}
free variables uniformly. The shift-up pre-compensates the argument so that
when the uniform shift-down happens, the argument ends up correct \emph{and}
any surviving body variables (like an unreplaced free variable) also end up
correct. Both are handled by the same shift-down.
\end{intuition}

\begin{exercisebox}[Full beta reduction trace]
\stepcounter{exercise}\textbf{\theexercise.}\; Work through the beta reduction of $(\lam.\,0)\;(\texttt{bvar}\;0)$ (applying the identity function to a free variable).

Step 1: $\texttt{shift}\;1\;0\;(\texttt{bvar}\;0) = \texttt{bvar}\;1$ (free variable shifts up).

Step 2: $\texttt{subst}\;0\;(\texttt{bvar}\;1)\;(\texttt{bvar}\;0) = \texttt{bvar}\;1$ (index 0 matches, replaced).

Step 3: $\texttt{shift}\;(-1)\;0\;(\texttt{bvar}\;1) = \texttt{bvar}\;0$ (shift back down).

Result: $\texttt{bvar}\;0$. The identity function returned its argument unchanged. $\checkmark$
\end{exercisebox}

\begin{exercisebox}[A harder trace]
\stepcounter{exercise}\textbf{\theexercise.}\; Work through $(\lam.\, \lam.\, 1)\;(\texttt{bvar}\;0)$. This is the constant function $\lam x.\,\lam y.\, x$ applied to a free variable. The result should be $\lam.\, \texttt{bvar}\;1$ (a function that ignores its argument and returns the free variable).

\emph{Hint:} The body is $\lam.\, 1$. Start by shifting the argument, then substitute into the body (which means going under the inner $\lam$), then shift the result down.
\end{exercisebox}

And evaluation works the same way---search for a redex, fire the rule:

\begin{minted}{lean}
def step : Expr → Option Expr
  | .app (.lam body) arg => some (betaReduce body arg)
  | .app e₁ e₂ =>
    match step e₁ with
    | some e₁' => some (.app e₁' e₂)
    | none     => match step e₂ with
                  | some e₂' => some (.app e₁ e₂')
                  | none     => none
  | .lam e => match step e with
             | some e' => some (.lam e')
             | none    => none
  | _ => none

def eval (fuel : Nat) (e : Expr) : Expr :=
  match fuel with
  | 0         => e
  | fuel' + 1 => match step e with
                 | none    => e
                 | some e' => eval fuel' e'
\end{minted}

\subsubsection{Different Strategies}

With the de Bruijn evaluator working, we can explore different evaluation
strategies. Recall from Section~\ref{sec:strategies}: normal order evaluates
the outermost redex first, while call-by-value evaluates arguments before
applying functions.

\begin{exercisebox}[Strategy challenge]
\stepcounter{exercise}\textbf{\theexercise.}\; How would you modify \texttt{step} to implement call-by-value? The key difference: before firing $(\lam.\,\mathit{body})\;\mathit{arg} \reduce \ldots$, first check whether \texttt{arg} is already a value (a $\lam$). If not, reduce it first.
\end{exercisebox}

\begin{minted}{lean}
def isValue : Expr → Bool
  | .lam _ => true
  | _     => false

def cbvStep : Expr → Option Expr
  | .app (.lam body) arg =>
    if isValue arg then some (betaReduce body arg)
    else match cbvStep arg with
         | some arg' => some (.app (.lam body) arg')
         | none      => none
  | .app e₁ e₂ =>
    match cbvStep e₁ with
    | some e₁' => some (.app e₁' e₂)
    | none     => match cbvStep e₂ with
                  | some e₂' => some (.app e₁ e₂')
                  | none     => none
  | _ => none
\end{minted}

Test with $\mathbf{K}\;\mathbf{I}\;\Omega$: normal order finds the answer in
2 steps (discards $\Omega$ without evaluating it), while call-by-value diverges
(tries to evaluate $\Omega$ before passing it to $\mathbf{K}$).

%======================================================================
\section{The Simply Typed Lambda Calculus}\label{sec:stlc}
%======================================================================

The untyped $\lam$-calculus lets you write nonsense like $\Omega$. The
\textbf{simply typed lambda calculus} (STLC, $\lam_\to$) adds types to rule
out certain errors statically.

Church introduced the first typed lambda calculus in 1940, just a few years
after the untyped version. His motivation was not programming (computers barely
existed!) but logic. As Kleene and Rosser had shown in 1935, the untyped
$\lam$-calculus was inconsistent as a logic: you could derive any proposition,
including false ones, using tricks like the $\mathbf{Y}$ combinator. By adding
types, Church hoped to tame the system just enough to make it logically
consistent while keeping its computational expressiveness.

The version we study here---the ``simply typed'' calculus---is the simplest
typed $\lam$-calculus. It was formalized in its modern presentation by Curry
in the 1950s and independently by Howard in 1969 (who discovered its deep
connection to logic, which we explore in Section~\ref{sec:curryhoward}).
The cost: the STLC is not Turing-complete. But this is also its strength---
every well-typed STLC term terminates, making it safe to use as a logic.

\subsection{Types}

\begin{defnbox}[Simple Types]
\[
  \tau \;::=\; \alpha \;\mid\; \tau_1 \to \tau_2
\]
Base types $\alpha$ (like $\textsf{Nat}$, $\textsf{Bool}$) and function types.
The arrow $\to$ is right-associative: $A \to B \to C = A \to (B \to C)$.
\end{defnbox}

\subsection{Turnstile Notation and Typing Judgments}\label{sec:turnstile}

The symbol $\proves$ is the \textbf{turnstile}. It appears in \textbf{typing judgments}:

\begin{defnbox}[Reading the Turnstile $\proves$]
\[
  \Gamma \proves e : \tau
\]
reads: ``under typing context $\Gamma$, expression $e$ has type $\tau$.''

\medskip
\begin{tabular}{ll}
$\Gamma$ & \textbf{Typing context}: a list of bindings like $x{:}\textsf{Nat},\; f{:}\textsf{Nat} \to \textsf{Bool}$ \\
$\proves$ & The \textbf{turnstile}: separates assumptions (left) from conclusion (right) \\
$e : \tau$ & The conclusion: $e$ has type $\tau$ \\
\end{tabular}
\end{defnbox}

\begin{intuition}[Programmer's Analogy]
$\Gamma$ is the set of variables in scope with their types.
$\Gamma \proves e : \tau$ means ``given these type bindings, the expression
type-checks at type $\tau$.'' The empty context $\varnothing \proves e : \tau$
means $e$ is a closed, self-contained, well-typed program.
\end{intuition}

Related symbols you may encounter:
\begin{center}
\renewcommand{\arraystretch}{1.4}
\begin{tabular}{ccl}
\textbf{Symbol} & \textbf{Name} & \textbf{Meaning} \\
\hline
$\proves$ & turnstile & ``we can derive'' (syntactic) \\
$\Proves$ & double turnstile & ``semantically entails'' \\
$\nproves$ & negated turnstile & ``we cannot derive'' \\
$\entails$ & models & ``satisfies'' (model theory) \\
\end{tabular}
\end{center}

\subsection{Typing Rules as Inference Rules}\label{sec:typrules}

An inference rule has premises above the line and a conclusion below:
\[
\dfrac{\text{premise}_1 \qquad \text{premise}_2}
      {\text{conclusion}}
\;\textsc{RuleName}
\]

Read it top-down: ``if all the premises hold, then the conclusion follows.''
But there is a more useful way to read it---\textbf{bottom-up}: ``to prove the
conclusion, I need to establish these premises.'' The bottom-up reading turns
type checking into a \emph{detective story}: you start with the thing you want
to prove and work backwards, asking ``what do I need to know?'' at each step.

The STLC has three rules, one for each shape of term:

\bigskip
\begin{center}
\AxiomC{$(x : \tau) \in \Gamma$}
\RightLabel{\;\textsc{Var}}
\UnaryInfC{$\Gamma \proves x : \tau$}
\DisplayProof
\end{center}

\textsc{Var} says: ``To type a variable, look it up in the context.'' This is
the base case---no further premises need proving. It's like checking that a
variable is declared before you use it.

\begin{center}
\AxiomC{$\Gamma,\; x : \tau_1 \proves e : \tau_2$}
\RightLabel{\;\textsc{Abs}}
\UnaryInfC{$\Gamma \proves (\lam x{:}\tau_1.\, e) : \tau_1 \to \tau_2$}
\DisplayProof
\end{center}

\textsc{Abs} says: ``To type a lambda, \emph{assume} the parameter has its
declared type, add that assumption to the context, and check the body. The
result is a function type.'' Notice how $\Gamma$ grows: we \emph{extend} it
with $x : \tau_1$ when we go inside the lambda. This mirrors scoping---inside
the function body, the parameter is in scope.

\begin{center}
\AxiomC{$\Gamma \proves e_1 : \tau_1 \to \tau_2$}
\AxiomC{$\Gamma \proves e_2 : \tau_1$}
\RightLabel{\;\textsc{App}}
\BinaryInfC{$\Gamma \proves e_1\; e_2 : \tau_2$}
\DisplayProof
\end{center}

\textsc{App} says: ``To type an application, the left side must have a function
type, and the right side must match the function's domain. The result type is
the function's codomain.'' This has \emph{two} premises---both must be
established.

\subsection{Building a Proof Tree: The Detective Method}\label{sec:detective}

Reading proof trees becomes easy once you learn the bottom-up method. Let's
work through a non-trivial example in full detail.

\paragraph{The goal.} Type the term $(\lam f{:}A{\to}B.\, \lam x{:}A.\, f\;x)$.
In English: ``a function that takes a function $f$ and a value $x$, then applies
$f$ to $x$.'' We want to show
$\varnothing \proves (\lam f{:}A{\to}B.\, \lam x{:}A.\, f\;x) : (A \to B) \to A \to B$.

\paragraph{Step 1: What shape is the outermost term?} It's a lambda
$\lam f{:}A{\to}B.\, \ldots$, so the \textsc{Abs} rule must apply. That rule
says: assume $f : A \to B$, add it to the context, and check the body. Our
remaining obligation is:
\[
  f : A \to B \proves (\lam x{:}A.\, f\;x) : A \to B
\]

\paragraph{Step 2: The body is another lambda.} Shape: $\lam x{:}A.\, f\;x$.
Apply \textsc{Abs} again: assume $x : A$, add it. Remaining obligation:
\[
  f : A \to B,\; x : A \proves f\;x : B
\]

\paragraph{Step 3: Now we have an application $f\;x$.} The \textsc{App} rule
applies. It requires two things: type the left side ($f$) and type the right
side ($x$). This splits into two sub-obligations:
\begin{enumerate}[leftmargin=2em]
\item $f : A \to B,\; x : A \proves f : A \to B$ \quad(what type does $f$ have?)
\item $f : A \to B,\; x : A \proves x : A$ \quad(what type does $x$ have?)
\end{enumerate}

\paragraph{Step 4: Both are variables---use \textsc{Var}.} Look each one up in
the context $\{f : A \to B,\; x : A\}$. Both are there. Done.

\paragraph{Assembling the tree.} Now read the steps in reverse---bottom to top---
and you have the proof tree:

\begin{center}
\AxiomC{$(f : A{\to}B) \in \{f : A{\to}B,\; x : A\}$}
\RightLabel{\;\textsc{Var}}
\UnaryInfC{$f : A \to B,\; x : A \proves f : A \to B$}
\AxiomC{$(x : A) \in \{f : A{\to}B,\; x : A\}$}
\RightLabel{\;\textsc{Var}}
\UnaryInfC{$f : A \to B,\; x : A \proves x : A$}
\RightLabel{\;\textsc{App}}
\BinaryInfC{$f : A \to B,\; x : A \proves f\;x : B$}
\RightLabel{\;\textsc{Abs}}
\UnaryInfC{$f : A \to B \proves (\lam x{:}A.\, f\;x) : A \to B$}
\RightLabel{\;\textsc{Abs}}
\UnaryInfC{$\varnothing \proves (\lam f{:}A{\to}B.\, \lam x{:}A.\, f\;x) : (A \to B) \to A \to B$}
\DisplayProof
\end{center}

Read the tree upward from the bottom and watch the context grow:
\begin{itemize}[leftmargin=2em]
\item At the bottom, the context is $\varnothing$ (nothing assumed).
\item The first \textsc{Abs} peels off $f{:}A{\to}B$ and adds it:
  the context becomes $f : A \to B$.
\item The second \textsc{Abs} peels off $x{:}A$ and adds it:
  the context becomes $f : A \to B,\; x : A$.
\item The \textsc{App} and \textsc{Var} rules all operate in this full context.
\end{itemize}

Every line in the tree shows \emph{exactly} what's assumed at that point---nothing
is hidden.

\begin{intuition}[How to read any proof tree]
The recipe:
\begin{enumerate}[leftmargin=2em]
\item \textbf{Start at the bottom.} That's the statement you want to prove.
\item \textbf{Look at the rule name} on the right side of each line. It tells
  you \emph{why} this step is valid.
\item \textbf{Read upward.} Each level decomposes the problem into simpler
  pieces. Lambdas peel off a binder and extend the context. Applications split
  into two sub-problems. Variables bottom out by looking things up.
\item \textbf{The leaves (top) are axioms.} They're either variable lookups
  (\textsc{Var}) or assumptions. No premises above them---they're self-evident.
\end{enumerate}
Bottom-up, a proof tree is a plan. Top-down, it's a certificate that the plan
works. Type checkers work bottom-up (``what rule applies here?''). Humans
verifying proofs work top-down (``do these premises justify this conclusion?'').
\end{intuition}

\subsection{More Worked Examples}

\begin{workedexample}[Proof tree for the identity: $\varnothing \proves \lam x{:}A.\, x \;:\; A \to A$]
This is the simplest possible proof tree. The term has one lambda (one
\textsc{Abs}) wrapping one variable (one \textsc{Var}):
\begin{center}
\AxiomC{$(x : A) \in \{x : A\}$}
\RightLabel{\;\textsc{Var}}
\UnaryInfC{$x : A \proves x : A$}
\RightLabel{\;\textsc{Abs}}
\UnaryInfC{$\varnothing \proves (\lam x{:}A.\, x) : A \to A$}
\DisplayProof
\end{center}
\end{workedexample}

\begin{workedexample}[Proof tree for $\mathbf{K}$: $\varnothing \proves \lam x{:}A.\, \lam y{:}B.\, x \;:\; A \to B \to A$]
Two lambdas (two \textsc{Abs} steps), one variable lookup. Notice how the
context grows at each lambda: $\varnothing \to \{x:A\} \to \{x:A,\, y:B\}$.
\begin{center}
\AxiomC{$(x : A) \in \{x:A,\, y:B\}$}
\RightLabel{\;\textsc{Var}}
\UnaryInfC{$x:A,\; y:B \proves x : A$}
\RightLabel{\;\textsc{Abs}}
\UnaryInfC{$x:A \proves (\lam y{:}B.\, x) : B \to A$}
\RightLabel{\;\textsc{Abs}}
\UnaryInfC{$\varnothing \proves (\lam x{:}A.\, \lam y{:}B.\, x) : A \to B \to A$}
\DisplayProof
\end{center}
\end{workedexample}

\begin{workedexample}[Proof tree for function application]
Derive $f : A \to B,\; x : A \proves f\;x : B$:
\begin{center}
\AxiomC{$(f : A{\to}B) \in \{f : A{\to}B,\; x : A\}$}
\RightLabel{\;\textsc{Var}}
\UnaryInfC{$f : A \to B,\; x : A \proves f : A \to B$}
\AxiomC{$(x : A) \in \{f : A{\to}B,\; x : A\}$}
\RightLabel{\;\textsc{Var}}
\UnaryInfC{$f : A \to B,\; x : A \proves x : A$}
\RightLabel{\;\textsc{App}}
\BinaryInfC{$f : A \to B,\; x : A \proves f\; x : B$}
\DisplayProof
\end{center}
\end{workedexample}

\begin{workedexample}[A type error: no derivation exists]
Try typing $(\lam x{:}\textsf{Nat}.\, x)\;(\lam y{:}\textsf{Bool}.\, y)$:

The function has type $\textsf{Nat} \to \textsf{Nat}$, but the argument has
type $\textsf{Bool} \to \textsf{Bool}$. \textsc{App} requires
$\textsf{Nat} = \textsf{Bool} \to \textsf{Bool}$, which fails.
\textbf{No derivation exists}: $\nproves$

This is the type system doing its job: it rejected a ``nonsense'' application
where you'd be treating a boolean function as if it were a number.
\end{workedexample}

\subsection{Notational Shorthand in Proof Trees}\label{sec:shorthand}

The proof trees above are fully explicit: every line spells out the entire context.
In textbooks and papers, you'll almost always see a compressed version that uses a
single letter like $\Gamma$ to stand for a context that's been defined elsewhere.
Knowing how to read the compressed form---and more importantly, how to mentally
expand it---is an essential skill.

\paragraph{The convention.} Authors write something like
``Let $\Gamma = f : A \to B,\; x : A$'' and then use $\Gamma$ throughout the tree.
For example, the proof tree for
$\varnothing \proves (\lam f{:}A{\to}B.\, \lam x{:}A.\, f\;x) : (A \to B) \to A \to B$
is commonly written:

\begin{center}
\AxiomC{$(f : A{\to}B) \in \Gamma$}
\RightLabel{\;\textsc{Var}}
\UnaryInfC{$\Gamma \proves f : A \to B$}
\AxiomC{$(x : A) \in \Gamma$}
\RightLabel{\;\textsc{Var}}
\UnaryInfC{$\Gamma \proves x : A$}
\RightLabel{\;\textsc{App}}
\BinaryInfC{$\Gamma \proves f\;x : B$}
\RightLabel{\;\textsc{Abs}}
\UnaryInfC{$f : A \to B \proves (\lam x{:}A.\, f\;x) : A \to B$}
\RightLabel{\;\textsc{Abs}}
\UnaryInfC{$\varnothing \proves (\lam f{:}A{\to}B.\, \lam x{:}A.\, f\;x) : (A \to B) \to A \to B$}
\DisplayProof
\end{center}
where $\Gamma = f : A \to B,\; x : A$.

\medskip
This is exactly the same tree as the one in Section~\ref{sec:detective}, just
compressed. Every occurrence of $\Gamma$ in the upper lines stands for
$f : A \to B,\; x : A$.

\paragraph{Why authors do this.} As terms get larger, contexts accumulate
many bindings and the tree becomes unreadably wide. The $\Gamma$ convention
keeps it manageable. You'll also sometimes see $\Gamma, x : \tau$ to mean
``the context $\Gamma$ extended with $x : \tau$''---this makes \textsc{Abs}
steps especially compact.

\paragraph{How to read the compressed form.} When you encounter $\Gamma$
in a proof tree:
\begin{enumerate}[leftmargin=2em]
\item Find the definition (usually a ``where $\Gamma = \ldots$'' line nearby).
\item Mentally substitute the full context back in.
\item Verify each rule still makes sense with the full context visible.
\end{enumerate}
If you ever feel lost reading a proof tree, the first thing to try is expanding
every $\Gamma$ back to its definition.

\subsection{Key Theorems of the STLC}

The STLC enjoys three fundamental properties that together give it the
property known as \textbf{type safety}: well-typed programs don't ``go wrong.''

\begin{theorem}[Type Preservation / Subject Reduction]
If\; $\Gamma \proves e : \tau$ \;and\; $e \reduce e'$,\; then\; $\Gamma \proves e' : \tau$.
\end{theorem}

In other words, beta reduction never changes a term's type. If you start with
a well-typed expression of type $\textsf{Nat}$, every intermediate step of
evaluation will also have type $\textsf{Nat}$. Types are \emph{invariants}
of computation.

\begin{theorem}[Progress]
If\; $\varnothing \proves e : \tau$ and $e$ is not a value, then $\exists\, e'$ with $e \reduce e'$.
\end{theorem}

A well-typed closed term is never ``stuck.'' It's either already a value (a
lambda), or there's a reduction step it can take. You'll never reach a state
like $x\;y$ where nothing can happen and you don't have a result---that would
mean you have a free variable, which type checking in the empty context
prevents.

\begin{theorem}[Strong Normalization]
Every well-typed STLC term reaches a normal form. No infinite loops.
\end{theorem}

This is the deepest result. It says that every well-typed term terminates.
This is possible \emph{because} the STLC rejects self-application: you can't
type $\lam x.\, x\;x$ (you'd need $x : A$ and $x : A \to B$ simultaneously),
and self-application is the mechanism behind infinite loops ($\Omega$) and
recursion ($\mathbf{Y}$).

Together: preservation says types are stable, progress says computation doesn't
get stuck, and strong normalization says it eventually finishes. The cost: the
STLC is not Turing-complete---no infinite loops means some computable functions
can't be expressed. Sections~\ref{sec:systemf-sec}--\ref{sec:dependent-sec} explore how to get expressive
power back while keeping (some of) these safety properties.

\begin{exercisebox}[Exercises: Types]
\stepcounter{exercise}\textbf{\theexercise.}\; Write the proof tree for typing the $\mathbf{S}$ combinator:
$\varnothing \proves \lam x{:}(A \to B \to C).\,\lam y{:}(A \to B).\,\lam z{:}A.\, x\;z\;(y\;z) : (A \to B \to C) \to (A \to B) \to A \to C$.

\stepcounter{exercise}\textbf{\theexercise.}\; Why can't $\lam x.\, x\;x$ be typed in the STLC? What type would $x$ need?

\stepcounter{exercise}\textbf{\theexercise.}\; The term $\lam f{:}(A \to A).\, \lam x{:}A.\, f\;(f\;x)$ is the Church numeral $\church{2}$. What is its type? What does this tell you about the relationship between Church numerals and the STLC?
\end{exercisebox}

%======================================================================


\subsection{Implementing It: A Type Checker in Lean}\label{sec:typechecker}
%======================================================================

\subsubsection{The Motivation}

Some $\lam$-terms are ``nonsense''---$\Omega$ loops forever, $\lam x.\, x\;x$
applies a value to itself (what type would $x$ need?). Can we reject these
before running them?

\begin{exercisebox}[Type system design]
\stepcounter{exercise}\textbf{\theexercise.}\; Think about what information you would need to decide whether a term ``makes sense.'' If you see $e_1\;e_2$, what must be true about $e_1$ and $e_2$ for the application to be valid? What if you see $\lam x.\, e$---what do you need to know about $x$?
\end{exercisebox}

Think through the application case first, since it's the one that goes wrong.
In $e_1\;e_2$, the term $e_1$ is being \emph{called as a function}. So $e_1$
must have a function type: it takes some input type $A$ and produces some output
type $B$. The term $e_2$ is the argument, so it must have type $A$---matching
the function's input. The result has type $B$.

For a lambda $\lam x.\, e$: you need to know what type $x$ has (this is the
annotation), and then you check whether the body $e$ type-checks under the
assumption that $x$ has that type. If the body has type $B$, then the whole
lambda has type $A \to B$.

For a variable: you need to look it up in a context---a list of ``which
variables have which types.'' With de Bruijn indices, the context is just a
list of types, indexed by position.

This gives you three rules, one per constructor---exactly the inference rules
from Section~\ref{sec:stlc}. Each rule becomes one branch of the type checker:

\begin{minted}{lean}
inductive Ty where
  | base  : String → Ty
  | arrow : Ty → Ty → Ty
  deriving Repr, DecidableEq

-- Typed terms: lambdas carry type annotations
inductive TExpr where
  | bvar : Nat → TExpr
  | lam  : Ty → TExpr → TExpr    -- annotation on the parameter
  | app  : TExpr → TExpr → TExpr
  deriving Repr
\end{minted}

\begin{minted}{lean}
-- Γ ⊢ e : τ    (Γ is List Ty, index 0 = innermost binder)
def typeCheck (ctx : List Ty) : TExpr → Option Ty
  | .bvar n          => ctx.get? n                -- Var: look up in Γ
  | .lam paramTy body =>                          -- Abs: extend Γ, check body
    match typeCheck (paramTy :: ctx) body with
    | some bodyTy => some (.arrow paramTy bodyTy)
    | none        => none
  | .app e₁ e₂ => do                             -- App: match domain type
    let funTy ← typeCheck ctx e₁
    let argTy ← typeCheck ctx e₂
    match funTy with
    | .arrow paramTy retTy =>
      if paramTy == argTy then some retTy else none
    | _ => none
\end{minted}

\begin{intuition}[Code $\leftrightarrow$ Math]
Each branch implements one inference rule from Section~\ref{sec:typrules}:

\texttt{.bvar n} $\leftrightarrow$ \textsc{Var} (lookup in $\Gamma$), \quad
\texttt{.lam} $\leftrightarrow$ \textsc{Abs} (extend $\Gamma$, check body), \quad
\texttt{.app} $\leftrightarrow$ \textsc{App} (match domain type).

Returning \texttt{some $\tau$} means $\Gamma \proves e : \tau$;
returning \texttt{none} means the term is ill-typed.
\end{intuition}

%======================================================================
\section{Scott Encodings}\label{sec:scott}
%======================================================================

\subsection{An Alternative to Church}

Church encodings represent data as \emph{folds}: a Church numeral iterates
$f$ a total of $n$ times. This is elegant but has a practical problem:
some operations are expensive. Predecessor, as we saw, requires $O(n)$ work
to ``rebuild'' the count from scratch. Pattern matching on a Church-encoded
sum type similarly requires deconstructing the entire fold.

\textbf{Scott encodings}, introduced by Dana Scott in the 1960s, take a
different approach. Scott was one of the architects of denotational semantics---
the project of giving mathematical meaning to programming languages. Where
Church asked ``how do I encode data as iterable functions?'', Scott asked
``how do I encode data as pattern-matchable functions?'' The answer turned out
to be simpler and more efficient for many operations, and it is closer to
what real compilers actually do internally.

Each constructor becomes a
function that takes one handler per constructor and calls the appropriate one.
This is essentially what a \texttt{match} or \texttt{case} expression does
at runtime in most compilers---so Scott encodings are closer to how data is
actually implemented in practice.

\subsection{Reading the Notation}

Scott encodings introduce a few notational conventions that differ from the
Church encoding sections:

\begin{center}
\renewcommand{\arraystretch}{1.6}
\begin{tabular}{cp{10cm}}
\textbf{Notation} & \textbf{Meaning} \\
\hline
$\church{n}_S$ &
  The Scott encoding of the number $n$. The subscript $S$ distinguishes it from
  the Church numeral $\church{n}$. Without the subscript, $\church{n}$ always
  means the Church version. \\

$\textsc{Pred}_S$ &
  The Scott version of the predecessor function. Again, the subscript
  distinguishes it from the Church predecessor $\textsc{Pred}$. \\

$\lam z\, s.\, \ldots$ &
  The parameters $z$ and $s$ stand for ``zero handler'' and ``successor
  handler.'' A Scott numeral takes these two arguments and calls whichever one
  matches its shape---$z$ if it's zero, $s$ if it's a successor. \\

$\lam n\, c.\, \ldots$ &
  For Scott lists, $n$ stands for ``nil handler'' and $c$ for ``cons handler.''
  Same idea: the list calls whichever handler matches. \\

$\church{n+1}_S$ &
  The Scott encoding of ``the successor of $n$.'' This is a \emph{pattern} in
  the definition, not an expression to evaluate. It says: ``any Scott numeral
  that represents a non-zero value has this form.'' \\
\end{tabular}
\end{center}

The key mental model: a Scott-encoded value is a function that takes
\textbf{one handler per constructor} and calls the matching one.
For natural numbers, there are two constructors (zero and successor),
so every Scott numeral takes two arguments. For lists, there are two
constructors (nil and cons), so every Scott list takes two arguments. This is
directly analogous to a \texttt{match} or \texttt{case} expression in a
programming language:

\begin{minted}{lean}
-- A Scott numeral IS this match, packaged as a function
match n with
| zero   => z            -- call the zero handler
| succ p => s p          -- call the successor handler with predecessor p
\end{minted}

\subsection{Scott Booleans}

Scott booleans look identical to Church booleans (both are choosing between
two options):
\begin{align*}
\textsc{STrue}  &\triangleq \lam t\, f.\, t &
\textsc{SFalse} &\triangleq \lam t\, f.\, f
\end{align*}

No difference here---the distinction appears with more complex types.

\subsection{Scott Numerals}

A natural number is either zero or the successor of another number. In a
pattern match, we provide one handler for each case:

\begin{defnbox}[Scott Numerals]
\vspace{-0.8em}
\begin{align*}
\church{0}_S &\triangleq \lam z\, s.\, z \\
\church{n+1}_S &\triangleq \lam z\, s.\, s\;\church{n}_S
\end{align*}
\end{defnbox}

Zero takes two handlers and calls the first ($z$). A successor $n+1$ takes
two handlers and calls the second ($s$) with the \textbf{predecessor}
$\church{n}_S$ as an argument. This is the crucial difference from Church
encoding: the predecessor is directly available, not reconstructed.

\begin{align*}
\church{0}_S &= \lam z\, s.\, z \\
\church{1}_S &= \lam z\, s.\, s\;\church{0}_S \\
\church{2}_S &= \lam z\, s.\, s\;\church{1}_S \\
\church{3}_S &= \lam z\, s.\, s\;\church{2}_S
\end{align*}

Now $\textsc{Pred}$ is trivial---just call the number with handlers that
extract the predecessor:
\[
\textsc{Pred}_S \triangleq \lam n.\, n\; \church{0}_S\; (\lam p.\, p)
\]

If $n$ is zero, return zero. If $n = s\;p$, the successor handler receives $p$
directly. Compare this to the complex pair-walking trick needed for Church
numerals!

\begin{workedexample}[$\textsc{Pred}_S\;\church{3}_S = \church{2}_S$]
\begin{align*}
  \textsc{Pred}_S\;\church{3}_S
  &= (\lam n.\, n\;\church{0}_S\;(\lam p.\, p))\;(\lam z\, s.\, s\;\church{2}_S)
    \rmark{expand definitions} \\
  &\reduce (\lam z\, s.\, s\;\church{2}_S)\;\church{0}_S\;(\lam p.\, p)
    \rmark{plug in $\church{3}_S$} \\
  &\reduces (\lam p.\, p)\;\church{2}_S
    \rmark{it's a successor, so call $s$ with the predecessor} \\
  &\reduce \church{2}_S \quad\checkmark
    \rmark{identity returns it unchanged}
\end{align*}
\end{workedexample}

\begin{intuition}[Church vs.\ Scott]
Church encoding: ``I am the number 3 because I can apply things 3 times.''
You know \emph{how to iterate} but not how to decompose.

Scott encoding: ``I am the number 3 because I am the successor of 2.''
You know \emph{how to pattern-match} and have direct access to subcomponents.

Church encodings give you \texttt{foldr} for free. Scott encodings give you
\texttt{case} for free. Most real compilers (GHC, OCaml) use representations
closer to Scott encoding internally, because pattern matching is the more
fundamental operation---you can build folds from pattern matching, but building
efficient pattern matching from folds is harder.
\end{intuition}

\subsection{Scott-Encoded Lists}

\begin{align*}
\textsc{SNil}  &\triangleq \lam n\, c.\, n \\
\textsc{SCons} &\triangleq \lam h\, t\, n\, c.\, c\; h\; t
\end{align*}

$\textsc{SNil}$ selects the nil handler. $\textsc{SCons}\;h\;t$ selects the
cons handler and passes it the head $h$ and tail $t$ directly. Extracting the
head or tail is a single reduction step---no folding required.

%======================================================================
\section{The Curry--Howard Correspondence}\label{sec:curryhoward}
%======================================================================

\subsection{Types Are Propositions; Programs Are Proofs}

There is a deep and beautiful connection between typed lambda calculus and
logic. It is arguably the most profound idea in this entire document.

The story unfolds over three decades. In 1934, Haskell Curry noticed that the
types of certain combinators corresponded to axiom schemes in propositional
logic---for instance, the type of $\mathbf{K}$, namely $A \to B \to A$,
is a well-known logical axiom. He noted this as a curiosity but did not develop
it further.

In 1958, Curry and his collaborator Robert Feys wrote up the observation more
carefully in their monumental textbook \emph{Combinatory Logic, Volume I}. But
the full picture remained unclear until 1969, when William Howard circulated an
unpublished manuscript (finally published in 1980) making the correspondence
explicit and systematic: types \emph{are} logical propositions, and programs
\emph{are} proofs. Every rule of typed lambda calculus corresponds precisely to
a rule of natural deduction in logic. Howard extended the correspondence beyond
simple types to include universal quantification, existential types, and
disjunction.

The correspondence says:

\begin{center}
\renewcommand{\arraystretch}{1.5}
\begin{tabular}{ll}
\textbf{Type theory} & \textbf{Logic} \\
\hline
Type $\tau$               & Proposition \\
Term $e : \tau$           & Proof of the proposition \\
Function type $A \to B$   & Implication ``$A$ implies $B$'' \\
Inhabited type            & Provable proposition \\
Empty type (no terms)     & False / unprovable proposition \\
$\beta$-reduction         & Proof simplification \\
\end{tabular}
\end{center}

This is not a loose analogy. It is a precise, formal correspondence. Every
valid typing derivation in the STLC \emph{is} a proof in propositional logic,
and every proof \emph{is} a well-typed program. This correspondence is the
intellectual foundation of modern proof assistants---Lean, Coq, Agda, and
Idris all exploit it directly, verifying mathematical proofs by type-checking
programs.

\subsection{Examples}

\paragraph{The identity function proves $A \Rightarrow A$.}
The term $\lam x{:}A.\, x$ has type $A \to A$. Under Curry--Howard, this is
a proof that ``if $A$, then $A$''---trivially true, since assuming $A$ gives
you $A$. The proof is the program: take the evidence and return it unchanged.

\paragraph{The K combinator proves $A \Rightarrow B \Rightarrow A$.}
$\lam x{:}A.\, \lam y{:}B.\, x$ has type $A \to B \to A$. This proves
``if $A$, then (if $B$, then $A$)''---given evidence for $A$ and $B$, you
can produce evidence for $A$ by ignoring $B$. The ``weakening'' or
``thinning'' rule of logic.

\paragraph{The S combinator proves $(A \Rightarrow B \Rightarrow C) \Rightarrow (A \Rightarrow B) \Rightarrow A \Rightarrow C$.}
This is the logical axiom scheme known as the \emph{distribution} of
implication. The program witnesses the proof: given a function from $A$ and $B$
to $C$, and a function from $A$ to $B$, and an $A$, produce a $C$ by applying
them in the right combination.

\paragraph{Modus ponens is function application.}
The logical rule ``from $A$ and $A \Rightarrow B$, conclude $B$'' corresponds
exactly to the \textsc{App} typing rule: if you have $e_1 : A \to B$ and
$e_2 : A$, then $e_1\;e_2 : B$.

\begin{intuition}[Why this matters]
The Curry--Howard correspondence is the foundation of proof assistants like
Lean, Coq, and Agda. When you write a proof in Lean, you are literally writing
a lambda calculus program. The type checker verifies the proof by type-checking
the program. This is why our lambda calculus interpreter and Lean itself are
built on the same ideas: they \emph{are} the same ideas.

The correspondence also explains why the STLC is not Turing-complete: every
well-typed STLC term normalizes (strong normalization theorem). Under
Curry--Howard, this means every proof terminates---which is a property you
want! A logic where you can prove anything (including false) by writing an
infinite loop would be useless as a logic.
\end{intuition}

\begin{workedexample}[Reading a proof tree as a logical derivation]
The typing derivation for $\lam x{:}A.\, \lam y{:}B.\, x : A \to B \to A$
is simultaneously a derivation in propositional logic:

\begin{center}
\AxiomC{$(A) \in \text{assumptions}$}
\RightLabel{\;\textsc{Assumption}}
\UnaryInfC{$A$}
\RightLabel{\;\textsc{$\Rightarrow$-Intro} (discharge $B$)}
\UnaryInfC{$B \Rightarrow A$}
\RightLabel{\;\textsc{$\Rightarrow$-Intro} (discharge $A$)}
\UnaryInfC{$A \Rightarrow B \Rightarrow A$}
\DisplayProof
\end{center}

This is exactly the same tree structure as the \textsc{Var}/\textsc{Abs} typing
derivation. The \textsc{Abs} rule in type theory is the
$\Rightarrow$-Introduction rule in logic. The \textsc{Var} rule is the
assumption rule.
\end{workedexample}

\begin{exercisebox}[Exercises: Going Deeper]
\stepcounter{exercise}\textbf{\theexercise.}\; Define $\textsc{Succ}_S$ for Scott numerals and verify that $\textsc{Succ}_S\;\church{2}_S = \church{3}_S$.

\stepcounter{exercise}\textbf{\theexercise.}\; Define $\textsc{IsZero}_S$ for Scott numerals. How much simpler is it compared to the Church version?

\stepcounter{exercise}\textbf{\theexercise.}\; Under Curry--Howard, what logical proposition does the type $(A \to B) \to (B \to C) \to (A \to C)$ correspond to? Write a $\lam$-term that has this type.

\stepcounter{exercise}\textbf{\theexercise.}\; Why is there no closed term of type $A \to B$ (where $A$ and $B$ are distinct base types) in the STLC? What does this say about the corresponding logical proposition?

\stepcounter{exercise}\textbf{\theexercise.}\; The type $\forall \alpha.\, \alpha$ (from System F, not the STLC) is uninhabited. Under Curry--Howard, what proposition does it represent?
\end{exercisebox}

%======================================================================

%======================================================================
\section{Intuitionistic Logic}\label{sec:intuitionistic-sec}
%======================================================================

Before we can understand dependent types, we need to understand the
philosophical revolution they grew out of. This detour into the foundations of
mathematics may seem far from programming, but it leads directly to the ideas
that make Lean, Coq, and Agda work.


\subsection{The Classical View}


Before we can understand dependent types, we need to understand the
philosophical revolution they grew out of. This detour into the foundations of
mathematics may seem far from programming, but it leads directly to the ideas
that make Lean, Coq, and Agda work.

\paragraph{The classical view.}
In classical logic---the logic most people learn in school and most
mathematicians use daily---every proposition is either true or false, regardless
of whether anyone can prove it. The Goldbach conjecture (``every even number
greater than 2 is the sum of two primes'') is either true or false \emph{right
now}, even though nobody has proved it either way. Truth exists independently
of our ability to verify it.

Classical logic gives you several powerful principles:

\begin{center}
\renewcommand{\arraystretch}{1.4}
\begin{tabular}{lp{8cm}}
\textbf{Law of Excluded Middle (LEM)} & For any proposition $P$: either $P$ or
  $\lnot P$. There is no middle ground. \\
\textbf{Double Negation Elim.\ (DNE)} & If $\lnot\lnot P$, then $P$. If it's
  impossible for $P$ to be false, then $P$ is true. \\
\textbf{Proof by Contradiction} & To prove $P$, assume $\lnot P$ and derive a
  contradiction. Since $\lnot P$ led to absurdity, $P$ must be true. \\
\end{tabular}
\end{center}

These feel obviously correct, and mathematicians use them constantly. But the
constructivists noticed something unsettling about them.

\paragraph{The constructivist objection.}
Consider this classical proof: ``There exist irrational numbers $a$ and $b$
such that $a^b$ is rational.''

\emph{Proof.} We know $\sqrt{2}$ is irrational. Now consider the number
$\sqrt{2}^{\sqrt{2}}$---we have no idea whether this is rational or irrational.
But by the law of excluded middle, it must be one or the other.

\emph{Case 1:} Suppose $\sqrt{2}^{\sqrt{2}}$ happens to be rational.
Then take $a = \sqrt{2}$ and $b = \sqrt{2}$. Both are irrational, and
$a^b = \sqrt{2}^{\sqrt{2}}$, which we assumed is rational. Done.

\emph{Case 2:} Suppose $\sqrt{2}^{\sqrt{2}}$ is irrational.
Then take $a = \sqrt{2}^{\sqrt{2}}$ (irrational by our assumption in this case)
and $b = \sqrt{2}$ (irrational). Now compute $a^b$ using the exponent law
$(x^y)^z = x^{y \cdot z}$:
\[
  a^b \;=\; \bigl(\sqrt{2}^{\sqrt{2}}\bigr)^{\sqrt{2}}
      \;=\; \sqrt{2}^{\;\sqrt{2}\,\cdot\,\sqrt{2}}
      \;=\; \sqrt{2}^{\;2}
      \;=\; 2
\]
And $2$ is rational. Done.

In both cases we found irrational $a$, $b$ with $a^b$ rational. So such a
pair exists. $\square$

The proof is valid in classical logic. But notice: \emph{we don't know which
case holds}. Is $\sqrt{2}^{\sqrt{2}}$ rational or irrational? The proof doesn't
tell us---it just says ``one of these two cases must be true, and either way
we win.'' We proved ``there exist $a$ and $b$'' without being able to point to
\emph{which} $a$ and $b$ actually work. (It turns out $\sqrt{2}^{\sqrt{2}}$ is
irrational, by the Gelfond--Schneider theorem---but proving that requires
entirely different and much harder mathematics that this proof never invokes.)

We have an existence proof that doesn't produce the thing whose existence it
claims.

For the constructivists---led by L.~E.~J.\ Brouwer in the 1910s and 1920s---this
is not a real proof. Brouwer argued that mathematical objects don't exist in
some Platonic realm waiting to be discovered. They exist only when a
mathematician \emph{constructs} them. To prove ``there exists an $x$ with
property $P$,'' you must exhibit a specific $x$ and demonstrate that it has
property $P$. A proof by contradiction that merely shows non-existence is
absurd---without pointing to a witness---is not a valid existence proof.

\paragraph{What intuitionistic logic rejects.}
Brouwer's student Arend Heyting (1930) formalized these ideas into
\textbf{intuitionistic logic}: a logic that modifies the classical rules to
enforce constructivity. The key differences:

\begin{center}
\renewcommand{\arraystretch}{1.4}
\begin{tabular}{lcc}
\textbf{Principle} & \textbf{Classical} & \textbf{Intuitionistic} \\
\hline
$P \lor \lnot P$ (excluded middle)       & \checkmark & \textsf{not assumed} \\
$\lnot\lnot P \Rightarrow P$ (double negation) & \checkmark & \textsf{not valid} \\
Proof by contradiction for existence     & \checkmark & \textsf{not valid} \\
$\lnot\lnot\lnot P \Rightarrow \lnot P$ & \checkmark & \checkmark \\
Everything else in propositional logic   & \checkmark & \checkmark \\
\end{tabular}
\end{center}

Intuitionistic logic does not say excluded middle is \emph{false}. It says
excluded middle is \emph{not assumed}---it might hold for specific propositions
(e.g., ``$n$ is even or $n$ is odd'' is constructively valid because you can
check by dividing), but you can't assert it for \emph{every} proposition by
default.

\paragraph{The BHK interpretation.}
Brouwer, Heyting, and Kolmogorov developed an interpretation of what proofs
\emph{mean} in intuitionistic logic, now called the \textbf{BHK interpretation}:

\begin{center}
\renewcommand{\arraystretch}{1.4}
\begin{tabular}{lp{9cm}}
\textbf{To prove\ldots} & \textbf{You must provide\ldots} \\
\hline
$A \land B$ (``$A$ and $B$'') & A proof of $A$ \emph{and} a proof of $B$
  (a pair). \\
$A \lor B$ (``$A$ or $B$'') & A proof of $A$ \emph{or} a proof of $B$,
  together with a tag saying \emph{which} (a tagged union). \\
$A \Rightarrow B$ (``$A$ implies $B$'') & A \emph{method} that transforms any
  proof of $A$ into a proof of $B$ (a function). \\
$\exists x.\, P(x)$ (``there exists'') & A specific witness $a$ together with
  a proof that $P(a)$ holds (a dependent pair). \\
$\forall x.\, P(x)$ (``for all'') & A method that, given any $x$, produces a
  proof of $P(x)$ (a dependent function). \\
$\lnot A$ (``not $A$'') & A method that transforms any proof of $A$ into a
  proof of $\bot$ (absurdity)---i.e., a proof that $A$ leads to contradiction. \\
\end{tabular}
\end{center}

Read that table again carefully. Do the right-hand descriptions remind you of
anything?

\emph{A pair. A tagged union. A function. A dependent pair. A dependent
function.} These are \textbf{data types}. The BHK interpretation, formulated
decades before computers, is describing \emph{programs}. A proof of
``$A$ and $B$'' is a pair. A proof of ``$A$ implies $B$'' is a function. A
proof of ``there exists $x$ with $P(x)$'' is a struct with a field for the
witness and a field for the evidence.

\paragraph{Why AND is a product type (and multiplication).}
A proof of $A \land B$ consists of a proof of $A$ \emph{and} a proof of $B$,
packaged together as a pair $(a, b)$ where $a : A$ and $b : B$. The type of
such pairs is the \textbf{product type} $A \times B$.

Why ``product''? Because the number of values in $A \times B$ is the
\emph{product} of the number of values in $A$ and $B$. If
$A = \textsf{Bool}$ (2 values) and $B = \textsf{Bool}$ (2 values), then
$A \times B$ has $2 \times 2 = 4$ values:
$(\texttt{true},\texttt{true})$,
$(\texttt{true},\texttt{false})$,
$(\texttt{false},\texttt{true})$,
$(\texttt{false},\texttt{false})$.

In Lean, the product type is \texttt{Prod A B} (written \texttt{A × B}), and
the conjunction $A \land B$ is defined as a structure with two fields---which
is exactly a product type. They are the same thing.

\paragraph{Why OR is a sum type (and addition).}
A proof of $A \lor B$ consists of \emph{either} a proof of $A$ \emph{or} a
proof of $B$, together with a tag saying which one you're providing. This is
a \textbf{sum type} (also called a tagged union, disjoint union, or coproduct):
$A + B$.

Why ``sum''? Because the number of values in $A + B$ is the \emph{sum} of the
number of values in $A$ and $B$. If $A = \textsf{Bool}$ (2 values) and
$B = \textsf{Unit}$ (1 value), then $A + B$ has $2 + 1 = 3$ values:
$\texttt{inl}(\texttt{true})$, $\texttt{inl}(\texttt{false})$,
$\texttt{inr}(\texttt{()})$.

In Lean, disjunction $A \lor B$ is an inductive type with two constructors
\texttt{Or.inl : A → A ∨ B} and \texttt{Or.inr : B → A ∨ B}---which is
exactly a sum type.

\paragraph{Why implication is a function type (and exponentiation).}
A proof of $A \Rightarrow B$ is a method that transforms any proof of $A$
into a proof of $B$---a function $A \to B$. The number of functions from $A$
to $B$ is $|B|^{|A|}$---exponentiation. If $A = \textsf{Bool}$ (2 values) and
$B = \textsf{Bool}$ (2 values), there are $2^2 = 4$ functions from
\textsf{Bool} to \textsf{Bool}: always-true, always-false, identity, and
negation.

This gives us a complete arithmetic of types:
\begin{center}
\renewcommand{\arraystretch}{1.4}
\begin{tabular}{lccc}
\textbf{Logic} & \textbf{Type} & \textbf{Arithmetic} & \textbf{Lean} \\
\hline
$A \land B$ & $A \times B$ (product) & $|A| \cdot |B|$ & \texttt{A × B} \\
$A \lor B$ & $A + B$ (sum) & $|A| + |B|$ & \texttt{A ⊕ B}, \texttt{Sum} \\
$A \Rightarrow B$ & $A \to B$ (function) & $|B|^{|A|}$ & \texttt{A → B} \\
$\top$ (true) & $\textsf{Unit}$ & $1$ & \texttt{Unit} \\
$\bot$ (false) & $\textsf{Empty}$ & $0$ & \texttt{Empty} \\
$\lnot A$ & $A \to \textsf{Empty}$ & $0^{|A|}$ & \texttt{A → Empty} \\
\end{tabular}
\end{center}

The last row is illuminating: $0^{|A|} = 0$ when $|A| > 0$ (there are no
functions from a nonempty type to an empty type) and $0^0 = 1$ (there is
exactly one function from the empty type to itself---the identity on nothing).
This matches the logic: $\lnot A$ is unprovable when $A$ is provable, and
$\lnot \bot$ has exactly one proof.

Even the laws of arithmetic correspond to logical laws: $A \times (B + C) =
A \times B + A \times C$ is distributivity of AND over OR. $A \to (B \times C)
\cong (A \to B) \times (A \to C)$ says a function returning a pair is the same
as a pair of functions---which is the logical tautology
$(A \Rightarrow B \land C) \Leftrightarrow (A \Rightarrow B) \land
(A \Rightarrow C)$.

This is the Curry--Howard correspondence, seen from the logic side. The
constructivists were describing computation before computation was formalized.

\paragraph{Why programmers should care.}
The connection to programming is not a metaphor. In a dependently typed
proof assistant like Lean:

\begin{minted}{lean}
-- "A and B" is a struct (pair)
theorem and_intro (ha : A) (hb : B) : A ∧ B := ⟨ha, hb⟩

-- "A or B" is a tagged union (sum type)
theorem or_left (ha : A) : A ∨ B := Or.inl ha

-- "A implies B" is a function
theorem modus_ponens (hab : A → B) (ha : A) : B := hab ha

-- "there exists n such that n > 5" needs a witness
theorem exists_gt_5 : ∃ n : Nat, n > 5 := ⟨6, by omega⟩
\end{minted}

The proof of \texttt{exists\_gt\_5} must provide the number $6$---an actual
witness---and a proof that $6 > 5$. You cannot prove it by contradiction. You
cannot say ``suppose no such $n$ exists'' and derive absurdity. You must hand
over the goods.

This is intuitionistic logic in action. Lean's kernel is an intuitionistic type
theory. Every proof must construct its evidence. This is why proofs in Lean
are programs: the BHK interpretation \emph{is} the type system.

\begin{intuition}[Classical logic in Lean]
Lean does allow classical reasoning---but only when you explicitly opt in by
importing \texttt{Classical}. When you write \texttt{Classical.em P}
(excluded middle for $P$), or use the \texttt{by\_contra} tactic, Lean adds
the axiom of choice or excluded middle as an additional assumption. The proof
still type-checks, but it is marked as ``noncomputable''---Lean can verify it
but cannot extract a runnable program from it, because the proof may use
existence claims without witnesses. This is the type-theoretic manifestation
of the constructivist objection: classical proofs don't always compute.
\end{intuition}

\paragraph{Why ``intuitionistic''?}
The name comes from Brouwer's philosophy of \textbf{intuitionism}, which held
that mathematics is a creation of the human mind, not a description of an
external Platonic reality. Mathematical objects exist only when constructed by
mental activity. A proof must be a mental construction that can, in principle,
be carried out step by step. Brouwer himself was quite radical---he rejected
formalization entirely and thought mathematics should be done purely through
intuition, without symbols. The irony is that his followers formalized his
ideas into one of the most precisely axiomatized logical systems in existence,
and that formalization became the foundation of machine-checked proof.


%======================================================================
\section{System F: Polymorphism}\label{sec:systemf-sec}
%======================================================================


\subsection{The Problem of Repetition}

The STLC has an embarrassing limitation. Consider the identity function. In the
untyped calculus, $\lam x.\, x$ works on anything---numbers, booleans, other
functions. But in the STLC, types are rigid: you need $\lam x{:}\textsf{Nat}.\, x$
for natural numbers, $\lam x{:}\textsf{Bool}.\, x$ for booleans,
$\lam x{:}(\textsf{Nat} \to \textsf{Nat}).\, x$ for functions on naturals,
and so on---infinitely many copies of the ``same'' function. This isn't a minor
inconvenience. It means fundamental operations like composition,
$\lam f.\, \lam g.\, \lam x.\, f\;(g\;x)$, cannot be given a single type. You
cannot even type Church numerals properly: $\church{2}$ applies its first
argument twice, but the STLC insists on knowing \emph{which type} the argument
has.

The question is: can we add ``generics'' to the type system without losing the
safety properties (progress, preservation, strong normalization) that make the
STLC useful?

\subsection{Two Independent Discoveries}

The answer came from two completely different intellectual traditions, arriving
at the same system independently.

\paragraph{Girard (1972): from proof theory.} Jean-Yves Girard, a French
logician, was working on the proof theory of second-order arithmetic---a system
where you can quantify not just over numbers (``for all $n$, \ldots'') but over
\emph{properties} of numbers (``for all predicates $P$, \ldots''). He needed a
typed lambda calculus that could express this second-order quantification, and he
constructed what he called ``System F'' (the ``F'' was simply the next letter
in his naming sequence---his earlier systems were called System D and System E).

Girard's breakthrough was proving that System F satisfies strong normalization.
This was a major result in proof theory: it gave a constructive consistency
proof for second-order arithmetic, extending the work of G\"odel and Gentzen.
The proof technique---now called the ``Girard--Tait method'' or ``method of
reducibility candidates''---involved defining a semantic notion of
``well-behaved term'' and showing by induction that all typed terms qualify.

\paragraph{Reynolds (1974): from programming.} John Reynolds, an American
computer scientist, was independently studying how to formalize the idea
that some programs are ``generic'' in their type. He defined what he called
the ``polymorphic lambda calculus'' with the same syntax and typing rules
as Girard's System F, but motivated entirely by programming: a polymorphic
function should work uniformly across all types, without inspecting which
type it receives.

Reynolds also discovered a deep consequence of this uniformity, later refined
by Philip Wadler (1989) into the celebrated \textbf{parametricity theorem}
(``theorems for free''): the type of a polymorphic function severely constrains
what it can do. A function of type $\forall\alpha.\, \alpha \to \alpha$ \emph{must}
be the identity---there is literally nothing else a polymorphic function of
that type can do, because it cannot inspect its argument's type. The type alone
proves the behavior.

\subsection{The Key Idea}

System F allows functions to take \emph{types} as arguments, not just
values. This lets you write a single identity function that works at any type:

\[
\Lam\alpha.\, \lam x{:}\alpha.\, x \;\;:\;\; \forall\alpha.\, \alpha \to \alpha
\]

The $\Lam\alpha$ (capital lambda) abstracts over a type variable $\alpha$.
The $\forall\alpha$ in the type says ``for any type $\alpha$.'' When you
\emph{apply} this function, you first supply a type, then a value:

\begin{align*}
  (\Lam\alpha.\, \lam x{:}\alpha.\, x)\;[\textsf{Nat}]\;5
    &\reduce (\lam x{:}\textsf{Nat}.\, x)\;5 \\
    &\reduce 5
\end{align*}

The $[\textsf{Nat}]$ is a \textbf{type application}: it instantiates $\alpha$
to $\textsf{Nat}$, giving back the Nat-specific identity function.

\begin{intuition}[System F in programming languages]
System F is the theoretical foundation for generics. When you write
\texttt{fn identity<T>(x: T) -> T} in Rust, you are writing a System F term.
The angle brackets \texttt{<T>} are $\Lam T$. The \texttt{where} clauses of
Rust traits, Haskell type classes, and Java's \texttt{<T extends Comparable>}
are all refinements of System F's type abstraction.

Most practical languages use a restricted form called \textbf{Hindley--Milner}
(HM) type inference, discovered independently by Hindley (1969) and Milner
(1978). HM restricts where $\forall$ can appear (only at the outermost level of
a let-binding) but in exchange gives you \emph{complete type inference}: the
compiler can figure out all the types without any annotations. This is why
Haskell and ML rarely need type signatures---the HM algorithm infers them
automatically. Full System F, by contrast, has undecidable type inference
(proved by Joe Wells, 1999): you sometimes \emph{must} provide annotations.
\end{intuition}

\subsection{The Philosophical Significance}

System F reveals something deep: \textbf{data can emerge from pure logic}.
In System F, you can encode natural numbers, booleans, pairs, and lists
\emph{entirely within the type system}:

\[
  \textsf{Nat}_F \triangleq \forall\alpha.\, (\alpha \to \alpha) \to \alpha \to \alpha
\]

This is the type of Church numerals! In the STLC, Church numerals have
type $(A \to A) \to A \to A$ for a \emph{fixed} type $A$. In System F, the
$\forall\alpha$ lets a single numeral work at any type. Similarly:
\begin{align*}
  \textsf{Bool}_F &\triangleq \forall\alpha.\, \alpha \to \alpha \to \alpha \\
  \textsf{Pair}_F\;A\;B &\triangleq \forall\alpha.\, (A \to B \to \alpha) \to \alpha
\end{align*}

These aren't tricks---they are the \emph{natural types} that arise when you ask
``what is the most general type a Church boolean can have?'' The answer
\emph{is} the universally quantified type above. The data types we take for
granted in programming (integers, booleans, lists) can all be understood as
patterns of polymorphic function application. From the perspective of proof
theory, Girard's encoding showed that second-order logic is rich enough to
define all the data structures of mathematics internally.

\subsection{Implementing System F}

System F extends the STLC with two new constructs: \textbf{type abstraction}
($\Lam\alpha.\, e$) and \textbf{type application} ($e\;[\tau]$). The types
grow too: we add $\forall\alpha.\, \tau$ (universally quantified types) and
type variables $\alpha$.

\begin{minted}{lean}
-- Types: add type variables and universal quantification
inductive FTy where
  | tvar  : Nat → FTy                   -- type variable (de Bruijn)
  | base  : String → FTy                -- base types like Nat, Bool
  | arrow : FTy → FTy → FTy             -- τ₁ → τ₂
  | forall_ : FTy → FTy                 -- ∀α. τ  (α is index 0 in τ)
  deriving Repr, DecidableEq

-- Terms: add type abstraction and type application
inductive FExpr where
  | bvar  : Nat → FExpr                 -- term variable (de Bruijn)
  | lam   : FTy → FExpr → FExpr         -- λ(x:τ). e
  | app   : FExpr → FExpr → FExpr       -- e₁ e₂
  | tlam  : FExpr → FExpr               -- Λα. e  (type abstraction)
  | tapp  : FExpr → FTy → FExpr         -- e [τ]  (type application)
  deriving Repr
\end{minted}

The type checker needs two operations on types: \textbf{shifting} and
\textbf{substitution}---the same operations we needed for terms with de Bruijn
indices, but now operating on type variables inside types.

\begin{minted}{lean}
-- Shift type variables in a type
def FTy.shift (d : Int) (c : Nat) : FTy → FTy
  | .tvar n      => if n ≥ c then .tvar (Int.toNat (n + d)) else .tvar n
  | .base s      => .base s
  | .arrow t₁ t₂ => .arrow (t₁.shift d c) (t₂.shift d c)
  | .forall_ t   => .forall_ (t.shift d (c + 1))  -- under a ∀: bump cutoff

-- Substitute a type for type variable k
def FTy.subst (k : Nat) (s : FTy) : FTy → FTy
  | .tvar n      => if n == k then s
                     else if n > k then .tvar (n - 1) else .tvar n
  | .base s      => .base s
  | .arrow t₁ t₂ => .arrow (t₁.subst k s) (t₂.subst k s)
  | .forall_ t   => .forall_ (t.subst (k + 1) (s.shift 1 0))
\end{minted}

Now the type checker. It has five cases---three from the STLC (Var, Abs, App)
plus two new ones:

\begin{minted}{lean}
def fTypeCheck (ctx : List FTy) : FExpr → Option FTy
  -- Same as STLC:
  | .bvar n       => ctx.get? n
  | .lam ty body  =>
    match fTypeCheck (ty :: ctx) body with
    | some bodyTy => some (.arrow ty bodyTy)
    | none        => none
  | .app e₁ e₂   => do
    let funTy ← fTypeCheck ctx e₁
    let argTy ← fTypeCheck ctx e₂
    match funTy with
    | .arrow dom cod => if dom == argTy then some cod else none
    | _              => none
  -- NEW: type abstraction (Λα. e)
  | .tlam body    =>
    -- Shift all types in ctx up (they're now under one more ∀)
    let ctx' := ctx.map (·.shift 1 0)
    match fTypeCheck ctx' body with
    | some bodyTy => some (.forall_ bodyTy)
    | none        => none
  -- NEW: type application (e [τ])
  | .tapp e ty    => do
    let eTy ← fTypeCheck ctx e
    match eTy with
    | .forall_ bodyTy => some (bodyTy.subst 0 ty)  -- substitute τ for α
    | _               => none
\end{minted}

The two new rules are:

\begin{center}
\AxiomC{$\Gamma' \proves e : \tau$}
\RightLabel{\;\textsc{TAbs}}
\UnaryInfC{$\Gamma \proves (\Lam\alpha.\, e) : \forall\alpha.\, \tau$}
\DisplayProof
\qquad
\AxiomC{$\Gamma \proves e : \forall\alpha.\, \tau$}
\RightLabel{\;\textsc{TApp}}
\UnaryInfC{$\Gamma \proves e\;[\sigma] : \tau[\alpha := \sigma]$}
\DisplayProof
\end{center}

\textsc{TAbs} says: to type $\Lam\alpha.\, e$, check $e$ in a context where
all existing types are shifted (because we've introduced a new type variable).
\textsc{TApp} says: if $e$ has type $\forall\alpha.\, \tau$, then applying it
to a type $\sigma$ gives $\tau$ with $\alpha$ replaced by $\sigma$. That
substitution in the code is \texttt{bodyTy.subst 0 ty}---exactly what we built.

\begin{minted}{lean}
-- Test: the polymorphic identity  Λα. λ(x:α). x : ∀α. α → α
#eval fTypeCheck [] (.tlam (.lam (.tvar 0) (.bvar 0)))
-- Result: some (FTy.forall_ (FTy.arrow (FTy.tvar 0) (FTy.tvar 0)))
--       = some (∀α. α → α)  ✓
\end{minted}


%======================================================================
\section{System F$\omega$: Type Operators}\label{sec:fomega}
%======================================================================

System F lets terms depend on types (polymorphism). But what about types that
depend on \emph{other types}? We've already seen the intuition in our
discussion of higher-kinded types (Section~\ref{sec:hkt}): \textsf{List} is
not a type---it's a \emph{function on types}. Give it \textsf{Nat} and it
returns \textsf{List Nat}. System F$\omega$ makes this kind of type-level
computation a first-class concept.

\subsection{From Types to Kind Systems}

In the STLC, every type has the same status: \textsf{Nat}, \textsf{Bool},
$\textsf{Nat} \to \textsf{Bool}$---these are all ``types.'' But in System
F$\omega$, we need to distinguish between types that are ``complete'' (you can
have a value of that type) and type constructors that are ``waiting for an
argument.''

This distinction is formalized with \textbf{kinds}---the ``types of types'':

\begin{center}
\renewcommand{\arraystretch}{1.4}
\begin{tabular}{lll}
\textbf{Kind} & \textbf{Meaning} & \textbf{Examples} \\
\hline
$\star$ & A proper type (has values) & \textsf{Nat}, \textsf{Bool},
  $\textsf{List}\;\textsf{Nat}$ \\
$\star \to \star$ & Takes a type, returns a type & \textsf{List},
  \textsf{Option}, \textsf{IO} \\
$\star \to \star \to \star$ & Takes two types, returns a type &
  \textsf{Either}, \textsf{Map}, \textsf{Pair} \\
$(\star \to \star) \to \star$ & Takes a type constructor, returns a type &
  \textsf{Fix} (fixed point of a functor) \\
\end{tabular}
\end{center}

Kinds are to types what types are to terms. Just as a type system prevents you
from applying an integer where a function is expected ($3\;5$ is a type error),
a \emph{kind system} prevents you from applying a type where a type constructor
is expected ($\textsf{Nat}\;\textsf{Bool}$ is a \emph{kind} error---\textsf{Nat}
has kind $\star$, so it doesn't accept arguments).

\subsection{Type-Level Functions}

In the STLC and System F, types are static---they can contain type variables,
but they can't \emph{compute}. In System F$\omega$, you can define
\textbf{type-level lambda abstractions}:

\[
  \lam\alpha{:}\star.\, \textsf{Pair}\;\alpha\;\alpha
  \;\;:\;\; \star \to \star
\]

This is a type-level function: give it a type $\tau$, and it returns
$\textsf{Pair}\;\tau\;\tau$ (pairs where both components have the same type).
You can apply it: $(\lam\alpha{:}\star.\, \textsf{Pair}\;\alpha\;\alpha)\;
\textsf{Nat} \;=\; \textsf{Pair}\;\textsf{Nat}\;\textsf{Nat}$.

This is exactly beta reduction, but at the type level. The computation rule is
the same---substitution. What changes is that we allow this computation to
happen \emph{inside types}, not just inside terms.

\subsection{What System F$\omega$ Enables}

The practical payoff is the ability to abstract over type constructors, as we
saw in Section~\ref{sec:hkt}. Here is a more formal view:

In System F, you can write:
\[
  \Lam\alpha{:}\star.\, \lam(x:\alpha).\, x
  \;\;:\;\; \forall(\alpha{:}\star).\, \alpha \to \alpha
\]
The type variable $\alpha$ ranges over types of kind $\star$.

In System F$\omega$, you can also write:
\[
  \Lam(f{:}\star{\to}\star).\, \Lam(\alpha{:}\star).\, \lam(x:f\;\alpha).\, x
  \;\;:\;\; \forall(f{:}\star{\to}\star).\, \forall(\alpha{:}\star).\, f\;\alpha \to f\;\alpha
\]
Now $f$ ranges over type constructors of kind $\star \to \star$---things like
\textsf{List}, \textsf{Option}, \textsf{Tree}. Inside the term, we apply $f$
to $\alpha$ to get the concrete type $f\;\alpha$. This is exactly the
``\texttt{Functor<F>}'' pattern we discussed earlier, expressed in the formal
calculus.

\subsection{System F$\omega$ in Practice}

Haskell is the mainstream language closest to System F$\omega$. Its kind
system, while implicit in early versions, was made explicit with the
\texttt{KindSignatures} and \texttt{PolyKinds} extensions:

\begin{minted}{haskell}
-- Haskell infers: Functor :: (Type -> Type) -> Constraint
class Functor (f :: Type -> Type) where
  fmap :: (a -> b) -> f a -> f b

-- Higher-kinded type variable: t has kind (Type -> Type) -> Type -> Type
class MonadTrans (t :: (Type -> Type) -> Type -> Type) where
  lift :: Monad m => m a -> t m a
\end{minted}

The \texttt{MonadTrans} example is striking: \texttt{t} is a type constructor
that takes a type constructor as an argument and returns a type constructor.
Its kind is $(\star \to \star) \to \star \to \star$---a second-order type-level
function. This level of abstraction is impossible in System F and is the reason
Haskell can define monad transformers generically.

\begin{intuition}[Where System F$\omega$ sits]
System F$\omega$ occupies a sweet spot: it's powerful enough for generic
programming abstractions (functors, monads, monad transformers) but still
has decidable type checking and retains strong normalization. It doesn't
have dependent types, so you can't express propositions as types---but you
also don't need termination proofs or proof obligations. Most working
programmers who use advanced generics (Haskell, Scala with higher-kinded
types) are operating at roughly the System F$\omega$ level.
\end{intuition}


%======================================================================
\section{The Lambda Cube}\label{sec:cube-sec}
%======================================================================

We've now seen three extensions beyond the STLC: System F (terms depend on
types), System F$\omega$ (types depend on types), and we're about to see
dependent types (types depend on terms). Henk Barendregt organized these
extensions into a unified framework.


\subsection{The Question: What Can Depend on What?}

Every type system answers a fundamental question: \textbf{what kinds of things
can be parameterized by what other kinds of things?} In the STLC, the only
abstraction is functions: terms parameterized by terms. System F adds a second
kind: terms parameterized by types (polymorphism). But there are two other
natural possibilities we haven't explored.

\textbf{Types parameterized by types} (type operators). Consider the type
$\textsf{List}$. By itself, $\textsf{List}$ is not a type---it's a
\emph{function on types}: give it $\textsf{Nat}$ and it returns the type
$\textsf{List}\;\textsf{Nat}$. In Haskell notation, $\textsf{List}$ has
``kind'' $\star \to \star$---it takes a type and returns a type. This is
\textbf{higher-kinded polymorphism}, and it's what lets Haskell define
\texttt{Functor}, \texttt{Monad}, and other abstractions that work over
\emph{any container type}, not just lists.

\textbf{Types parameterized by terms} (dependent types). The type
$\textsf{Vec}\;n$ depends on a \emph{value} $n$---a number, not a type. This
is a function from terms to types, and it's the most powerful extension.

\subsection{Barendregt's Classification}

In 1991, Henk Barendregt (the Dutch mathematician who also proved key properties
of the untyped $\lam$-calculus and literally wrote the book on it---his 1984
monograph \emph{The Lambda Calculus: Its Syntax and Semantics} is the definitive
reference) organized these four kinds of dependency into a three-dimensional
cube. Each axis represents one ``new'' kind of dependency beyond the
STLC's terms-to-terms:

\begin{center}
\renewcommand{\arraystretch}{1.6}
\begin{tabular}{lp{7cm}l}
\textbf{Dependency} & \textbf{What it means} & \textbf{System} \\
\hline
Terms $\to$ Terms & Ordinary functions (already in STLC) & $\lam_\to$ \\
Types $\to$ Terms & Polymorphism: a term that works at many types & System F ($\lam 2$) \\
Types $\to$ Types & Type operators: a type built from other types & $\lam\underline\omega$ \\
Terms $\to$ Types & Dependent types: a type computed from a value & $\lam P$ \\
\end{tabular}
\end{center}

Each corner of the cube enables some subset of these three ``extra'' axes.
The STLC ($\lam_\to$) sits at the origin with none enabled. System F enables
one axis. The corner with all three enabled is the \textbf{Calculus of
Constructions} (CoC), introduced by Thierry Coquand and G\'erard Huet in 1988.

\subsection{Why a Cube?}

The cube structure reveals a beautiful pattern. Each axis is independent: you
can add polymorphism without dependent types, or dependent types without type
operators. The systems at each corner satisfy a clean compositional property---
the features compose without interfering. This is why the cube has $2^3 = 8$
corners, each a coherent type system.

Barendregt's insight was not that these systems existed (most had been
studied individually), but that they formed a \emph{unified family} with a
common structure. Every system in the cube has the same syntax (variables,
abstraction, application), the same computation rule (beta reduction), and
the same style of typing rules (contexts, judgments, inference rules). What
varies is only \emph{which sorts of things can appear in which positions} in
the typing rules.

This is philosophically significant: it suggests that the fundamental notion
of ``computation by substitution'' (beta reduction) is stable, and what changes
across type theories is only the question of what we \emph{allow} to be
substituted where.

\subsection{Types $\to$ Types in Practice}\label{sec:hkt}

The Types $\to$ Types axis is the subtlest to understand. Let's make it
concrete by comparing what you \emph{can} and \emph{cannot} express in
different languages.

\paragraph{What System F gives you: type application.}
In System F (and TypeScript, Java, Rust, etc.), you can \emph{apply} a type
constructor to an argument: \texttt{Array<string>}, \texttt{List Nat},
\texttt{Vec<i32>}. This is just instantiating a type variable. But the type
constructor \texttt{Array} itself is not a value you can abstract over---it's
baked into the definition.

\paragraph{What System F does \emph{not} give you: abstracting over the
constructor.} In System F, type variables range over \emph{types} like
\textsf{Nat}, \textsf{Bool}, \textsf{List Nat}---complete types you could have
a value of. But \textsf{List} by itself is not a complete type. It's a
\emph{function on types}: give it \textsf{Nat} and it returns
\textsf{List Nat}. It has ``kind'' $\star \to \star$ (a type that takes a type
and produces a type). System F has no way to write ``for all type-level
functions $F$ of kind $\star \to \star$.''

\paragraph{Why it matters: the Functor example.}
In Haskell, which has higher-kinded types, you can write \texttt{fmap} once and
have it work for \emph{any} container:

\begin{minted}{haskell}
class Functor f where                    -- f has kind Type -> Type
  fmap :: (a -> b) -> f a -> f b

-- Works for every container:
fmap (+1) [1, 2, 3]                      -- [2, 3, 4]       (f = List)
fmap (+1) (Just 5)                       -- Just 6           (f = Maybe)
fmap show (Node 1 Leaf Leaf)             -- Node "1" ...     (f = Tree)
\end{minted}

The variable \texttt{f} ranges over \emph{type constructors}: \texttt{List},
\texttt{Maybe}, \texttt{Tree}. The type \texttt{f a} means ``apply whatever
container \texttt{f} is to the element type \texttt{a}.'' The return type
\texttt{f b} means ``the same container, with transformed contents.'' The type
system tracks which container you started with and guarantees you get the same
shape back.

In TypeScript, you cannot express this. A first attempt might try:

\begin{minted}{typescript}
// NOT VALID TypeScript:
interface Functor<F<_>> {
  map<A, B>(fa: F<A>, f: (a: A) => B): F<B>;
}
\end{minted}

The \texttt{F<\_>} syntax does not exist in TypeScript. You might try a
workaround with \texttt{extends}:

\begin{minted}{typescript}
interface Functor<F extends Array<unknown> | Set<unknown>> {
  map<A, B>(fa: F, f: (a: A) => B): ???;
  //                                ^^^
  // Problem: F is already Array<unknown> — a concrete type.
  // We can't "re-open" it and change the inner type from A to B.
  // We can't express "same container, different contents."
}
\end{minted}

The problem: \texttt{Array<unknown>} is a \emph{complete type}. TypeScript has
forgotten that it was built by applying the constructor \texttt{Array} to
something. You can't say ``take whatever constructor built \texttt{F} and
re-apply it to \texttt{B}.'' The constructor is gone---consumed during type
application.

\paragraph{What it would look like with higher-kinded types.}
If TypeScript had the Types $\to$ Types axis, you could write:

\begin{minted}{typescript}
// FANTASY TypeScript with higher-kinded types:
interface Functor<F<_>> {                    // F takes one type param
  map<A, B>(fa: F<A>, f: (a: A) => B): F<B>;
}

const ArrayFunctor: Functor<Array> = {       // Array : Type -> Type
  map: (fa, f) => fa.map(f)
};
const PromiseFunctor: Functor<Promise> = {
  map: (fa, f) => fa.then(f)
};

// Generic over ANY container:
function double<F<_>>(func: Functor<F>, fa: F<number>): F<number> {
  return func.map(fa, x => x * 2);
}

double(ArrayFunctor, [1, 2, 3])       // type: Array<number>
double(PromiseFunctor, fetchCount())  // type: Promise<number>
\end{minted}

The crucial line is \verb|double<F<_>>|: a function generic over the
\emph{container itself}. The return type \texttt{F<number>} says ``whatever
container you gave me, you get the same shape back.'' Without higher-kinded
types, you need a separate \texttt{doubleArray}, \texttt{doublePromise},
\texttt{doubleSet}, and so on---exactly the ``code repetition'' problem that
polymorphism was supposed to solve, but at one level up.

\begin{intuition}[The pattern]
Each axis of the Lambda Cube removes one kind of code repetition:
\begin{enumerate}[leftmargin=2em]
\item \textbf{System F} (Types $\to$ Terms): Write one function for all
  \emph{element types}. Instead of \texttt{identityNat}, \texttt{identityBool},
  \ldots, write \texttt{identity<T>}.
\item \textbf{System F$\omega$} (Types $\to$ Types): Write one function for all
  \emph{container types}. Instead of \texttt{mapArray}, \texttt{mapPromise},
  \ldots, write \texttt{map<F>}.
\item \textbf{Dependent types} (Terms $\to$ Types): Write one type for all
  \emph{values}. Instead of \texttt{Vec3}, \texttt{Vec4}, \ldots, write
  \texttt{Vec n}.
\end{enumerate}
\end{intuition}

\subsection{Simulation Is Not the Same as Native Support}\label{sec:simulation}

A natural reaction to the Lambda Cube is: ``But my language \emph{can} express
higher-kinded types!'' For example, TypeScript's \texttt{fp-ts} library
simulates Functor using a clever encoding: string literal types serve as names
for type constructors, and TypeScript's module augmentation acts as a global
registry that maps each name back to its concrete type. The result looks
remarkably close to Haskell's \texttt{Functor}:

\begin{minted}{typescript}
// HKT.ts — a "lookup table" from string names to type constructors
interface URI2HKT<A> {}
type Kind<F extends keyof URI2HKT<any>, A> = URI2HKT<A>[F];

// Functor.ts — F is a string key, Kind resolves it
interface Functor1<F extends keyof URI2HKT<any>> {
  map: <A, B>(f: (a: A) => B, fa: Kind<F, A>) => Kind<F, B>;
}

// Option.ts — register Option in the lookup table
declare module './HKT' {
  interface URI2HKT<A> { Option: Option<A>; }
}
\end{minted}

It works. But it does not make TypeScript a System F$\omega$ language. The
Lambda Cube boundaries are not about what can be \emph{simulated}---they are
about what the type system \emph{natively understands and can reason about}.

\paragraph{The analogy with the untyped calculus.}
Consider: the untyped $\lam$-calculus can compute anything---it's Turing
complete. Church numerals simulate natural numbers, Church booleans simulate
booleans, the $\mathbf{Y}$ combinator simulates recursion. Does that make the
untyped calculus ``a system with natural numbers and booleans and recursion''?
In a sense yes---the encodings work. But in the important sense no: the system
doesn't \emph{know} those things are numbers. It can't check that you didn't
accidentally add a boolean to a list. The STLC doesn't add any computational
power---it adds \emph{knowledge}. The type checker understands what is a number
and what is a function, and it can verify that you use each correctly.

The same principle applies at every level of the cube:

\begin{center}
\renewcommand{\arraystretch}{1.4}
\begin{tabular}{p{4cm}p{9cm}}
\textbf{Level} & \textbf{What ``native support'' means} \\
\hline
System F & The compiler understands that a type variable ranges over
  types. It can check universality, infer instances, and optimize based on
  parametricity. \\
System F$\omega$ & The compiler understands that a type variable ranges over
  \emph{type constructors}. It can check kind correctness ($\star \to \star$
  vs.\ $\star$), report errors in terms of type constructors, and derive
  instances. \\
Dependent types & The compiler understands that types can contain
  \emph{values}. It can evaluate expressions during type checking, compare
  types up to computation, and verify proofs. \\
\end{tabular}
\end{center}

In the \texttt{fp-ts} encoding, TypeScript thinks \texttt{F} is a
\emph{string}. It has no concept of kinds. It doesn't know the encoding
represents a type constructor---it's just checking that a string literal is a
valid key in an interface. The \emph{programmer} knows the abstraction is
there; the \emph{compiler} doesn't. This means: error messages talk about
string keys rather than type constructors; the compiler can't verify functor
laws or kind correctness; and optimizations based on parametricity are
impossible.

\begin{intuition}[The OOP analogy]
You can simulate object-oriented programming in C with structs containing
function pointers. The Linux kernel does this extensively. Does that make C
an object-oriented language? It can \emph{express} OO patterns, but the
language doesn't \emph{understand} them---there's no inheritance checking, no
virtual dispatch guarantees, no access control enforced by the compiler. The
patterns work because the programmer maintains the discipline, not because
the type system enforces it.

Similarly, you can simulate higher-kinded types in TypeScript, or dependent
types in Java (with enough generics abuse). The simulation works---but the
compiler is enforcing a \emph{different, weaker} set of guarantees than a
system that natively supports the feature. The Lambda Cube describes what the
type system \emph{understands as first-class concepts}: what it can reason
about, check, and enforce---not what clever encodings can approximate.
\end{intuition}


%======================================================================
\section{Dependent Types: Types That Compute}\label{sec:dependent-sec}
%======================================================================

\subsection{Martin-L\"of's Type Theory}\label{sec:martinlof}

Per Martin-L\"of, a Swedish logician and philosopher, took the BHK
interpretation literally and built a formal system around it. His
\textbf{intuitionistic type theory} (1971, with major revisions in 1973 and
1984) made the connection between proofs and programs fully precise: each
logical connective becomes a type constructor, and each proof rule becomes a
typing rule.

In Martin-L\"of's system, the type $\Sigma(x{:}A).\, B(x)$ (the dependent pair
or ``existential'') represents ``there exists an $x : A$ such that $B(x)$.'' A
term of this type is a \emph{pair} $(a, b)$ where $a : A$ and
$b : B(a)$---literally a witness $a$ together with proof $b$ that the witness
satisfies the property. You cannot claim existence without producing the thing
that exists.

The type $\Pi(x{:}A).\, B(x)$ (the dependent function or ``universal'')
represents ``for all $x : A$, $B(x)$.'' A term of this type is a function
that, given any $a : A$, produces a proof of $B(a)$. A ``for all'' claim must
come with a uniform method that works for every input.

This is the BHK interpretation, formalized as a type system:

\begin{center}
\renewcommand{\arraystretch}{1.4}
\begin{tabular}{lcl}
\textbf{Logic} & & \textbf{Martin-L\"of type theory} \\
\hline
$A \land B$ & $\longleftrightarrow$ & $A \times B$ (product type) \\
$A \lor B$ & $\longleftrightarrow$ & $A + B$ (sum type) \\
$A \Rightarrow B$ & $\longleftrightarrow$ & $A \to B$ (function type) \\
$\exists x{:}A.\, B(x)$ & $\longleftrightarrow$ & $\Sigma(x{:}A).\, B(x)$ (dependent pair) \\
$\forall x{:}A.\, B(x)$ & $\longleftrightarrow$ & $\Pi(x{:}A).\, B(x)$ (dependent function) \\
$\bot$ (false) & $\longleftrightarrow$ & Empty type (no constructors) \\
$\lnot A$ & $\longleftrightarrow$ & $A \to \bot$ (function to empty type) \\
\end{tabular}
\end{center}

The last row is illuminating: ``not $A$'' is defined as a function from $A$ to
the empty type. If you could produce a value of type $A$, this function would
give you a value of the empty type---which is impossible, since the empty type
has no values. So having such a function means $A$ itself must be uninhabited
(unprovable). Negation is not a primitive---it's \emph{derived} from functions
and the empty type.

\subsection{From Philosophy to Software}

Martin-L\"of's type theory remained a topic in mathematical logic for over a
decade. The transition to software came through two parallel efforts.

\textbf{Coquand and Huet (1985--1988).} Thierry Coquand, a French logician
(and student of Girard), and G\'erard Huet at INRIA developed the
\textbf{Calculus of Constructions} (CoC)---a type theory that sat at the
``top corner'' of the Lambda Cube, combining polymorphism, type operators, and
dependent types in a single unified system. Their insight was that
Martin-L\"of's ideas could be organized into a clean, minimalist calculus with
just a few rules. The CoC became the theoretical foundation of the \textbf{Coq}
proof assistant (1989, named as a pun on Coquand's name and the French word
for ``rooster'').

\textbf{The Automath and NuPRL traditions.} Independently, de Bruijn's Automath
project (1968, the same de Bruijn who invented de Bruijn indices) had already
explored dependent types for formalizing mathematics. Robert Constable's NuPRL
system at Cornell (1986) implemented Martin-L\"of's type theory directly. These
systems demonstrated that dependent types could be made practical, not just
theoretical.

Lean 4 (2021, Leonardo de Moura at Microsoft Research) continues this lineage,
implementing the \textbf{Calculus of Inductive Constructions} (CIC)---the CoC
extended with inductive types and a universe hierarchy.

\subsection{How Dependent Types Work}

In the STLC, types and terms live in separate worlds---a type
like $\textsf{Nat} \to \textsf{Bool}$ mentions other types but never mentions
specific values. With dependent types, the wall comes down.

The classic example is a vector with its length in the type:
\[
  \textsf{Vec} : \textsf{Type} \to \textsf{Nat} \to \textsf{Type}
\]
$\textsf{Vec}\;\textsf{Nat}\;3$ is the type of lists of exactly three natural
numbers. $\textsf{Vec}\;\textsf{Bool}\;0$ is the type of empty boolean lists.
The length $3$ or $0$ is a \emph{value} appearing inside a \emph{type}.

This lets you express invariants that simpler type systems can't:
\begin{minted}{lean}
-- This type GUARANTEES the output has the same length as the inputs
def zipWith (f : α → β → γ) : Vec α n → Vec β n → Vec γ n
\end{minted}

If you try to zip vectors of different lengths, the type checker rejects it at
compile time---not at runtime. The \texttt{n} in the type signature is a value
(a natural number), not just a type variable.

\subsection{The Unification: Types Are Terms}

Here is the conceptual leap. In the STLC and System F, there are two
separate ``worlds'': the world of terms (which compute) and the world of
types (which classify). These worlds interact through typing judgments
($e : \tau$), but they are fundamentally different kinds of objects.

Dependent types \emph{erase this distinction}. A type is just a term that
happens to classify other terms. $\textsf{Nat}$ is a term (of type
$\textsf{Type}$). $\textsf{Vec}\;\textsf{Nat}\;3$ is a term (built by
applying the function $\textsf{Vec}$ to the arguments $\textsf{Nat}$ and $3$).
Even the function type $A \to B$ is a term (a special case of $\Pi(x{:}A).\,B$
where $B$ doesn't mention $x$).

This unification means the type checker must be able to \emph{evaluate terms}
during type checking. To decide if $\textsf{Vec}\;\textsf{Nat}\;(2+1)$ and
$\textsf{Vec}\;\textsf{Nat}\;3$ are the same type, the type checker must
compute $2+1$ and see that it equals $3$. The type checker is an evaluator.
This is why our dependent type checker implementation includes \texttt{whnf}
(evaluation) and \texttt{beq} (comparison after evaluation)---they are not
optional features but essential to the core algorithm.

\begin{intuition}[The Curry--Howard payoff]
With dependent types, the Curry--Howard correspondence reaches its full
potential. A type like $\forall (n : \textsf{Nat}).\, n + 0 = n$ is
simultaneously:
\begin{enumerate}[leftmargin=2em]
\item A \textbf{type} (you can declare variables of this type)
\item A \textbf{proposition} (``for all $n$, $n + 0 = n$'')
\end{enumerate}
A term inhabiting this type is simultaneously a \textbf{program} and a
\textbf{proof}. This is exactly how Lean works: when you write
\texttt{theorem add\_zero (n : Nat) : n + 0 = n}, you are defining a term
of a dependent type. The proof is a program; the type checker verifies it.
\end{intuition}

\subsection{The Price of Power: Undecidability and Type : Type}

Dependent types come with subtleties that simpler systems avoid.

\textbf{Type checking requires evaluation.} Since types can contain arbitrary
computation, type checking in a dependently typed system is only decidable if
all programs terminate. This is why Lean requires termination proofs for
recursive functions (and why you fought with \texttt{termination\_by} when
implementing the interpreter!). If a non-terminating computation could appear
in a type, the type checker might loop forever trying to normalize it.

\textbf{The Type : Type paradox.} If \texttt{Type} has type \texttt{Type}, you
get a variant of Russell's paradox. Girard discovered this in 1972 (``Girard's
paradox''): in a system with \texttt{Type : Type}, you can construct a term of
any type, meaning every proposition is ``provable''---the system is logically
inconsistent. The solution is a \textbf{universe hierarchy}: $\textsf{Type}_0 :
\textsf{Type}_1 : \textsf{Type}_2 : \ldots$, where each level can only refer
to levels below it. Lean implements this with its \texttt{Sort} hierarchy. Our
toy implementation uses a single \texttt{univ} for simplicity, which is why it
is not logically consistent---but the core algorithm is the same.

\subsection{Implementing a Dependent Type Checker}

This is where the magic happens. In a dependently typed system, \textbf{types
and terms are the same thing}. There is no separate grammar for types---a type
is just an expression that happens to classify other expressions. This
simplifies the syntax dramatically:

\begin{minted}{lean}
-- In dependent types, terms and types are unified.
-- "Type" is itself a term (a universe).
inductive DTerm where
  | bvar   : Nat → DTerm               -- variable (de Bruijn)
  | app    : DTerm → DTerm → DTerm     -- application: e₁ e₂
  | lam    : DTerm → DTerm → DTerm     -- λ(x:A). e  (A is a term!)
  | pi     : DTerm → DTerm → DTerm     -- Π(x:A). B  (dependent function type)
  | univ   : DTerm                      -- Type (the universe)
  deriving Repr, DecidableEq, Inhabited
\end{minted}

Notice: there is no separate \texttt{Ty} type. The \texttt{pi} constructor is
the dependent function type $\Pi(x{:}A).\, B$, where $B$ can mention $x$. When
$B$ does \emph{not} mention $x$, this is just the ordinary arrow $A \to B$. And
\texttt{univ} is the type of types: \texttt{Type} itself.

The type checker is more complex than before because it must \textbf{evaluate
terms during type checking}. When checking whether two types are equal, we
can't just compare syntax---we need to beta-reduce them first. For example,
$(\lam x.\, x)\;\textsf{Nat}$ and $\textsf{Nat}$ are different expressions
but the same type.

\begin{minted}{lean}
-- Shifting and substitution (same pattern as before)
def DTerm.shift (d : Int) (c : Nat) : DTerm → DTerm
  | .bvar n    => if n ≥ c then .bvar (Int.toNat (n + d)) else .bvar n
  | .app e₁ e₂ => .app (e₁.shift d c) (e₂.shift d c)
  | .lam ty e  => .lam (ty.shift d c) (e.shift d (c + 1))
  | .pi  ty e  => .pi  (ty.shift d c) (e.shift d (c + 1))
  | .univ      => .univ

def DTerm.subst (k : Nat) (s : DTerm) : DTerm → DTerm
  | .bvar n    => if n == k then s
                   else if n > k then .bvar (n - 1) else .bvar n
  | .app e₁ e₂ => .app (e₁.subst k s) (e₂.subst k s)
  | .lam ty e  => .lam (ty.subst k s) (e.subst (k+1) (s.shift 1 0))
  | .pi  ty e  => .pi  (ty.subst k s) (e.subst (k+1) (s.shift 1 0))
  | .univ      => .univ

-- Beta reduction (needed during type checking!)
partial def DTerm.whnf : DTerm → DTerm
  | .app (.lam _ body) arg => (body.subst 0 arg).whnf
  | e                      => e

-- Definitional equality: are two terms the same after reduction?
partial def DTerm.beq (a b : DTerm) : Bool :=
  match a.whnf, b.whnf with
  | .bvar n₁,    .bvar n₂    => n₁ == n₂
  | .app f₁ a₁,  .app f₂ a₂  => f₁.beq f₂ && a₁.beq a₂
  | .lam t₁ e₁,  .lam t₂ e₂  => t₁.beq t₂ && e₁.beq e₂
  | .pi  t₁ e₁,  .pi  t₂ e₂  => t₁.beq t₂ && e₁.beq e₂
  | .univ,        .univ        => true
  | _,            _            => false
\end{minted}

The \texttt{whnf} function (weak head normal form) reduces applications until
the top-level structure is exposed. The \texttt{beq} function compares two
terms after reduction---this is \textbf{definitional equality}, the core of
dependent type checking. Two types are considered equal if they reduce to the
same thing.

Now the type checker itself. The context maps variable indices to their types
(which are themselves \texttt{DTerm}s):

\begin{minted}{lean}
partial def dTypeCheck (ctx : List DTerm) : DTerm → Option DTerm
  -- Variable: look up its type in the context
  | .bvar n    => ctx.get? n
  -- Universe: Type has type Type (simplified; real systems use a hierarchy)
  | .univ      => some .univ
  -- Pi type: Π(x:A).B is valid if A : Type and B : Type (with x:A in scope)
  | .pi a b    => do
    let aTy ← dTypeCheck ctx a
    guard (aTy.whnf matches .univ)          -- A must be a type
    let bTy ← dTypeCheck (a :: ctx) b
    guard (bTy.whnf matches .univ)          -- B must be a type
    some .univ                               -- then Π(x:A).B : Type
  -- Lambda: λ(x:A).e has type Π(x:A).B if e:B with x:A in scope
  | .lam a body => do
    let aTy ← dTypeCheck ctx a
    guard (aTy.whnf matches .univ)
    let bodyTy ← dTypeCheck (a :: ctx) body
    some (.pi a bodyTy)
  -- Application: (e₁ e₂) — e₁ must have a Π type, e₂ must match the domain
  | .app e₁ e₂ => do
    let funTy ← dTypeCheck ctx e₁
    match funTy.whnf with
    | .pi dom cod =>
      let argTy ← dTypeCheck ctx e₂
      guard (dom.beq argTy)                  -- argument type must match domain
      some (cod.subst 0 e₂)                 -- CRUCIAL: substitute the actual
                                             -- argument into the codomain!
    | _ => none
\end{minted}

\begin{intuition}[The magic line]
The most important line in the entire type checker is in the application case:
\texttt{cod.subst 0 e₂}. This is where dependent types earn their name. In
the STLC, applying $f : A \to B$ to $x : A$ gives type $B$---the return type
is fixed. But in a dependent system, applying $f : \Pi(x{:}A).\, B$ to $a : A$
gives type $B[x := a]$---the return type \emph{depends on the argument}. This
single substitution is what makes dependent types able to express ``the output
list has the same length as the input'' or ``this function returns a proof that
its output is sorted.''
\end{intuition}

Let's test it. We can encode the STLC's identity function and verify it:

\begin{minted}{lean}
-- λ(A : Type). λ(x : A). x  :  Π(A : Type). Π(_ : A). A
-- This is the dependently-typed polymorphic identity!
def polyId : DTerm :=
  .lam .univ (.lam (.bvar 0) (.bvar 0))

#eval dTypeCheck [] polyId
-- Result: some (DTerm.pi DTerm.univ (DTerm.pi (DTerm.bvar 0) (DTerm.bvar 1)))
-- Which reads: Π(A : Type). Π(x : A). A    ✓
\end{minted}

\paragraph{Why $\lambda$ and $\Pi$ look so similar (but aren't).}
The term and its type look almost identical, which can be confusing:
\begin{center}
\renewcommand{\arraystretch}{1.3}
\begin{tabular}{lll}
\textbf{Term} (the function) & $\lam(A : \textsf{Type}).\, \lam(x : A).\, x$
  & ``given $A$ and $x$, return $x$'' \\
\textbf{Type} (the signature) & $\Pi(A : \textsf{Type}).\, \Pi(\_ : A).\, A$
  & ``takes a type $A$ and an $A$, returns an $A$'' \\
\end{tabular}
\end{center}

$\lambda$ \emph{builds} a function. $\Pi$ \emph{describes} what kind of
function it is. Every $\lambda$ has a $\Pi$ as its type, the same way every
function in a programming language has a type signature. In everyday code, the
distinction is obvious:
\begin{minted}{typescript}
// The FUNCTION (λ):    (x: number) => x + 1
// The TYPE (Π):        (x: number) => number
\end{minted}
They look similar because they both mention the parameter name and type. But
one is the \emph{code} (what to compute) and the other is the \emph{contract}
(what goes in and what comes out).

\paragraph{Why the de Bruijn indices differ.}
In the \emph{term} \texttt{.lam .univ (.lam (.bvar 0) (.bvar 0))}, both uses
of \texttt{.bvar 0} mean ``the nearest enclosing binder'':

\begin{minted}{text}
  λ(A : Type). λ(x : A). x
  ↑ binder 1    ↑ binder 0  ↑ bvar 0 = nearest binder = x  ✓
                ↑ the annotation A here is (.bvar 0),
                  because at this point only the outer λ is in scope,
                  so A (the outer λ's variable) is index 0
\end{minted}

But in the \emph{type} \texttt{DTerm.pi .univ (DTerm.pi (.bvar 0) (.bvar 1))},
the final \texttt{.bvar 1} reaches past the inner binder:

\begin{minted}{text}
  Π(A : Type). Π(_ : A).  A
  ↑ binder 1    ↑ binder 0  ↑ bvar 1 = skip inner binder,
                               grab outer binder = A  ✓
                ↑ (.bvar 0) = A, because only the outer Π is in scope here
\end{minted}

The term has \texttt{.bvar 0} at the end because its final variable is
$x$ (the \emph{inner} $\lambda$'s variable---the nearest binder). The type has
\texttt{.bvar 1} at the end because its return type is $A$ (the \emph{outer}
$\Pi$'s variable), and from inside two nested $\Pi$ binders, the outer one is
at index 1.


%======================================================================
\section{Type Checking Is Proof Checking}\label{sec:proofcheck-sec}
%======================================================================


Here is the punchline. In a dependently typed system, propositions are types
and proofs are terms. So \textbf{checking a proof is just type checking a
term}. There is no separate proof-verification engine---the type checker
\emph{is} the proof verifier. Let's see this in action.

We can encode simple logical connectives as types. The proposition
``$A$ implies $B$'' is the function type $A \to B$ (a non-dependent Pi type).
The proposition ``for all $x : A$, $P(x)$'' is $\Pi(x:A).\, P\;x$ (a
dependent Pi type). A \emph{proof} of these propositions is a term that
inhabits the type.

\begin{minted}{lean}
-- Proposition: A → A  ("A implies A")
-- Proof: λ(A : Type). λ(a : A). a  (the identity function!)
-- Type checker confirms: this term has type Π(A:Type). A → A  ✓
-- Therefore: the proposition "A implies A" is proved.

-- Proposition: A → B → A  ("A implies (B implies A)")
-- Proof: λ(A:Type). λ(B:Type). λ(a:A). λ(b:B). a
def kProof : DTerm :=
  .lam .univ (.lam .univ (.lam (.bvar 1) (.lam (.bvar 1) (.bvar 1))))

-- The type checker returns: Π(A:Type). Π(B:Type). Π(_:A). Π(_:B). A
-- Which IS the proposition A → B → A.
-- The fact that type checking succeeded IS the proof verification.
#eval dTypeCheck [] kProof
\end{minted}

The key insight: when the type checker accepts a term $e$ at type $\tau$, it
has verified that $e$ is a valid proof of the proposition $\tau$. When it
rejects a term, it means the ``proof'' is invalid. There is nothing else to
check---\textbf{type safety is logical consistency}.

This is exactly what happens inside Lean when you write:

\begin{minted}{lean}
-- Lean's theorem keyword is just a "def" checked by this same machinery
theorem identity (A : Prop) (a : A) : A := a

-- Lean's type checker verifies:
--   a : A       (variable lookup)
--   λa. a : A → A  (Abs rule)
--   etc.
-- If it type-checks, the proof is valid. Period.
\end{minted}

The entire 4539-line Lean kernel that verifies every proof in Lean's
\texttt{mathlib} library works on essentially this principle: terms are proofs,
types are propositions, and type checking is proof verification. What we've
built above is a toy version of the same idea---simplified, but the core
mechanism is identical.

\section{Parsing Lambda Calculus}\label{sec:parser}
%======================================================================

\subsection{What Parsing Is}

Every programming language exists in two forms. The \textbf{concrete syntax}
is what humans write: \verb|\x. x y|. The \textbf{abstract syntax} is what
the machine manipulates: \texttt{lam "x" (app (var "x") (var "y"))}. A
\textbf{parser} bridges these: it takes a string of characters and produces
a tree.

Parsing is one of the oldest problems in computer science. The theory was
developed in the 1950s--1970s alongside formal language theory, and it splits
into two parts:

\textbf{Lexing} (tokenization): splitting a character stream into tokens.
The string \verb*|\x y. x y| becomes tokens
$[\texttt{LAMBDA}, \texttt{ID}(x), \texttt{ID}(y), \texttt{DOT},
\texttt{ID}(x), \texttt{ID}(y)]$. Whitespace is consumed; meaningful chunks
are identified. For our simple grammar, we'll do this inline rather than as a
separate pass.

\textbf{Parsing} (syntax analysis): organizing tokens into a tree according to
a grammar. The token stream above becomes
$\texttt{lam}(\text{``x''}, \texttt{lam}(\text{``y''},
\texttt{app}(\texttt{var}(\text{``x''}), \texttt{var}(\text{``y''}))))$.

\paragraph{Context-free grammars.}
A grammar specifies which strings are valid programs. The standard formalism is
a \textbf{context-free grammar} (CFG), defined by Noam Chomsky in 1956
(originally for natural languages, but quickly adopted by computer scientists).
A CFG consists of:

\begin{enumerate}[leftmargin=2em]
\item \textbf{Terminals}: literal characters or tokens (\verb|\|, \texttt{.},
  identifiers)
\item \textbf{Nonterminals}: abstract categories (\textit{expr}, \textit{app},
  \textit{atom})
\item \textbf{Production rules}: how nonterminals expand into sequences of
  terminals and nonterminals
\item A \textbf{start symbol}: the nonterminal that represents a complete
  program
\end{enumerate}

Two major approaches exist for building a parser from a grammar:
\textbf{parser generators} (tools like yacc, ANTLR, or Menhir that read a
grammar file and output parser code) and \textbf{hand-written parsers}
(recursive descent or parser combinators). We'll use the hand-written approach
because it's more instructive---you see every decision the parser makes.

\begin{exercisebox}[Grammar design]
\stepcounter{exercise}\textbf{\theexercise.}\; What syntax would feel natural for writing $\lam$-terms? Think about: how do you write a lambda? How do you group things? What does application look like? Write out a grammar (like the BNF from Section~\ref{sec:syntax}) for your syntax.
\end{exercisebox}

\subsection{The Grammar}

Here is one clean design:

\begin{minted}{text}
expr   ::= '\' ident+ '.' expr    -- lambda: \x y z. body
         | app
app    ::= atom+                  -- one or more atoms, left-assoc
atom   ::= '(' expr ')'           -- grouped expression
         | ident                   -- variable name
ident  ::= [a-zA-Z_][a-zA-Z0-9_']*
\end{minted}

Multi-argument lambdas (\verb|\x y. e| desugars to nested $\lam$'s),
left-associative application, and parentheses for grouping. This grammar
is simple enough to parse without any special tools---no need for yacc, ANTLR,
or lex. Two approaches work cleanly: recursive descent and parser combinators.

\subsection{Approach 1: Recursive Descent}

The simplest technique is \textbf{recursive descent}: each grammar rule becomes
a function, and the functions call each other recursively. The structure of
the code mirrors the structure of the grammar.

\begin{exercisebox}[Parser design]
\stepcounter{exercise}\textbf{\theexercise.}\; The \texttt{parseExpr} function must decide: is this a lambda or an application? What character tells you which one? The \texttt{parseApp} function must handle left-associative application: \texttt{f x y} means \texttt{(f x) y}. How would you accumulate applications left-to-right?
\end{exercisebox}

Think about \texttt{parseExpr} first. You're looking at the input and need to
decide what you're about to parse. If the next character is \verb|\|, you know
it's a lambda. Otherwise, it must be an application (a sequence of atoms).
That's it---just one character of lookahead tells you which branch to take.

Now think about \texttt{parseApp}. Application syntax is just \emph{juxtaposition}:
\texttt{f x y} means ``apply $f$ to $x$, then apply the result to $y$.'' This
is left-associative: $f\;x\;y = (f\;x)\;y$. So you parse the first atom, then
keep parsing more atoms and wrapping them in left-associative \texttt{app} nodes.
If you start with $f$, then read $x$, you have $\texttt{app}(f, x)$. Then read $y$:
$\texttt{app}(\texttt{app}(f, x), y)$. Stop when there are no more atoms.

Finally, \texttt{parseAtom} handles the two base cases: either a parenthesized
expression (start with \texttt{(}, parse recursively, expect \texttt{)}) or a
bare identifier.

\begin{minted}{lean}
-- Dispatch: lambda or application?
partial def parseExpr (s : PState) : ParseResult Term :=
  match peek s with
  | some '\\' => parseLambda s     -- starts with \: must be a lambda
  | _         => parseApp s        -- otherwise: try application

-- Lambda: consume \, collect parameter names, consume ., parse body
-- Fold right: [x, y, z] → lam x (lam y (lam z body))
partial def parseLambda (s : PState) : ParseResult Term := ...

-- Application: parse atoms left-to-right, accumulate with app
partial def parseApp (s : PState) : ParseResult Term :=
  match parseAtom s with
  | some (first, s) => go s first   -- got one atom, try for more
  | none => none
  where go s acc := match parseAtom s with
    | some (next, s') => go s' (.app acc next)  -- left-associative!
    | none => some (acc, s)

-- Atom: parenthesized expression or identifier
partial def parseAtom (s : PState) : ParseResult Term := ...
\end{minted}

This works well for our simple grammar. But notice the manual threading of
\texttt{PState} through every function---you pass it in, get an updated state
back, pass that to the next function. If parsing fails, you return
\texttt{none} and the caller must handle it. This ``plumbing'' obscures the
structure. As grammars grow more complex, the plumbing dominates the code.

\subsection{Approach 2: Parser Combinators (The Parsec Idea)}

\textbf{Parser combinators} solve the plumbing problem by treating parsers as
\emph{values}---first-class objects that can be combined with operators, just
like numbers combine with $+$ and $\times$. The idea was pioneered by Philip
Wadler (1985), developed by Graham Hutton and Erik Meijer (1996), and brought
to practical maturity by Daan Leijen's \textbf{Parsec} library for Haskell
(2001).

\subsubsection{The Core Abstraction}

A parser is a function: it takes an input string and a position, and either
succeeds (returning a result and the new position) or fails:

\begin{minted}{lean}
-- A parser that produces values of type α
def Parser (α : Type) := String → Nat → Option (α × Nat)
-- Input string → position → maybe (result, new position)
\end{minted}

This is a simple type, but it encodes a powerful idea: a parser is a
\emph{computation with two effects}---it consumes input (state) and it can
fail (partiality). The entire Parsec approach comes from combining these
computations with a small set of operators.

\subsubsection{Primitive Combinators}

We start with combinators that do almost nothing by themselves:

\begin{minted}{lean}
-- `pure a`: always succeeds, produces `a`, consumes nothing
def pure (a : α) : Parser α := fun s pos => some (a, pos)

-- `fail`: always fails
def fail : Parser α := fun _ _ => none

-- `char c`: succeeds if the next character is `c`
def char (c : Char) : Parser Char := fun s pos =>
  if s.get ⟨pos⟩ == c then some (c, pos + 1) else none

-- `satisfy p`: succeeds if the next character satisfies predicate `p`
def satisfy (p : Char → Bool) : Parser Char := fun s pos =>
  let c := s.get ⟨pos⟩
  if p c then some (c, pos + 1) else none
\end{minted}

Each one is tiny. The power comes from \emph{composition}.

\subsubsection{Sequential Composition: Bind}

The most important combinator is \texttt{bind}: run one parser, feed its
result to a function that returns another parser, then run that:

\begin{minted}{lean}
-- `p >>= f`: run p, feed its result to f, run the resulting parser
def bind (p : Parser α) (f : α → Parser β) : Parser β := fun s pos =>
  match p s pos with
  | some (a, pos') => f a s pos'    -- p succeeded: continue with f
  | none           => none          -- p failed: propagate failure
\end{minted}

This is \emph{monadic sequencing}. All the state-threading plumbing we did
manually in recursive descent is now hidden inside \texttt{bind}. Each parser
receives the updated position automatically.

\subsubsection{Choice: Or-Else}

\begin{minted}{lean}
-- `p <|> q`: try p; if it fails, try q (backtracking)
def orElse (p q : Parser α) : Parser α := fun s pos =>
  match p s pos with
  | some r => some r        -- p succeeded
  | none   => q s pos       -- p failed: try q at the same position
\end{minted}

This is \textbf{backtracking}: if $p$ fails, we rewind to the original
position and try $q$. The parser explores alternatives without the programmer
manually managing the position.

\subsubsection{Repetition}

\begin{minted}{lean}
-- `many p`: apply p zero or more times, collect results
partial def many (p : Parser α) : Parser (List α) := fun s pos =>
  match p s pos with
  | none         => some ([], pos)
  | some (a, pos') =>
    match many p s pos' with
    | some (rest, pos'') => some (a :: rest, pos'')
    | none               => some ([a], pos')

-- `many1 p`: one or more (fails if p doesn't match at least once)
def many1 (p : Parser α) : Parser (List α) :=
  p >>= fun first =>
  many p >>= fun rest =>
  pure (first :: rest)
\end{minted}

\subsubsection{Whitespace and Tokens}

Real parsers need to skip whitespace between tokens:

\begin{minted}{lean}
-- Skip whitespace
def spaces : Parser Unit :=
  many (satisfy Char.isWhitespace) >>= fun _ => pure ()

-- Parse p, then skip trailing whitespace
def token (p : Parser α) : Parser α :=
  p >>= fun a =>
  spaces >>= fun _ =>
  pure a

-- Parse a specific character, skip trailing whitespace
def symbol (c : Char) : Parser Char := token (char c)
\end{minted}

The \texttt{token} combinator is a small but important pattern: it wraps
any parser to consume trailing whitespace, so the rest of your grammar never
mentions whitespace explicitly.

\subsubsection{Parsers as Functor, Applicative, and Monad}

In Haskell (and real Parsec implementations), parsers aren't just used with
bare functions like \texttt{bind} and \texttt{pure}. They implement the
\textbf{typeclass hierarchy}: Functor, Applicative, and Monad. This unlocks
powerful syntax and generic combinators.

\begin{minted}{lean}
-- Functor: transform the result of a parser
instance : Functor Parser where
  map f p := fun s pos =>
    match p s pos with
    | some (a, pos') => some (f a, pos')
    | none           => none

-- Applicative: run two parsers, combine their results
instance : Applicative Parser where
  pure a := fun s pos => some (a, pos)
  seq pf pa := fun s pos =>
    match pf s pos with
    | some (f, pos') =>
      match pa () s pos' with
      | some (a, pos'') => some (f a, pos'')
      | none            => none
    | none => none

-- Monad: sequence parsers where the second depends on the first's result
instance : Monad Parser where
  bind p f := fun s pos =>
    match p s pos with
    | some (a, pos') => f a s pos'
    | none           => none

-- Alternative: choice between parsers
instance : Alternative Parser where
  failure := fun _ _ => none
  orElse p q := fun s pos =>
    match p s pos with
    | some r => some r
    | none   => q () s pos
\end{minted}

Why does the typeclass hierarchy matter? Because each level gives you
\emph{different combinators for free}:

\textbf{Functor} gives you \texttt{<\$>} (map): transform a parser's result
without changing what it consumes. If \texttt{digit} parses a character,
\texttt{Char.toNat <\$> digit} parses a character and converts it to a number.

\textbf{Applicative} gives you \texttt{<*>} (apply): run two parsers in
sequence and combine their results. This is enough for parsers that don't need
to inspect intermediate results:
\begin{minted}{lean}
-- Parse a pair: (expr, expr)
def pairP : Parser (Term × Term) :=
  (fun _ a _ b _ => (a, b))          -- how to combine results
    <$> symbol '(' <*> exprP         -- ( expr
    <*> symbol ',' <*> exprP         -- , expr
    <*> symbol ')'                   -- )
\end{minted}

\textbf{Monad} gives you \texttt{>>=} (bind) and, crucially, \texttt{do}
notation. This is the most common style in practice because it reads like
imperative code:

\begin{minted}{lean}
-- The lambda parser in do-notation:
def lamP : Parser Term := do
  let _ ← symbol '\\'
  let params ← many1 ident
  let _ ← symbol '.'
  let body ← exprP
  pure (params.foldr Term.lam body)

-- Compare to the >>= version:
def lamP' : Parser Term :=
  symbol '\\' >>= fun _ =>
  many1 ident >>= fun params =>
  symbol '.'  >>= fun _ =>
  exprP       >>= fun body =>
  pure (params.foldr Term.lam body)
\end{minted}

The two are identical---\texttt{do} notation is syntactic sugar for chains
of \texttt{>>=}. But the \texttt{do} version is more readable, especially for
longer parsers. This is what real Parsec code looks like.

\textbf{Applicative vs.\ Monad.} Why have both? Applicative parsing is
\emph{less powerful}---the structure of the parse is fixed before any input is
consumed. Monadic parsing lets the second parser \emph{depend on the result}
of the first (e.g., ``if we parsed a keyword \texttt{let}, now parse a binding;
if we parsed \texttt{if}, now parse a condition''). But applicative parsing is
easier to analyze statically, which some parser libraries exploit for
optimization. Most practical parsers use monadic combinators because the
expressiveness is essential.

\begin{intuition}[Parsers are monads]
Look at the types: \texttt{Parser α} wraps a computation that might fail.
\texttt{pure} lifts a value into the parser. \texttt{bind} sequences two
parsers. This is exactly the monad pattern! Parsec's insight was that monadic
composition gives you the right abstraction for parsing: sequencing
(\texttt{>>=}), choice (\texttt{<|>}), and failure propagation---all handled
automatically.

If you've used JavaScript Promises, you've used the same pattern:
\texttt{fetch(url).then(parse).then(validate)} sequences computations that
might fail, just like \texttt{parseExpr >>= parseDot >>= parseBody}.
\end{intuition}

\subsubsection{Building Our Lambda Parser}

With these building blocks, the parser mirrors the grammar almost directly.
Now that we have \texttt{do} notation, let's use it:

\begin{minted}{lean}
-- Parse an identifier: one or more alphanumeric/underscore characters
def ident : Parser String := token do
  let chars ← many1 (satisfy fun c => c.isAlphanum || c == '_' || c == '\'')
  pure (String.mk chars)

-- Parse a keyword (an identifier matching a specific string)
def keyword (s : String) : Parser Unit := do
  let name ← ident
  if name == s then pure () else failure

-- Parse a lambda: '\' ident+ '.' expr
partial def lamP : Parser Term := do
  let _ ← symbol '\\'
  let params ← many1 ident       -- collect parameter names
  let _ ← symbol '.'
  let body ← exprP               -- parse the body
  pure (params.foldr Term.lam body) -- desugar: [x,y] → lam x (lam y body)

-- Parse an application: atom+, left-associative
partial def appP : Parser Term := do
  let atoms ← many1 atomP
  pure (atoms.tail.foldl Term.app atoms.head)

-- Parse an atom: parenthesized expression or identifier
partial def atomP : Parser Term :=
  (do let _ ← symbol '('; let e ← exprP; let _ ← symbol ')'; pure e)
  <|> (do let name ← ident; pure (Term.var name))

-- Parse an expression: lambda or application
partial def exprP : Parser Term := lamP <|> appP
\end{minted}

Compare this to the recursive descent version. The manual state threading is
gone---\texttt{bind} handles it. Error propagation is gone---\texttt{Option}
handles it. What remains is just the grammar, expressed as composition of small
parsers. The code \emph{reads like} the grammar.

\subsubsection{Backtracking and Committed Choice}

Our simple \texttt{<|>} always backtracks on failure. Parsec introduced a
crucial refinement: \textbf{committed choice}. Once a parser has consumed
input successfully, it \emph{commits} to that branch---it won't backtrack.
This gives two benefits:

\textbf{Better error messages.} If you write \verb|\x. y| and the parser
successfully consumes the \verb|\| and \texttt{x} but fails on what comes
next, you want ``expected \texttt{.} at position 3,'' not ``expected
identifier''  (which is what you'd get if it backtracked and tried the
\texttt{app} branch instead).

\textbf{Better performance.} Without committed choice, the parser might try
exponentially many alternatives. With it, once any input is consumed, the
decision is final---no more alternatives are explored.

Parsec provides a \texttt{try} combinator that explicitly enables backtracking
for cases where you need it:

\begin{minted}{lean}
-- `try p`: run p, but if it fails AFTER consuming input, backtrack
-- (our simple version always backtracks; Parsec only backtracks with `try`)
def try_ (p : Parser α) : Parser α := fun s pos =>
  match p s pos with
  | some r => some r
  | none   => none   -- reset to original position (in Parsec, this matters)
\end{minted}

The rule of thumb: if two alternatives start with the same token, wrap the
first in \texttt{try}. Otherwise, let committed choice give you speed and
error quality for free.

\subsubsection{From Wadler to Parsec: The History}

Parser combinators have a rich intellectual history:

\textbf{Wadler (1985)} showed that recursive descent parsers can be expressed
as higher-order functions in a lazy functional language, and that monadic
sequencing captures the essence of sequential parsing.

\textbf{Hutton and Meijer (1996)} wrote the influential paper ``Monadic Parser
Combinators,'' which crystallized the approach: a parser is a monad, and
the monad interface (\texttt{return}, \texttt{>>=}) gives you sequencing
for free.

\textbf{Leijen (2001)} created Parsec, which solved the two practical problems
that had kept combinators out of production use: error messages (by tracking
source positions and ``expected'' labels) and performance (by introducing
committed choice). Parsec demonstrated that combinator parsers could handle
real programming languages like Haskell itself.

The idea spread everywhere: Scala's \texttt{scala-parser-combinators}, Rust's
\texttt{nom} and \texttt{combine} crates, Python's \texttt{parsy}, and
JavaScript's \texttt{parsimmon} all follow the Parsec pattern. Lean itself has
a sophisticated parser combinator library built into its front end. The idea
of ``small composable units combined with operators'' is one of functional
programming's most successful exports to mainstream software engineering.

\subsection{Named-to-De-Bruijn Conversion}

Once we have a named \texttt{Term}, we convert to de Bruijn. The key: maintain
a list of names where each name's \emph{position} is its de Bruijn index.

\begin{exercisebox}[Conversion challenge]
\stepcounter{exercise}\textbf{\theexercise.}\; If the context is \texttt{["x", "y"]} (meaning $x$ was bound most recently, $y$ before that), what de Bruijn index should $x$ get? What about $y$?

\emph{Answer:} $x$ is at position 0, $y$ at position 1.
\end{exercisebox}

\begin{minted}{lean}
def toDeBruijn (ctx : List String := []) : Term → Option Expr
  | Term.var x =>
    match ctx.indexOf? x with
    | some i => some (Expr.bvar i)    -- position IS the index
    | none   => none                  -- unbound variable
  | Term.lam x body =>
    match toDeBruijn (x :: ctx) body with   -- push x at index 0
    | some e => some (Expr.lam e)
    | none   => none
  | Term.app e₁ e₂ => do
    let e₁' ← toDeBruijn ctx e₁
    let e₂' ← toDeBruijn ctx e₂
    some (Expr.app e₁' e₂')
\end{minted}

\subsection{A Definitions Environment}

Typing Church encodings by hand is painful. We add a definitions map: when
the parser sees a known name like \texttt{succ} or \texttt{true}, it expands
it to the corresponding $\lam$-term. This is a simple tree-walk replacement
(respecting shadowing). See Appendix~\ref{app:code} for the full
\texttt{expandDefs} and \texttt{churchDefs} implementations.

%======================================================================
\section{Putting It All Together: The REPL}\label{sec:repl}
%======================================================================

\subsection{What a REPL Is}

REPL stands for \textbf{Read--Eval--Print Loop}. The idea is ancient---Lisp
had one in 1958---and simple: read user input, evaluate it, print the result,
repeat. Every interactive language environment uses this pattern: Python's
\texttt{>>>} prompt, Node.js's console, GHCi for Haskell, and Lean's
\texttt{\#eval}.

The name describes the four phases precisely:

\begin{enumerate}[leftmargin=2em]
\item \textbf{Read}: Accept a line of text from the user.
\item \textbf{Eval}: Parse it, transform it, and evaluate it.
\item \textbf{Print}: Display the result (or an error message).
\item \textbf{Loop}: Go back to step 1.
\end{enumerate}

What makes a REPL more than just ``run code and show output'' is the
\emph{loop}: state accumulates across iterations. Definitions made in one
line are available in the next. This turns a REPL from a calculator into
an interactive environment for exploration.

\subsection{The Pipeline}

Our REPL ties every component we've built into a single pipeline. Each arrow
represents a function we implemented in an earlier section:

\begin{center}
\texttt{string}
$\;\xrightarrow{\;\text{parse}\;}\;$
\texttt{Term}
$\;\xrightarrow{\;\text{expand}\;}\;$
\texttt{Term}
$\;\xrightarrow{\;\text{toDeBruijn}\;}\;$
\texttt{Expr}
$\;\xrightarrow{\;\text{eval}\;}\;$
\texttt{Expr}
$\;\xrightarrow{\;\text{display}\;}\;$
\texttt{string}
\end{center}

\begin{center}
\renewcommand{\arraystretch}{1.4}
\begin{tabular}{lll}
\textbf{Stage} & \textbf{What it does} & \textbf{Section} \\
\hline
Parse & String $\to$ named AST & Section~\ref{sec:parser} \\
Expand & Substitute known definitions & Section~\ref{sec:parser} \\
To de Bruijn & Named $\to$ de Bruijn indices & Section~\ref{sec:debruijn} \\
Eval & Beta-reduce to normal form & Section~\ref{sec:beta} \\
Display & De Bruijn AST $\to$ readable string & --- \\
\end{tabular}
\end{center}

The pipeline function itself is a composition of these stages, where each stage
can fail:

\begin{minted}{lean}
def pipeline (defs : Defs) (input : String) : Option String := do
  let named    ← parse input           -- stage 1: parse
  let expanded := expandDefs defs named -- stage 2: expand definitions
  let expr     ← toDeBruijn [] expanded -- stage 3: named → de Bruijn
  let result   := eval 1000 expr       -- stage 4: evaluate (with fuel)
  pure (display result)                -- stage 5: display
\end{minted}

Notice the \texttt{do} notation---this is the same monadic sequencing we saw
in the parser combinators. The \texttt{←} operator extracts a value from an
\texttt{Option}, and if any step returns \texttt{none}, the whole pipeline
short-circuits and returns \texttt{none}. Monads everywhere.

If any stage fails (parse error, unbound variable), the REPL reports the error
and waits for the next input---it never crashes.

\subsection{State and Commands}

Beyond evaluating expressions, our REPL supports commands. The key design
decision: \texttt{defs} is threaded through the loop as state. When you
\texttt{:let two = succ (succ zero)}, the name \texttt{two} is added to
\texttt{defs} and available in all subsequent inputs.

\begin{minted}{lean}
-- The REPL loop
partial def repl (defs : Defs) : IO Unit := do
  IO.print "λ> "
  let input ← (← IO.getStdin).getLine
  let input := input.trim
  match input with
  | ":quit"  => pure ()
  | ":defs"  => printDefs defs; repl defs
  | _        =>
    if input.startsWith ":let " then
      -- :let name = expr — add a new definition
      match parseLetBinding input with
      | some (name, term) =>
        IO.println s!"{name} defined"
        repl (defs.insert name term)
      | none => IO.println "Parse error"; repl defs
    else
      -- Evaluate an expression
      match pipeline defs input with
      | some result => IO.println s!"Result:    {result}"
      | none        => IO.println "Error: parse or conversion failed"
      repl defs
\end{minted}

This is a \textbf{functional loop}: there is no mutable state. Instead, the
updated \texttt{defs} map is passed as an argument to the next recursive call.
Each iteration of the REPL is a fresh function call with potentially different
state. This pattern---looping via recursion, state via arguments---is the
functional alternative to \texttt{while} loops with mutable variables.

\subsection{Example Session}

\begin{minted}{text}
λ> \x. x
Parsed:    (λx. x)
De Bruijn: (λ. 0)
Result:    (λ. 0)

λ> succ (succ (succ zero))
Result:    (λ. (λ. (0 (0 (0 1)))))

λ> add (succ (succ zero)) (succ (succ (succ zero)))
Result:    (λ. (λ. (0 (0 (0 (0 (0 1)))))))

λ> :let two = succ (succ zero)
two defined

λ> mul two (succ two)
Result:    (λ. (λ. (0 (0 (0 (0 (0 (0 1))))))))

λ> and true false
Result:    (λ. (λ. 1))

λ> :let apply_twice = \f x. f (f x)
apply_twice defined

λ> apply_twice succ zero
Result:    (λ. (λ. (0 (0 1))))
\end{minted}

That last result---$(\lam.\, \lam.\, 0\;(0\;1))$---is the Church numeral
$\church{2}$: applying successor twice to zero gives 2.

\subsection{Extending to a Typed REPL}

Our REPL evaluates untyped lambda terms. But the components are modular enough
to extend. Let's sketch how to add each feature from the type system sections.

\subsubsection{Adding Type Checking}

Insert a type-checking stage between conversion and evaluation. The pipeline
becomes:

\begin{center}
\texttt{string}
$\;\xrightarrow{\;\text{parse}\;}\;$
\texttt{Term}
$\;\xrightarrow{\;\text{convert}\;}\;$
\texttt{Expr}
$\;\xrightarrow{\;\text{\textbf{typecheck}}\;}\;$
\texttt{Ty}
$\;\xrightarrow{\;\text{eval}\;}\;$
\texttt{Expr}
$\;\xrightarrow{\;\text{display}\;}\;$
\texttt{string}
\end{center}

\begin{minted}{lean}
def typedPipeline (defs : Defs) (input : String) : Option String := do
  let named    ← parse input
  let expanded := expandDefs defs named
  let expr     ← toDeBruijn [] expanded
  let ty       ← typeCheck [] expr      -- NEW: reject ill-typed terms
  let result   := eval 1000 expr
  pure s!"{display result} : {displayTy ty}"
\end{minted}

Now \texttt{(\textbackslash x. x x)} is rejected before evaluation---exactly
what the STLC was designed to do.

\subsubsection{Adding Type Annotations}

To type-check lambda terms, the type checker needs to know the types of
parameters. We extend both the AST and the parser.

\paragraph{Extending the AST.}
The named term type gains an optional type annotation on lambdas:

\begin{minted}{lean}
inductive Ty where
  | nat  : Ty
  | bool : Ty
  | arr  : Ty → Ty → Ty     -- function type: Ty → Ty
  deriving Repr

inductive Term where
  | var : String → Term
  | lam : String → Ty → Term → Term     -- now: λ(x : T). body
  | app : Term → Term → Term
\end{minted}

\paragraph{Parsing types.}
The type grammar is its own small parser. Function arrows are
\emph{right}-associative ($A \to B \to C$ means $A \to (B \to C)$):

\begin{minted}{lean}
-- Parse a base type: Nat, Bool, or (type)
partial def atomTyP : Parser Ty :=
  (do keyword "Nat"; pure .nat)
  <|> (do keyword "Bool"; pure .bool)
  <|> (do let _ ← symbol '('; let t ← typeP; let _ ← symbol ')'; pure t)

-- Parse a function type: base -> base -> ... (right-associative)
partial def typeP : Parser Ty := do
  let lhs ← atomTyP
  (do let _ ← token (char '-' >>= fun _ => char '>') -- parse "->"
      let rhs ← typeP                                 -- right-recursive!
      pure (.arr lhs rhs))
  <|> pure lhs
\end{minted}

\paragraph{Parsing typed lambdas.}
The lambda parser now expects a colon and type annotation after each parameter:

\begin{minted}{lean}
-- Parse one typed parameter: ident : type
partial def typedParam : Parser (String × Ty) := do
  let name ← ident
  let _ ← symbol ':'
  let ty ← typeP
  pure (name, ty)

-- Parse a typed lambda: \(x : Nat)(y : Bool). body
-- Each parameter is parenthesized with its type
partial def typedLamP : Parser Term := do
  let _ ← symbol '\\'
  let params ← many1 (do
    let _ ← symbol '('
    let p ← typedParam
    let _ ← symbol ')'
    pure p)
  let _ ← symbol '.'
  let body ← exprP
  pure (params.foldr (fun (name, ty) acc => .lam name ty acc) body)
\end{minted}

\paragraph{Example.}
With this grammar, the identity function on naturals is:

\begin{minted}{text}
λ> \(x : Nat). x
Parsed:    λ(x : Nat). x
Type:      Nat → Nat
Result:    (λ. 0)

λ> \(f : Nat -> Nat)(x : Nat). f x
Parsed:    λ(f : Nat → Nat). λ(x : Nat). f x
Type:      (Nat → Nat) → Nat → Nat
Result:    (λ. (λ. (1 0)))
\end{minted}

The pipeline now produces typed output. If you try to apply a boolean to a
function expecting a natural, the type checker rejects it before evaluation
even begins.

\subsubsection{Adding Commands}

A production REPL supports multiple commands:

\begin{center}
\renewcommand{\arraystretch}{1.4}
\begin{tabular}{lp{8cm}}
\textbf{Command} & \textbf{What it does} \\
\hline
\texttt{:type expr} & Show the type without evaluating (like GHCi's
  \texttt{:t}) \\
\texttt{:let name = expr} & Bind a name to a term \\
\texttt{:step expr} & Show each reduction step, not just the result \\
\texttt{:debruijn expr} & Show the de Bruijn representation \\
\texttt{:defs} & List all current definitions \\
\texttt{:quit} & Exit \\
\end{tabular}
\end{center}

Each command is just a different path through the pipeline. \texttt{:type} runs
parse $\to$ convert $\to$ typecheck and stops. \texttt{:step} runs parse $\to$
convert $\to$ eval-with-tracing, printing each intermediate form.

\begin{intuition}[What you've built]
Starting from three constructs (variable, lambda, application) and one rule
(beta reduction), you've built a system that can do arithmetic, logic,
recursion, and more---all from functions. The interpreter follows the exact
same structure as industrial implementations of functional languages, just
without the optimizations. GHC (Haskell), OCaml's compiler, and Lean's own
kernel all follow this pattern: parse source text into an AST, elaborate it
into a core calculus, type-check, then evaluate by repeated reduction.

The pipeline you've built is a miniature version of what every programming
language implementation does. The components are the same---lexer, parser,
AST, type checker, evaluator, REPL loop---just at a smaller scale. If you
understand why each piece exists and how they connect, you understand the
architecture of every compiler and interpreter.
\end{intuition}

%======================================================================

%======================================================================
\section{The Landscape: From STLC to Lean}\label{sec:landscape-sec}
%======================================================================


\subsection{The Historical Arc}

The progression from STLC to modern proof assistants spans half a century of
interplay between logic, mathematics, and computer science. Here is the story
in compressed form:

\begin{center}
\renewcommand{\arraystretch}{1.5}
\begin{tabular}{llp{5cm}l}
\textbf{System} & \textbf{Year} & \textbf{Key addition} & \textbf{People} \\
\hline
STLC & 1940 & Simple types & Church \\
System F & 1972 & Polymorphism ($\forall\alpha$) & Girard, Reynolds \\
ML type inference & 1978 & Hindley--Milner & Milner, Hindley \\
Martin-L\"of TT & 1971--84 & Dependent types & Martin-L\"of \\
CoC & 1985 & All axes of the cube & Coquand, Huet \\
CIC & 1989 & $+$ inductive types & Paulin-Mohring \\
Coq & 1989 & First mature proof assistant & Huet, Coquand et al. \\
Agda & 1999--2007 & Dependent types with patterns & Norell, Bove \\
Lean 4 & 2021 & CIC $+$ practical PL features & de Moura \\
\end{tabular}
\end{center}

Each row builds on the ones above it. Church invented simple types to fix the
logical inconsistency Kleene and Rosser found. Girard added polymorphism to
capture second-order logic. Martin-L\"of added dependent types to capture
constructive mathematics. Coquand unified them all. And de Moura made the
result into a practical programming language.

\subsection{What Each Extension Costs}

Every extension pays for its expressiveness. The trade-offs reveal a deep
tension between power and decidability:

\textbf{STLC:} Type checking is decidable and simple (our 15-line type
checker). Type \emph{inference} is also decidable---you don't need any
annotations. But the system is too weak: you can't write generic code, you
can't express recursion, and Church numerals can't be given a single type.

\textbf{System F:} Type checking is still decidable. But type
\emph{inference} is undecidable (Wells, 1999)---sometimes you must provide
type annotations. The Hindley--Milner restriction trades full System F for
complete inference, which is why Haskell and ML can often omit annotations
while Java and Rust cannot.

\textbf{Dependent types:} Type \emph{checking} requires evaluation, so it's
only decidable if all programs terminate. This is the origin of Lean's
termination checker and its insistence on structural recursion. Without
termination, the type checker could loop forever. The reward: types can express
arbitrary mathematical propositions, and type checking is proof verification.

\subsection{The Philosophical Thread}

The deepest lesson of this progression is that \textbf{types are the
language of invariants}. At every level, the question is the same: what
properties can you express and verify statically?

In the STLC, you can say ``this function takes a number and returns a
boolean.'' In System F, you can say ``this function works at every type
uniformly.'' With dependent types, you can say ``this sorting function
returns a list that is a permutation of the input and is in ascending
order''---and the type checker will verify it.

This perspective connects back to Hilbert's original question that started
this entire story (Section~\ref{sec:what}). Hilbert wanted a mechanical
procedure that could verify mathematical proofs. The Entscheidungsproblem
asked: is there an algorithm that, given a statement in formal logic, decides
whether the statement is provable?

Church and Turing showed the answer is ``no'' in general. But dependent type
theory gives a remarkable partial answer: \emph{if you express your
propositions as types, then type checking is your proof verifier}. You can't
automatically \emph{find} proofs (that's still undecidable), but you can
automatically \emph{check} them. The type checker is the mechanical proof
verifier Hilbert was looking for---with the caveat that someone still has to
write the proof (the program) by hand.

This is exactly what proof assistants do. When a mathematician formalizes a
theorem in Lean, they are writing a program. When Lean's type checker accepts
it, Lean has mechanically verified the proof. The 4539-line Lean kernel---a
dependent type checker at its heart---is the closest thing we have to
Hilbert's dream, built on the lambda calculus that Church invented to show
the dream was impossible in its original form.

\begin{intuition}[The unifying theme]
Every system in this progression is a lambda calculus. The syntax is the same
three constructs: variables, abstraction, application. The computation rule is
the same: beta reduction. What changes is the \emph{type system}---the rules
that determine which terms are well-formed. The STLC is the starting point.
Each extension adds a new kind of abstraction (over types, over type operators,
or over values inside types), which makes the system more expressive while
(carefully) preserving the properties that make it useful.

This is why understanding the STLC matters: it's not just a toy system. It's
the foundation that every modern typed language builds on.
\end{intuition}

%======================================================================
\section{Common Mistakes and FAQ}\label{sec:faq}
%======================================================================

This section collects the most common points of confusion for newcomers. If
something earlier didn't quite click, you may find the answer here.

\subsection{``Can I write $\lam(x, y).\, e$ for a function of two arguments?''}

No. Lambda calculus functions take exactly one argument---this is fundamental to
the system, not just a syntactic limitation. The grammar allows only $\lam x.\,
e$ with a single variable after the $\lam$. Multi-argument functions are
handled via currying (Section~\ref{sec:currying}): $\lam x.\, \lam y.\, e$ is
a function that takes $x$ and returns \emph{another function} that takes $y$.

The shorthand $\lam x\, y.\, e$ is just notation for $\lam x.\, \lam y.\, e$.
It does not change the underlying system.

\subsection{``What's the difference between $=$ and $\aeq$?''}

$=$ is used informally to mean ``is the same term'' or ``we define this to be.''
$\aeq$ ($\alpha$-equivalence) is a specific technical relation: two terms are
$\alpha$-equivalent if they differ only in the names of their bound variables.

For example, $\lam x.\, x \;\aeq\; \lam y.\, y$ because they are the same
function with different parameter names. But we would not write
$\lam x.\, x = \lam y.\, y$ in a formal context, because they are literally
different strings of symbols.

In practice, most treatments (and our de Bruijn implementation) identify
$\alpha$-equivalent terms---they are considered ``the same'' for all purposes.

\subsection{``Why does my reduction seem to go in circles?''}

Some terms genuinely don't terminate. The classic example is $\Omega = (\lam
x.\, x\;x)\;(\lam x.\, x\;x)$, which reduces to itself forever
(Section~\ref{sec:beta}).

If you're reducing by hand and keep getting the same term back, you probably
have a non-terminating term. Try applying a different evaluation strategy---but
be aware that some terms have \emph{no} normal form under any strategy.

If you're using the interpreter and it returns the same thing you put in, the
``fuel'' (step limit) may have run out. Try increasing it.

\subsection{``Is $\lam x.\, x\; x$ typeable?''}

No, not in the simply typed lambda calculus (STLC). For $x\;x$ to be
well-typed, $x$ must simultaneously be a function (type $A \to B$) and its own
argument (type $A$). This requires $A = A \to B$, an infinite type. The STLC
has no infinite types and no recursive types, so this term is untypeable.

This is why the $\mathbf{Y}$ combinator (which contains $x\;x$) also cannot be
typed in the STLC. If you want recursion in a typed system, you need either
an explicit \texttt{fix} operator, recursive types, or a more expressive type
system like System F.

\subsection{``What's the difference between $\reduce$ and $\reduces$?''}

$\reduce$ (single-headed arrow) means \textbf{one step} of $\beta$-reduction:
exactly one redex was reduced. $\reduces$ (double-headed arrow) means
\textbf{zero or more steps}: the term reduces to the target after some
(possibly zero) number of single steps. So $e \reduces e'$ means ``there is a
chain $e \reduce e_1 \reduce e_2 \reduce \cdots \reduce e'$.'' If the chain
has zero steps, it means $e = e'$.

\subsection{``$\church{0}$ and $\FL$ are the same term. Isn't that a problem?''}

Yes, they are both $\lam t\, f.\, f$ (select the second argument).
$\textsc{Nil}$ is also the same term. This is a genuine limitation of Church
encoding: different ``types'' of data can collide because there are no types in
untyped $\lam$-calculus to distinguish them.

In practice, this means Church-encoded programs must be disciplined about which
``type'' of value they expect. A function expecting a boolean could be given a
number and would not ``know the difference.'' This is one of the motivations for
adding a type system (Section~\ref{sec:stlc}), which ensures such confusions
are caught at compile time.

\subsection{``Why do I need de Bruijn indices? Named variables seem fine.''}

Named variables are fine for pen-and-paper work, but they have a nasty
implementation problem: $\alpha$-conversion. Every time you perform
substitution, you must check for variable capture and potentially rename
binders. This makes the code more complex and error-prone.

De Bruijn indices eliminate the problem entirely. Two $\alpha$-equivalent
terms are \emph{literally identical} when written with de Bruijn indices:
$\lam x.\, x$ and $\lam y.\, y$ both become $\lam.\, 0$. No renaming is ever
needed, and structural equality is $\alpha$-equality. The trade-off is that
de Bruijn terms are harder for humans to read, but easier for machines to
manipulate.

\subsection{``Can I use lambda calculus as an actual programming language?''}

Technically yes---it is Turing-complete, so it can compute anything. In
practice, raw $\lam$-calculus is far too slow and unreadable for real work.
Every functional programming language is essentially $\lam$-calculus with
ergonomic and performance improvements bolted on: named functions, pattern
matching, type checking, native data types, efficient compilation, and so on.

The value of studying $\lam$-calculus is not to use it directly, but to
understand the \emph{foundations} that all these languages are built on.

\subsection{``Why is it called `lambda'? What does the $\lam$ symbol mean?''}

The origin is surprisingly mundane. Church originally used a circumflex
($\hat{x}$) to mark bound variables, then switched to $\Lambda x$ for
typographical reasons, which was eventually simplified to $\lam x$. The symbol
$\lam$ has no deep mathematical meaning---it is just a marker that says
``a function is being defined here.''

%======================================================================
\part{Building a Lambda Calculus Interpreter}
%======================================================================

You've learned the theory. Now imagine you want to \emph{build} it---a program
that takes a $\lam$-expression as input, reduces it, and shows you the result.
This part guides you through the process of designing and implementing such a
program.

The approach here is deliberately \emph{not} ``here's the code, memorize it.''
Instead, each section poses a design challenge, gives you the tools and
intuition to solve it yourself, and then reveals the answer. \textbf{Try each
challenge before reading on.} Even a few minutes of honest struggle---
sketching on paper, writing pseudo-code, getting stuck and thinking about
\emph{why} you're stuck---will teach you far more than passively reading the
solution.

The reason this works is that every piece of implementation code is a direct
translation of the theory from the earlier sections. If you understand why
substitution has six cases, you can derive the substitution function. If you
understand why de Bruijn indices exist, you can figure out what shifting must
do. The code doesn't come from nowhere---it comes from the math, and the math
comes from the problems the math was invented to solve. Our goal is for you to
see each piece of code as \emph{inevitable}: the only thing the implementation
\emph{could} look like, given the theory.

The complete, copy-pasteable implementation is collected in
Appendix~\ref{app:code}. But \textbf{don't skip ahead}---the appendix is a
reference, not a tutorial. The learning happens here.

We use Lean\,4 as our implementation language because its own type system is
built on a $\lam$-calculus, giving a nice meta-circular elegance. But the ideas
translate to any language. If you prefer, read the Lean code as pseudocode---the
patterns (algebraic data types, pattern matching, recursion) exist in Python,
Haskell, Rust, OCaml, TypeScript, and more.

%======================================================================
\section{Where to Go from Here}
%======================================================================

\textbf{Prove correctness.} Lean is a proof assistant---prove that evaluation
preserves types ($\Gamma \proves e : \tau \;\wedge\; e \reduce e'
\;\Longrightarrow\; \Gamma \proves e' : \tau$).

\textbf{System F.} Add polymorphism: $\Lam\alpha.\, \lam x{:}\alpha.\, x
\;:\; \forall\alpha.\, \alpha \to \alpha$.

\textbf{Dependent types.} Implement a language where types depend on values,
exploring the Curry--Howard correspondence between proofs and programs.

\textbf{Compilation.} Replace tree-walking with a stack machine, G-machine, or
CPS transform for real performance.

\newpage
%======================================================================
\appendix
\section{Quick Reference Card}\label{app:ref}
%======================================================================

\subsection*{Core Syntax and Rules}
\vspace{-0.5em}
\begin{align*}
  e &::= x \mid \lam x.\, e \mid e_1\; e_2 \\
  \alpha &:\quad \lam x.\, e \;\aeq\; \lam y.\, e[x := y] \quad(y \text{ fresh}) \\
  \beta &:\quad (\lam x.\, e)\; a \;\reduce\; e[x := a] \\
  \eta &:\quad \lam x.\, f\; x \;\ef\; f \quad(x \notin \FV(f))
\end{align*}

\subsection*{Famous Combinators}
\vspace{-0.5em}
\begin{align*}
  \mathbf{I} &= \lam x.\, x     &  \mathbf{K} &= \lam x\, y.\, x     &  \mathbf{S} &= \lam f\, g\, x.\, f\;x\;(g\;x) \\
  \Omega &= (\lam x.\, x\;x)^2 &
  \mathbf{Y} &= \lam f.\, (\lam x.\, f(x\;x))^2
\end{align*}

\subsection*{Church Encodings}
\vspace{-0.5em}
\begin{align*}
  \TR &= \lam t\, f.\, t   &  \FL &= \lam t\, f.\, f \\
  \church{n} &= \lam f\, x.\, f^n\;x &
  \textsc{Succ} &= \lam n\, f\, x.\, f\;(n\;f\;x) \\
  \textsc{Add} &= \lam m\, n\, f\, x.\, m\;f\;(n\;f\;x) &
  \textsc{Mul} &= \lam m\, n\, f.\, m\;(n\;f) \\
  \textsc{Pair} &= \lam a\, b\, f.\, f\;a\;b &
  \textsc{Fst} &= \lam p.\, p\;\TR
\end{align*}

\subsection*{De Bruijn}
\vspace{-0.5em}
\[
  (\lam.\, e)\; s \;\reduce\; \uparrow^{-1}_0\!\bigl(e[0 := \uparrow^1_0(s)]\bigr)
\]

\subsection*{STLC Typing Rules}

\begin{center}
\AxiomC{$(x : \tau) \in \Gamma$}
\RightLabel{\;\textsc{Var}}
\UnaryInfC{$\Gamma \proves x : \tau$}
\DisplayProof
\quad
\AxiomC{$\Gamma, x : \tau_1 \proves e : \tau_2$}
\RightLabel{\;\textsc{Abs}}
\UnaryInfC{$\Gamma \proves \lam x{:}\tau_1.\, e : \tau_1 \to \tau_2$}
\DisplayProof
\quad
\AxiomC{$\Gamma \proves e_1 : \tau_1 \to \tau_2$}
\AxiomC{$\Gamma \proves e_2 : \tau_1$}
\RightLabel{\;\textsc{App}}
\BinaryInfC{$\Gamma \proves e_1\; e_2 : \tau_2$}
\DisplayProof
\end{center}

\bigskip

%======================================================================
\section{Glossary}\label{app:glossary}
%======================================================================

\renewcommand{\arraystretch}{1.5}
\begin{tabular}{lp{10cm}}
\textbf{Term} & \textbf{Definition} \\
\hline
Abstraction & $\lam x.\, e$: a function with parameter $x$ and body $e$. \\
$\alpha$-conversion & Renaming bound variables; $\lam x.\,x \aeq \lam y.\,y$. \\
$\beta$-reduction & $(\lam x.\, e)\;a \reduce e[x := a]$. The computation rule. \\
$\eta$-conversion & $\lam x.\, f\;x = f$ when $x \notin \FV(f)$. Extensionality. \\
Church encoding & Representing data as pure $\lam$-terms. \\
Church--Rosser & Normal forms are unique up to $\aeq$. \\
Combinator & A closed term ($\FV(e) = \varnothing$). \\
Context ($\Gamma$) & A list of variable-type bindings for judgments. \\
De Bruijn index & Nameless variable: a number counting binders outward. \\
Fixed-point combinator & $\mathbf{Y}$ with $\mathbf{Y}\;f = f\;(\mathbf{Y}\;f)$. \\
Inference rule & Premises above a line, conclusion below. \\
Normal form & A term with no redexes; fully reduced. \\
Redex & A reducible expression $(\lam x.\, e)\;a$. \\
Substitution & $e[x := a]$: replace free $x$ in $e$ with $a$, avoiding capture. \\
Turnstile ($\proves$) & Separates context from conclusion in $\Gamma \proves e : \tau$. \\
WHNF & A term that is a $\lam$ or has no top-level redex. \\
\end{tabular}

\newpage
%======================================================================
\section{Exercise Solutions}\label{app:solutions}
%======================================================================

Solutions are provided in condensed form. Working through the exercises
yourself before reading the solutions is strongly recommended.

\subsection*{Variables and Renaming (Section~\ref{sec:alpha})}

\textbf{5.1.}\; $\FV(\lam x.\, \lam y.\, x\; z\; y) = \{z\}$. The variables $x$ and $y$ are bound by their respective $\lam$'s. Only $z$ is free.

\medskip
\textbf{5.2.}\; Yes, $\lam f.\, \lam x.\, f\;(f\;x)$ is closed. $\FV = \varnothing$: both $f$ and $x$ are bound. This is the Church numeral $\church{2}$.

\medskip
\textbf{5.3.}\; The inner $\lam x$ shadows the outer one. The inner $x$ refers to the inner binder. Alpha-convert the outer: $\lam x.\, \lam x.\, x \;\aeq\; \lam y.\, \lam x.\, x$. This is $\mathbf{K}^* = \lam y.\, \lam x.\, x$: a function that takes two arguments and returns the \emph{second} (the opposite of $\mathbf{K}$). Equivalently, it is $\lam y.\, \mathbf{I}$.

\medskip
\textbf{5.4.}\; No. $\lam x.\, x\;y$ has $\FV = \{y\}$, but $\lam y.\, y\;y$ has $\FV = \varnothing$. Alpha conversion can only rename bound variables---it cannot change the set of free variables. These terms have different free variable sets, so they cannot be $\alpha$-equivalent.

\subsection*{Beta Reduction (Section~\ref{sec:beta})}

\textbf{6.1.}\; $(\lam x.\, x\;x)\;(\lam y.\, y) \;\reduce\; (\lam y.\, y)\;(\lam y.\, y) \;\reduce\; \lam y.\, y$. The normal form is the identity $\mathbf{I}$.

\medskip
\textbf{6.2.}\; Let $\omega = \lam y.\, y\;y$. Then:
$(\lam f.\, \lam x.\, f\;(f\;x))\;\omega \;\reduce\; \lam x.\, \omega\;(\omega\;x)$.
Now apply to some $a$:
$\omega\;(\omega\;a) = (\lam y.\, y\;y)\;(\omega\;a) \reduce (\omega\;a)\;(\omega\;a)$, which grows without bound.
The term is \emph{not} in normal form and does not terminate if we try to fully normalize the body.

\medskip
\textbf{6.3.}\; Two redexes: (1) the outer application $(\lam x.\, x)\;((\lam y.\, y)\;z)$ is a redex (the whole expression), and (2) the inner application $(\lam y.\, y)\;z$ inside the argument is also a redex.

\medskip
\textbf{6.4.}\; No, $\lam x.\, (\lam y.\, y)\;x$ is not in normal form---the body contains the redex $(\lam y.\, y)\;x$. Reducing: $(\lam y.\, y)\;x \reduce x$, giving $\lam x.\, x$. Then yes, we can $\eta$-reduce\ldots\ but $\lam x.\, x$ is already the identity; there is no $f$ to $\eta$-reduce to (the body is $x$, not $f\;x$ for some $f$ with $x \notin \FV(f)$). So the final normal form is $\lam x.\, x = \mathbf{I}$.

\medskip
\textbf{6.5.}\; We need $(\lam y.\, x\;y)[x := \lam y.\, y]$. This matches rule 5: we have $\lam y$ with $y \in \FV(\lam y.\, y) = \varnothing$\ldots\ actually $y \notin \FV(\lam y.\, y)$ since the $y$ in $\lam y.\, y$ is bound. So rule 4 applies (safe case): $\lam y.\, (x\;y)[x := \lam y.\, y] = \lam y.\, (\lam y.\, y)\;y$. Wait---this looks like capture! But it isn't: the substituted $\lam y.\, y$ has $\FV = \varnothing$, so $y \notin \FV(\lam y.\, y)$ and rule 4 is correct. The result $\lam y.\, (\lam y.\, y)\;y$ reduces to $\lam y.\, y$, which is correct ($\eta$-equivalent to the original).

\subsection*{Church Encodings (Section~\ref{sec:church})}

\textbf{9.1.}\; $\textsc{IsZero}\;(\textsc{Succ}\;\church{0})$: First, $\textsc{Succ}\;\church{0} \reduces \church{1} = \lam f\,x.\, f\;x$. Then $\textsc{IsZero}\;\church{1} = \church{1}\;(\lam\_.\,\FL)\;\TR \reduces (\lam\_.\,\FL)\;\TR \reduce \FL$.

\medskip
\textbf{9.2.}\; $\textsc{Add}\;\church{2}\;\church{3} = (\lam m\,n\,f\,x.\, m\;f\;(n\;f\;x))\;\church{2}\;\church{3}$. Substituting: $\lam f\,x.\, \church{2}\;f\;(\church{3}\;f\;x) \reduces \lam f\,x.\, f\;(f\;(f\;(f\;(f\;x)))) = \church{5}$.

\medskip
\textbf{9.3.}\; $\textsc{Sub} \triangleq \lam m\, n.\, n\; \textsc{Pred}\; m$. Apply $\textsc{Pred}$ a total of $n$ times to $m$. Since $\church{n}$ applies its first argument $n$ times to its second, $n\;\textsc{Pred}\;m$ does exactly this.

\medskip
\textbf{9.4.}\; $\textsc{Length} \triangleq \lam l.\, l\; (\lam \_\, \mathit{acc}.\, \textsc{Succ}\;\mathit{acc})\; \church{0}$. Fold over the list: for each element (ignored), increment the accumulator. Start from $\church{0}$.

\medskip
\textbf{9.5.}\; Yes, $\textsc{Nil} = \lam c\, n.\, n$ and $\church{0} = \lam f\, x.\, x$ are the same term (up to renaming: $c \leftrightarrow f$, $n \leftrightarrow x$). They are $\alpha$-equivalent. This illustrates a fundamental limitation of Church encoding in untyped $\lam$-calculus: different ``types'' of data can be identical terms. Only a type system can distinguish them.

\subsection*{Types (Section~\ref{sec:stlc})}

\textbf{11.1.}\; The proof tree for the $\mathbf{S}$ combinator has type $(A \to B \to C) \to (A \to B) \to A \to C$. It requires three \textsc{Abs} rules (one for each $\lam$), two \textsc{App} rules (for $x\;z$ and $y\;z$ and combining them), and three \textsc{Var} axioms ($x$, $y$, $z$):

\begin{center}
\small
\AxiomC{$x : A {\to} B {\to} C$}
\RightLabel{\textsc{V}}
\UnaryInfC{$\Gamma \proves x : A {\to} B {\to} C$}
\AxiomC{$z : A$}
\RightLabel{\textsc{V}}
\UnaryInfC{$\Gamma \proves z : A$}
\RightLabel{\textsc{App}}
\BinaryInfC{$\Gamma \proves x\;z : B {\to} C$}
\AxiomC{$y : A {\to} B$}
\RightLabel{\textsc{V}}
\UnaryInfC{$\Gamma \proves y : A {\to} B$}
\AxiomC{$z : A$}
\RightLabel{\textsc{V}}
\UnaryInfC{$\Gamma \proves z : A$}
\RightLabel{\textsc{App}}
\BinaryInfC{$\Gamma \proves y\;z : B$}
\RightLabel{\textsc{App}}
\BinaryInfC{$\Gamma \proves x\;z\;(y\;z) : C$}
\RightLabel{\textsc{Abs$\times$3}}
\UnaryInfC{$\proves \lam x\,y\,z.\, x\;z\;(y\;z) : (A {\to} B {\to} C) \to (A {\to} B) \to A \to C$}
\DisplayProof
\end{center}
where $\Gamma = x : A {\to} B {\to} C,\; y : A {\to} B,\; z : A$.

\medskip
\textbf{11.2.}\; For $x\;x$ to type-check, $x$ must have a function type $A \to B$ (left of the application) and also be a valid argument (type $A$). So we need $A \to B = A$, which means $A = A \to B = (A \to B) \to B = \ldots$\,---an infinite type. The STLC has only finite types, so $\lam x.\, x\;x$ is untypeable.

\medskip
\textbf{11.3.}\; $\lam f{:}(A \to A).\, \lam x{:}A.\, f\;(f\;x)$ has type $(A \to A) \to A \to A$. This is the same as $A \to A$ where we've substituted $A \to A$ for $A$. Every Church numeral $\church{n}$ has this same type. Under Curry--Howard, it's a proof of ``if $A$ implies $A$, then $A$ implies $A$''---which is trivially true, and there are infinitely many distinct proofs (one per numeral).

\subsection*{De Bruijn Substitution (Section~\ref{sec:dbimpl})}

\textbf{23.}\; The na\"ive version fails on $\lam.\,\lam.\, 0\;1\;2$:
indices $0$ and $1$ are bound (by the inner and outer $\lam$s respectively),
and index $2$ is the free variable that corresponds to external variable $0$.
\texttt{substNaive} with $k=0$ incorrectly replaces \texttt{bvar 0} (which is
the innermost $\lam$'s own parameter) with $s$, while leaving the actual target
(\texttt{bvar 2}) untouched.

\medskip
\textbf{24.}\; $\texttt{substV2}\;0\;(\texttt{bvar}\;0)\;(\lam.\, 1)$:
Named version is ``substitute $y$ for $x$ in $\lam z.\, x$.'' Result should
be $\lam z.\, y$, i.e., $\lam.\, 1$ (index 1 skips past $z$'s binder to
reach $y$). But $\texttt{substV2}$ produces $\lam.\, 0$---the replacement
$\texttt{bvar}\;0$ was dropped in unchanged, and inside the $\lam$ it points
to $z$ instead of $y$. The fix: shift $s$ up when entering the $\lam$.

\medskip
\textbf{25.}\; (a) Replace with $s$ when $n = k$; leave as $\texttt{bvar}\;n$
when $n \neq k$.
(b) Recurse: $\texttt{subst}\;k\;s$ into both sides independently.
(c) Both fixes apply: the target $k$ increments to $k + 1$ (fix 1) because the
variable we're hunting for is one more binder away, and the replacement $s$
gets shifted up by 1 via $\texttt{shift}\;1\;0\;s$ (fix 2) because its free
variables need to skip the new binder.
(d) Fix 1 without fix 2: correctly finds the target variable but inserts $s$
with wrong indices (capture bug from stage 2). Fix 2 without fix 1: shifts $s$
correctly but looks for the wrong variable inside the $\lam$ (would replace
bound variables or miss the target, like stage 1).

\medskip
\textbf{26.}\; $k$ goes $0 \to 1 \to 2 \to 3$. Yes, index $3$ inside three binders refers to the same variable as index $0$ outside. The substitution fires correctly.

\medskip
\textbf{27.}\; See the worked trace in the exercise box.

\medskip
\textbf{28.}\; Body is $\lam.\, 1$, argument is $\texttt{bvar}\;0$.

Step 1: $\texttt{shift}\;1\;0\;(\texttt{bvar}\;0) = \texttt{bvar}\;1$.

Step 2: $\texttt{subst}\;0\;(\texttt{bvar}\;1)\;(\lam.\, 1)$. We hit the $\lam$: recurse with $k = 1$, $s = \texttt{shift}\;1\;0\;(\texttt{bvar}\;1) = \texttt{bvar}\;2$.

Now: $\texttt{subst}\;1\;(\texttt{bvar}\;2)\;(\texttt{bvar}\;1)$. Is $1 = 1$? Yes! Result: $\texttt{bvar}\;2$.

After wrapping: $\lam.\, 2$.

Step 3: $\texttt{shift}\;(-1)\;0\;(\lam.\, 2)$. Inside the $\lam$, cutoff is 1. Is $2 \geq 1$? Yes: $2 - 1 = 1$. Result: $\lam.\, 1$.

Final: $\lam.\, \texttt{bvar}\;1$---a function that ignores its argument and returns the free variable. $\checkmark$

\subsection*{Going Deeper (Section~\ref{sec:curryhoward})}

\textbf{13.1.}\; $\textsc{Succ}_S \triangleq \lam n.\, \lam z\, s.\, s\; n$. This creates a new Scott numeral that, when pattern-matched, calls the successor handler $s$ with $n$ as the predecessor. Verification: $\textsc{Succ}_S\;\church{2}_S = (\lam n.\, \lam z\,s.\, s\;n)\;\church{2}_S \reduce \lam z\,s.\, s\;\church{2}_S = \church{3}_S$.

\medskip
\textbf{13.2.}\; $\textsc{IsZero}_S \triangleq \lam n.\, n\;\TR\;(\lam\_.\, \FL)$. If $n$ is zero, it selects the first handler ($\TR$). If $n$ is a successor, it calls the second handler (which ignores the predecessor and returns $\FL$). This is much simpler than the Church version---a single reduction step vs.\ the Church version's $O(n)$ iterations.

\medskip
\textbf{13.3.}\; $(A \to B) \to (B \to C) \to (A \to C)$ corresponds to the transitivity of implication: ``if $A$ implies $B$ and $B$ implies $C$, then $A$ implies $C$.'' A witnessing term: $\lam f{:}(A \to B).\, \lam g{:}(B \to C).\, \lam x{:}A.\, g\;(f\;x)$. This is function composition.

\medskip
\textbf{13.4.}\; There is no closed term of type $A \to B$ when $A \neq B$ because to produce a $B$, you would need either a variable of type $B$ or a function that returns $B$. But the only variable available is $x : A$, and there is no way to convert an $A$ into a $B$. Under Curry--Howard, this means the proposition ``$A$ implies $B$'' (for unrelated $A$ and $B$) is not provable---which is correct.

\medskip
\textbf{13.5.}\; $\forall \alpha.\, \alpha$ says ``for every type $\alpha$, I can produce a value of type $\alpha$.'' Under Curry--Howard, this is the proposition ``everything is true''---i.e., falsity ($\bot$). It is uninhabited because no single term can work at every type. (In System F, the term would need to simultaneously be a natural number, a boolean, a function, etc.)

\newpage
%======================================================================
\section{Complete Implementation}\label{app:code}
%======================================================================

This appendix collects the complete Lean\,4 implementation from the implementation sections in one
place, suitable for copy-pasting into a project. The code is organized by
module.

\subsection*{Project Layout}

\begin{minted}{text}
LambdaCalc/
├── lakefile.lean
├── LambdaCalc.lean
└── LambdaCalc/
    ├── Named.lean       -- Named terms, FV, substitution, eval
    ├── DeBruijn.lean    -- De Bruijn terms, shift, subst, eval
    ├── Parser.lean      -- Recursive descent parser
    ├── Defs.lean        -- Church encoding definitions + expandDefs
    ├── Typed.lean       -- STLC type checker
    └── Main.lean        -- REPL
\end{minted}

\subsection*{Named.lean --- Named Representation}

\begin{minted}{lean}
inductive Term where
  | var : String → Term
  | lam : String → Term → Term
  | app : Term → Term → Term
  deriving Repr, DecidableEq, Inhabited

open Term

def Term.toString : Term → String
  | .var x     => x
  | .lam x e   => s!"(λ{x}. {e.toString})"
  | .app e₁ e₂ => s!"({e₁.toString} {e₂.toString})"

instance : ToString Term := ⟨Term.toString⟩

def freeVars : Term → List String
  | .var x     => [x]
  | .lam x e   => (freeVars e).filter (· != x)
  | .app e₁ e₂ => (freeVars e₁ ++ freeVars e₂).eraseDups

def freshVar (avoid : List String) (base : String := "x") : String :=
  let rec go (n : Nat) : String :=
    let candidate := s!"{base}{n}"
    if candidate ∈ avoid then go (n + 1) else candidate
  if base ∉ avoid then base else go 0

def subst (x : String) (s : Term) : Term → Term
  | .var y     => if y == x then s else .var y
  | .app e₁ e₂ => .app (subst x s e₁) (subst x s e₂)
  | .lam y e   =>
    if y == x then .lam y e
    else if y ∉ freeVars s then
      .lam y (subst x s e)
    else
      let z := freshVar (freeVars s ++ freeVars e ++ [x])
      .lam z (subst x s (subst y (.var z) e))

def betaStep : Term → Option Term
  | .app (.lam x body) arg => some (subst x arg body)
  | .app e₁ e₂ =>
    match betaStep e₁ with
    | some e₁' => some (.app e₁' e₂)
    | none     => match betaStep e₂ with
                  | some e₂' => some (.app e₁ e₂')
                  | none     => none
  | .lam x e => match betaStep e with
               | some e' => some (.lam x e')
               | none    => none
  | _ => none

def eval (fuel : Nat) (t : Term) : Term :=
  match fuel with
  | 0         => t
  | fuel' + 1 => match betaStep t with
                 | none    => t
                 | some t' => eval fuel' t'
\end{minted}

\subsection*{DeBruijn.lean --- De Bruijn Representation}

\begin{minted}{lean}
inductive Expr where
  | bvar : Nat → Expr
  | lam  : Expr → Expr
  | app  : Expr → Expr → Expr
  deriving Repr, DecidableEq, Inhabited

open Expr

def Expr.toString : Expr → String
  | .bvar n    => s!"{n}"
  | .lam e     => s!"(λ. {e.toString})"
  | .app e₁ e₂ => s!"({e₁.toString} {e₂.toString})"

instance : ToString Expr := ⟨Expr.toString⟩

def shift (d : Int) (c : Nat) : Expr → Expr
  | .bvar n    => if n ≥ c then .bvar (Int.toNat (n + d)) else .bvar n
  | .lam e     => .lam (shift d (c + 1) e)
  | .app e₁ e₂ => .app (shift d c e₁) (shift d c e₂)

def subst (k : Nat) (s : Expr) : Expr → Expr
  | .bvar n    => if n == k then s else .bvar n
  | .lam e     => .lam (subst (k + 1) (shift 1 0 s) e)
  | .app e₁ e₂ => .app (subst k s e₁) (subst k s e₂)

def betaReduce (body arg : Expr) : Expr :=
  shift (-1) 0 (subst 0 (shift 1 0 arg) body)

-- Normal order (leftmost outermost)
def step : Expr → Option Expr
  | .app (.lam body) arg => some (betaReduce body arg)
  | .app e₁ e₂ =>
    match step e₁ with
    | some e₁' => some (.app e₁' e₂)
    | none     => match step e₂ with
                  | some e₂' => some (.app e₁ e₂')
                  | none     => none
  | .lam e => match step e with
             | some e' => some (.lam e')
             | none    => none
  | _ => none

-- Call-by-value
def isValue : Expr → Bool
  | .lam _ => true
  | _     => false

def cbvStep : Expr → Option Expr
  | .app (.lam body) arg =>
    if isValue arg then some (betaReduce body arg)
    else match cbvStep arg with
         | some arg' => some (.app (.lam body) arg')
         | none      => none
  | .app e₁ e₂ =>
    match cbvStep e₁ with
    | some e₁' => some (.app e₁' e₂)
    | none     => match cbvStep e₂ with
                  | some e₂' => some (.app e₁ e₂')
                  | none     => none
  | _ => none

def eval (fuel : Nat) (e : Expr) : Expr :=
  match fuel with
  | 0         => e
  | fuel' + 1 => match step e with
                 | none    => e
                 | some e' => eval fuel' e'
\end{minted}

\subsection*{Parser.lean --- Recursive Descent Parser}

\begin{minted}{lean}
structure PState where
  input : String
  pos   : String.Pos
  deriving Repr

abbrev ParseResult (α : Type) := Option (α × PState)

def skipWS (s : PState) : PState :=
  let rec go (p : String.Pos) : String.Pos :=
    if h : p < s.input.endPos then
      let c := s.input.get' p h
      if c == ' ' || c == '\t' then go (s.input.next' p h) else p
    else p
  { s with pos := go s.pos }

def peek (s : PState) : Option Char :=
  if h : s.pos < s.input.endPos
  then some (s.input.get' s.pos h) else none

def expect (s : PState) (c : Char) : ParseResult Unit :=
  if h : s.pos < s.input.endPos then
    if s.input.get' s.pos h == c
    then some ((), skipWS { s with pos := s.input.next' s.pos h })
    else none
  else none

def isIdentStart (c : Char) : Bool := c.isAlpha || c == '_'
def isIdentCont (c : Char) : Bool := c.isAlphanum || c == '_' || c == '\''

def parseIdent (s : PState) : ParseResult String :=
  let s := skipWS s
  if h : s.pos < s.input.endPos then
    let c := s.input.get' s.pos h
    if isIdentStart c then
      let rec go (p : String.Pos) : String.Pos :=
        if h2 : p < s.input.endPos then
          if isIdentCont (s.input.get' p h2)
          then go (s.input.next' p h2) else p
        else p
      let endP := go (s.input.next' s.pos h)
      let name := s.input.extract s.pos endP
      some (name, skipWS { s with pos := endP })
    else none
  else none

mutual
partial def parseExpr (s : PState) : ParseResult Term :=
  let s := skipWS s
  match peek s with
  | some '\\' => parseLambda s
  | _         => parseApp s

partial def parseLambda (s : PState) : ParseResult Term :=
  match expect s '\\' with
  | none => none
  | some (_, s) =>
    let rec parseParams (s : PState) (acc : List String)
        : ParseResult (List String) :=
      match parseIdent s with
      | some (name, s') => parseParams s' (acc ++ [name])
      | none => if acc.isEmpty then none else some (acc, s)
    match parseParams s [] with
    | none => none
    | some (params, s) =>
      match expect s '.' with
      | none => none
      | some (_, s) =>
        match parseExpr s with
        | none => none
        | some (body, s) =>
          let term := params.foldr (fun p acc => Term.lam p acc) body
          some (term, s)

partial def parseApp (s : PState) : ParseResult Term :=
  match parseAtom s with
  | none => none
  | some (first, s) =>
    let rec go (s : PState) (acc : Term) : Term × PState :=
      match parseAtom s with
      | some (next, s') => go s' (Term.app acc next)
      | none => (acc, s)
    let (result, s) := go s first
    some (result, s)

partial def parseAtom (s : PState) : ParseResult Term :=
  let s := skipWS s
  match peek s with
  | some '(' =>
    match expect s '(' with
    | none => none
    | some (_, s) =>
      match parseExpr s with
      | none => none
      | some (e, s) =>
        match expect s ')' with
        | none => none
        | some (_, s) => some (e, s)
  | _ =>
    match parseIdent s with
    | some (name, s) => some (Term.var name, s)
    | none => none
end

def parse (input : String) : Option Term :=
  let s : PState := skipWS { input := input, pos := 0 }
  match parseExpr s with
  | some (term, _) => some term
  | none           => none
\end{minted}

\subsection*{Defs.lean --- Definitions Environment}

\begin{minted}{lean}
def expandDefs (defs : List (String × Term)) : Term → Term
  | Term.var x =>
    match defs.lookup x with
    | some t => t
    | none   => Term.var x
  | Term.lam x e =>
    let defs' := defs.filter (fun (n, _) => n != x)
    Term.lam x (expandDefs defs' e)
  | Term.app e₁ e₂ =>
    Term.app (expandDefs defs e₁) (expandDefs defs e₂)

def churchDefs : List (String × Term) :=
  let t name := Term.var name
  let l := Term.lam
  let a := Term.app
  [ ("true",  l "t" (l "f" (t "t"))),
    ("false", l "t" (l "f" (t "f"))),
    ("and",   l "p" (l "q" (a (a (t "p") (t "q"))
                               (l "t" (l "f" (t "f")))))),
    ("or",    l "p" (l "q" (a (a (t "p")
                               (l "t" (l "f" (t "t"))))
                               (t "q")))),
    ("not",   l "p" (a (a (t "p")
                        (l "t" (l "f" (t "f"))))
                        (l "t" (l "f" (t "t"))))),
    ("zero",  l "f" (l "x" (t "x"))),
    ("succ",  l "n" (l "f" (l "x"
               (a (t "f") (a (a (t "n") (t "f")) (t "x")))))),
    ("add",   l "m" (l "n" (l "f" (l "x"
               (a (a (t "m") (t "f"))
                  (a (a (t "n") (t "f")) (t "x"))))))),
    ("mul",   l "m" (l "n" (l "f"
               (a (t "m") (a (t "n") (t "f")))))),
    ("id",    l "x" (t "x")),
    ("const", l "x" (l "y" (t "x")))
  ]

def toDeBruijn (ctx : List String := []) : Term → Option Expr
  | Term.var x =>
    match ctx.indexOf? x with
    | some i => some (Expr.bvar i)
    | none   => none
  | Term.lam x body =>
    match toDeBruijn (x :: ctx) body with
    | some e => some (Expr.lam e)
    | none   => none
  | Term.app e₁ e₂ => do
    let e₁' ← toDeBruijn ctx e₁
    let e₂' ← toDeBruijn ctx e₂
    some (Expr.app e₁' e₂')
\end{minted}

\subsection*{Typed.lean --- STLC Type Checker}

\begin{minted}{lean}
inductive Ty where
  | base  : String → Ty
  | arrow : Ty → Ty → Ty
  deriving Repr, DecidableEq

inductive TExpr where
  | bvar : Nat → TExpr
  | lam  : Ty → TExpr → TExpr
  | app  : TExpr → TExpr → TExpr
  deriving Repr

def typeCheck (ctx : List Ty) : TExpr → Option Ty
  | .bvar n          => ctx.get? n
  | .lam paramTy body =>
    match typeCheck (paramTy :: ctx) body with
    | some bodyTy => some (.arrow paramTy bodyTy)
    | none        => none
  | .app e₁ e₂ => do
    let funTy ← typeCheck ctx e₁
    let argTy ← typeCheck ctx e₂
    match funTy with
    | .arrow paramTy retTy =>
      if paramTy == argTy then some retTy else none
    | _ => none
\end{minted}

\subsection*{Main.lean --- REPL}

\begin{minted}{lean}
def evalString (defs : List (String × Term)) (input : String)
    : String := Id.run do
  let some parsed := parse input | return "Parse error"
  let expanded := expandDefs defs parsed
  let some db := toDeBruijn [] expanded | return "Unbound variable"
  let result := eval 10000 db
  s!"{result}"

partial def repl (defs : List (String × Term)
    := churchDefs) : IO Unit := do
  IO.print "λ> "
  let input ← (← IO.getStdin).getLine
  let input := input.trim
  if input == ":quit" then return
  if input == ":defs" then
    for (name, _) in defs do IO.print s!"{name} "
    IO.println ""; repl defs; return
  if input.startsWith ":let " then
    let rest := input.drop 5 |>.trim
    match rest.splitOn " = " with
    | [name, expr] =>
      match parse expr with
      | some term =>
        let term := expandDefs defs term
        IO.println s!"{name} defined"
        repl ((name, term) :: defs)
      | none => IO.println "Parse error in definition"; repl defs
    | _ => IO.println "Usage: :let name = expr"; repl defs
    return
  match parse input with
  | none => IO.println "Parse error"; repl defs
  | some namedTerm =>
    let expanded := expandDefs defs namedTerm
    IO.println s!"Parsed:    {expanded}"
    match toDeBruijn [] expanded with
    | none => IO.println "Error: unbound variable"; repl defs
    | some dbTerm =>
      IO.println s!"De Bruijn: {dbTerm}"
      IO.println s!"Result:    {eval 10000 dbTerm}"
      IO.println ""; repl defs

def main : IO Unit := repl
\end{minted}

\end{document}
